% La violence — Philosophie 1re Tech — Cours signé OnBuch+
% Programme officiel MINESEC. Compiler avec : tectonic N09-violence.tex
\newcommand{\DOCMATIERE}{Philosophie}
\newcommand{\DOCNIVEAU}{1re Tech}
\newcommand{\DOCMODULE}{Philosophie – classe de Première technique}
\newcommand{\DOCLECON}{N09}
\newcommand{\DOCTITRE}{La violence}
% OnBuch+ — préambule « cours signés » ENRICHI (compilé avec tectonic / XeLaTeX)
% Système visuel « Pop » d'OnBuch+ : fond crème (jamais blanc), bordures encre
% épaisses, ombres « dures » décalées, titres Archivo Black en capitales.
% Version MATHÉMATIQUES : graphes (pgfplots/tikz), boîtes pédagogiques
% (définition, propriété, méthode, attention, le savais-tu, exercice résolu,
% exercice, TP, bilan), page de garde et signature OnBuch+ ancrée en bas.
%
% Macros attendues dans main.tex AVANT \input{preamble} :
%   \DOCMATIERE  \DOCNIVEAU  \DOCMODULE  \DOCLECON (numéro)  \DOCTITRE
\documentclass[11pt]{article}

\usepackage{fontspec}
\usepackage[french]{babel}
\usepackage[a4paper,margin=2cm,top=3.4cm,bottom=2.8cm,headheight=1.4cm,footskip=1.3cm]{geometry}
\usepackage{amsmath,amssymb,amsfonts}
\usepackage{mathtools} % \xmapsto, \coloneqq... souvent utilisés par les modèles
\usepackage{cancel} % simplifications barrées (bilans de forces)
% \abs{z} (module, valeur absolue) et \norm{u} (norme), utilisés par les modèles
\providecommand{\abs}{}\let\abs\relax\DeclarePairedDelimiter{\abs}{\lvert}{\rvert}
\providecommand{\norm}{}\let\norm\relax\DeclarePairedDelimiter{\norm}{\lVert}{\rVert}
\usepackage[table]{xcolor} % [table] : fournit aussi \rowcolors (alternance de lignes)
\usepackage{pagecolor}
\usepackage{tikz}
\usetikzlibrary{arrows.meta,positioning,calc,shapes.geometric,shapes.misc,shapes.arrows,decorations.pathmorphing,decorations.pathreplacing,decorations.markings,patterns,shadows,mindmap,trees,fit,backgrounds,matrix,intersections,angles,quotes}
\usepackage{pgfplots}
\usepackage[version=4]{mhchem} % équations nucléaires \ce{^{235}_{92}U}
\usepackage[european,siunitx]{circuitikz} % schémas électriques
\pgfplotsset{compat=1.18}
\usepgfplotslibrary{fillbetween,groupplots} % aires entre courbes (calcul intégral)
\usepackage{enumitem}
\usepackage{fancyhdr}
% [most] charge toutes les bibliothèques tcolorbox (listings, minted, raster,
% documentation...) et gonfle inutilement la mémoire TeX sur un document long
% (~300 boîtes) : on ne charge que ce qui est réellement utilisé.
\usepackage[skins,breakable]{tcolorbox}
\usepackage{graphicx}
\usepackage{colortbl}
\usepackage{booktabs}
\usepackage{tabularx}
\usepackage{multirow}
\usepackage{diagbox} % case d'en-tête diagonale des tableaux à double entrée
% Colonnes extensibles centrée / gauche / droite (C, L, R), souvent utilisées par les modèles
\newcolumntype{C}{>{\centering\arraybackslash}X}
\newcolumntype{L}{>{\raggedright\arraybackslash}X}
\newcolumntype{R}{>{\raggedleft\arraybackslash}X}
\usepackage{array}
\usepackage{multicol}
\usepackage{siunitx}
% parse-numbers=false : accepte aussi \SI{2,5 \times 10^{-3}}{...}
\sisetup{locale=FR,output-decimal-marker={,},group-minimum-digits=4,parse-numbers=false}
\DeclareSIUnit\ohm{\ensuremath{\symup{\Omega}}} % U+2126 absent de la police
\DeclareSIUnit\division{div} % réglages d'oscilloscope (ms/div, V/div)
\DeclareSIUnit\tr{tr} % tours (vitesse de rotation, ex. \SI{1500}{\tr\per\minute})
\DeclareSIUnit\molar{M} % concentration molaire (ex. \SI{3}{\molar}, \SI{1}{\micro\molar})
\DeclareSIUnit\an{an} % année (le modèle confond parfois avec \year, un compteur TeX) — ex. \SI{5}{\kilo\meter\per\an}
\DeclareSIUnit\FCFA{FCFA} % franc CFA (ex. \SI{500}{\FCFA\per\kilo\gram})
\DeclareSIUnit\degreeBrix{\textdegree Brix} % teneur en sucre d'un sirop/jus (réfractométrie)
\DeclareSIUnit\month{mois} % le modèle utilise parfois \month, un compteur TeX, comme unité de durée
\DeclareSIUnit\microliter{\micro\liter} % le modèle écrit parfois \microliter en un seul token
\DeclareSIUnit\million{millions} % dénombrement (ex. \SI{15}{\million\per\mL} spermatozoïdes)
\usepackage{qrcode}
% \degree : commande fréquemment utilisée par les modèles pour un angle ou une
% température en dehors de \SI{}{} (non fournie par les paquets chargés).
\providecommand{\degree}{^{\circ}}
% \qed (fin de démonstration), utilisé par les modèles sans amsthm
\providecommand{\qed}{\hfill$\square$}
% \barre : double barre des valeurs interdites dans un tableau de variations (tkz-tab)
\providecommand{\barre}{\ensuremath{\|}}
% environnement proof (amsthm non chargé) utilisé par les modèles
\ifdefined\proof\else\newenvironment{proof}[1][Démonstration]{\par\noindent\textit{#1.}\ }{\qed\par}\fi
\usepackage{lastpage}
\usepackage{hyperref}
\hypersetup{hidelinks}

% ============================================================
% Palette « Pop » OnBuch+ — mêmes hex que la classe P (plus_ui.dart)
% ============================================================
\definecolor{poporange}{HTML}{F59321}
\definecolor{popdark}{HTML}{E07A0C}
\definecolor{popcream}{HTML}{FFF6E8}
\definecolor{popink}{HTML}{1C1714}
\definecolor{poppurple}{HTML}{7A5AE0}
\definecolor{popgreen}{HTML}{2E9E6A}
\definecolor{poppink}{HTML}{E0457B}
\definecolor{popblue}{HTML}{2F7DE1}
\definecolor{popgold}{HTML}{C98A12}
\definecolor{popmuted}{HTML}{8A7F76}
\definecolor{popline}{HTML}{EADFD2}
\definecolor{popgray}{HTML}{8A7F76}
\colorlet{popgrey}{popgray}
% Alias tolérants (le modèle nomme parfois les couleurs différemment)
\colorlet{popwhite}{white}
\colorlet{popblack}{popink}
\colorlet{popred}{poppink}
\colorlet{popyellow}{popgold}
% Teintes claires (fonds de tableaux/encadrés) — le modèle nomme parfois
% les couleurs avec un suffixe L (ex. popblueL) sans les avoir définies.
\colorlet{poporangeL}{poporange!20}
\colorlet{popdarkL}{popdark!20}
\colorlet{poppurpleL}{poppurple!20}
\colorlet{popgreenL}{popgreen!20}
\colorlet{poppinkL}{poppink!20}
\colorlet{popblueL}{popblue!20}
\colorlet{popgoldL}{popgold!20}
\colorlet{popredL}{popred!20}
\colorlet{popgrayL}{popgray!20}
% Teintes claires pour fonds de tableaux / zones
\colorlet{popblueL}{popblue!10!white}
\colorlet{poporangeL}{poporange!14!white}
\colorlet{popgreenL}{popgreen!12!white}
\colorlet{poppurpleL}{poppurple!12!white}
\colorlet{poppinkL}{poppink!12!white}
\colorlet{popgoldL}{popgold!14!white}

\pagecolor{popcream}

% ============================================================
% Typographie
% ============================================================
\defaultfontfeatures{Ligatures=TeX}
\setmainfont{PlusJakartaSans-Medium.ttf}[Path=../../../fonts/,BoldFont=PlusJakartaSans-Bold.ttf]
% Maths en sans-serif (Fira Math) avec chiffres et lettres droites de Plus
% Jakarta, pour que les formules se fondent dans le texte.
\usepackage[math-style=ISO,bold-style=ISO]{unicode-math}
\setmathfont{FiraMath-Regular.otf}[Path=../../../fonts/]
\setmathfont{PlusJakartaSans-Medium.ttf}[Path=../../../fonts/,range={up/num,up/{Latin,latin},\mathup}]
\usepackage{icomma} % virgule décimale sans espace en mode math
% Caractères mathématiques Unicode absents des polices de texte (Plus Jakarta,
% Archivo Black) : rendus par leur équivalent mathématique, en texte comme en
% formule — sinon ils s'affichent en carré vide (ex. « Arithmétique dans ℤ »).
\usepackage{newunicodechar}
\newunicodechar{ℝ}{\ensuremath{\mathbb{R}}}
\newunicodechar{ℤ}{\ensuremath{\mathbb{Z}}}
\newunicodechar{ℂ}{\ensuremath{\mathbb{C}}}
\newunicodechar{ℕ}{\ensuremath{\mathbb{N}}}
\newunicodechar{ℚ}{\ensuremath{\mathbb{Q}}}
\newunicodechar{𝔻}{\ensuremath{\mathbb{D}}}
\newunicodechar{α}{\ensuremath{\alpha}}
\newunicodechar{β}{\ensuremath{\beta}}
\newunicodechar{γ}{\ensuremath{\gamma}}
\newunicodechar{δ}{\ensuremath{\delta}}
\newunicodechar{ε}{\ensuremath{\varepsilon}}
\newunicodechar{θ}{\ensuremath{\theta}}
\newunicodechar{λ}{\ensuremath{\lambda}}
\newunicodechar{μ}{\ensuremath{\mu}}
\newunicodechar{σ}{\ensuremath{\sigma}}
\newunicodechar{τ}{\ensuremath{\tau}}
\newunicodechar{φ}{\ensuremath{\varphi}}
\newunicodechar{ψ}{\ensuremath{\psi}}
\newunicodechar{ω}{\ensuremath{\omega}}
\newunicodechar{Σ}{\ensuremath{\Sigma}}
% majuscules produites par \MakeUppercase dans les titres (α-aminés -> Α)
\newunicodechar{Α}{\ensuremath{\alpha}}
\newunicodechar{Β}{\ensuremath{\beta}}
\newunicodechar{Γ}{\ensuremath{\Gamma}}
\newunicodechar{Δ}{\ensuremath{\Delta}}
\newunicodechar{Ε}{\ensuremath{\varepsilon}}
\newunicodechar{Θ}{\ensuremath{\Theta}}
\newunicodechar{Λ}{\ensuremath{\Lambda}}
\newunicodechar{Μ}{\ensuremath{\mu}}
\newunicodechar{Τ}{\ensuremath{\tau}}
\newunicodechar{Φ}{\ensuremath{\Phi}}
\newunicodechar{Ψ}{\ensuremath{\Psi}}
\newunicodechar{Ω}{\ensuremath{\Omega}}
\newunicodechar{↦}{\ensuremath{\mapsto}}
\newunicodechar{∈}{\ensuremath{\in}}
\newunicodechar{∉}{\ensuremath{\notin}}
\newunicodechar{∧}{\ensuremath{\wedge}}
\newunicodechar{⇔}{\ensuremath{\Leftrightarrow}}
\newunicodechar{⟺}{\ensuremath{\Longleftrightarrow}}
\newunicodechar{⇒}{\ensuremath{\Rightarrow}}
\newunicodechar{⟹}{\ensuremath{\Longrightarrow}}
\newunicodechar{∘}{\ensuremath{\circ}}
\newunicodechar{∪}{\ensuremath{\cup}}
\newunicodechar{∩}{\ensuremath{\cap}}
\newunicodechar{≡}{\ensuremath{\equiv}}
\newunicodechar{⊂}{\ensuremath{\subset}}
\newunicodechar{⊥}{\ensuremath{\perp}}
\newunicodechar{∃}{\ensuremath{\exists}}
\newunicodechar{∀}{\ensuremath{\forall}}
\newunicodechar{∛}{\ensuremath{\sqrt[3]{\,}}}
\newunicodechar{∜}{\ensuremath{\sqrt[4]{\,}}}
\newunicodechar{⃗}{\ensuremath{{}^{\to}}}
\newunicodechar{ⁿ}{\ifmmode^{n}\else\textsuperscript{\ensuremath{n}}\fi}
\newunicodechar{ˣ}{\ifmmode^{x}\else\textsuperscript{\ensuremath{x}}\fi}
\newunicodechar{ᵇ}{\ifmmode^{b}\else\textsuperscript{\ensuremath{b}}\fi}
\newunicodechar{ʳ}{\ifmmode^{r}\else\textsuperscript{\ensuremath{r}}\fi}
\newunicodechar{ᵘ}{\ifmmode^{u}\else\textsuperscript{\ensuremath{u}}\fi}
\newunicodechar{ʸ}{\ifmmode^{y}\else\textsuperscript{\ensuremath{y}}\fi}
\newunicodechar{ᵃ}{\ifmmode^{a}\else\textsuperscript{\ensuremath{a}}\fi}
\newunicodechar{ᵉ}{\ifmmode^{e}\else\textsuperscript{\ensuremath{e}}\fi}
\newunicodechar{ᵏ}{\ifmmode^{k}\else\textsuperscript{\ensuremath{k}}\fi}
\newunicodechar{ᵖ}{\ifmmode^{p}\else\textsuperscript{\ensuremath{p}}\fi}
\newunicodechar{ᴺ}{\ifmmode^{N}\else\textsuperscript{\ensuremath{N}}\fi}
\newunicodechar{ᶜ}{\ifmmode^{c}\else\textsuperscript{\ensuremath{c}}\fi}
\newunicodechar{ʰ}{\ifmmode^{h}\else\textsuperscript{\ensuremath{h}}\fi}
\newunicodechar{ᵠ}{\ifmmode^{q}\else\textsuperscript{\ensuremath{q}}\fi}
\newunicodechar{ₙ}{\ifmmode_{n}\else\textsubscript{\ensuremath{n}}\fi}
\newunicodechar{ₐ}{\ifmmode_{a}\else\textsubscript{\ensuremath{a}}\fi}
\newunicodechar{ᵢ}{\ifmmode_{i}\else\textsubscript{\ensuremath{i}}\fi}
\newunicodechar{ᵣ}{\ifmmode_{r}\else\textsubscript{\ensuremath{r}}\fi}
\newunicodechar{ₖ}{\ifmmode_{k}\else\textsubscript{\ensuremath{k}}\fi}
\newunicodechar{ᵦ}{\ifmmode_{\beta}\else\textsubscript{\ensuremath{\beta}}\fi}
\newunicodechar{ₓ}{\ifmmode_{x}\else\textsubscript{\ensuremath{x}}\fi}
\newfontfamily\popdisplay{ArchivoBlack-Regular.ttf}[Path=../../../fonts/]
\newfontfamily\poplogo{SpaceGrotesk-Bold.ttf}[Path=../../../fonts/]
\newfontfamily\popsemi{PlusJakartaSans-SemiBold.ttf}[Path=../../../fonts/]
\newfontfamily\popxbold{PlusJakartaSans-ExtraBold.ttf}[Path=../../../fonts/]
\newfontfamily\popxboldls{PlusJakartaSans-ExtraBold.ttf}[Path=../../../fonts/,LetterSpace=3.0]

\setlist[enumerate]{leftmargin=1.7em,itemsep=5pt,topsep=4pt}
\setlist[itemize]{leftmargin=1.7em,itemsep=4pt,topsep=4pt,label={\color{poporange}\textbullet}}
\setlength{\parindent}{0pt}
\setlength{\parskip}{5pt}
\renewcommand{\arraystretch}{1.25}

% pgfplots : style Pop par défaut
\pgfplotsset{
  every axis/.append style={axis line style={popink,line width=1pt},tick style={popink},
    grid style={popline,line width=0.5pt},label style={font=\small},tick label style={font=\footnotesize}},
  popcurve/.style={poporange,line width=1.8pt,no markers,smooth},
}
% Même style disponible dans un \draw TikZ simple (hors pgfplots)
\tikzset{popcurve/.style={poporange,line width=1.8pt}}
\tikzset{popgrid/.style={popline,very thin}}
\colorlet{popcurve}{poporange} % le nom du style est parfois utilisé comme couleur

% ============================================================
% Pill / squiggle
% ============================================================
\newcommand{\poppill}[3]{%
  \tikz[baseline=-0.55ex]\node[draw=popink,line width=1.2pt,rounded corners=2.6pt,%
    fill=#2,inner xsep=6.5pt,inner ysep=3pt]{%
    \popxboldls\fontsize{7}{7}\selectfont\color{#3}\MakeUppercase{#1}};%
}
\newcommand{\popsquiggle}[1]{%
  \tikz[baseline=-0.4ex]{
    \draw[#1,line width=2.6pt,line cap=round]
      plot [smooth] coordinates {(0,0.09) (0.34,0.01) (0.68,0.21) (1.02,0.01) (1.36,0.10)};
  }%
}
% Mot-clé surligné (marqueur orange sous le texte)
\newcommand{\cle}[1]{{\popxbold\color{popdark}#1}}

% ============================================================
% En-tête / pied
% ============================================================
\pagestyle{fancy}
\fancyhf{}
\renewcommand{\headrulewidth}{0pt}
\renewcommand{\footrulewidth}{0pt}
\lhead{\raisebox{-0.15ex}{\poplogo\fontsize{13}{13}\selectfont\color{popink}OnBuch\color{poporange}+}}
\rhead{\poppill{\DOCMATIERE\ \textbullet\ \DOCNIVEAU\ \textbullet\ Leçon \DOCLECON}{white}{popink}}
\lfoot{\popxbold\fontsize{7.6}{7.6}\selectfont\color{popmuted}\MakeUppercase{Cours signé OnBuch+}\ \textbullet\ {\popsemi\DOCTITRE}}
\rfoot{\tikz[baseline=-0.6ex]\node[rounded corners=8.5pt,fill=popink,minimum height=17pt,minimum width=17pt,inner xsep=5pt,inner sep=0]{\popxbold\fontsize{8}{8}\selectfont\color{popcream}\thepage\,/\,\pageref*{LastPage}};}

% ============================================================
% Titres
% ============================================================
\newcommand{\chaptitre}[2]{%
  \begin{tcolorbox}[enhanced,colback=poporange,colframe=popink,boxrule=2.6pt,arc=11pt,
      drop shadow=popink,left=15pt,right=15pt,top=12pt,bottom=11pt,before skip=3pt,after skip=18pt]
    \poppill{#2}{popink}{popcream}\par\vspace{9pt}
    {\raggedright\hyphenpenalty=10000\exhyphenpenalty=10000\popdisplay\fontsize{22}{24}\selectfont\color{popcream}\MakeUppercase{#1}\par}
  \end{tcolorbox}
}
% Section numérotée : pastille orange avec le numéro + titre Archivo
\newcounter{popsec}
\newcommand{\coursec}[1]{%
  \stepcounter{popsec}\setcounter{popsub}{0}%
  \par\vspace{16pt}\noindent
  \begin{minipage}{\linewidth}
  \tikz[baseline=-0.7ex]\node[draw=popink,line width=1.6pt,fill=poporange,rounded corners=4pt,minimum size=22pt,inner sep=0]{\popdisplay\fontsize{12}{12}\selectfont\color{popcream}\thepopsec};%
  \hspace{8pt}{\popdisplay\fontsize{15}{17}\selectfont\color{popink}\MakeUppercase{#1}}\par
  \vspace{4pt}\noindent{\color{poporange}\rule{36pt}{3.6pt}}
  \end{minipage}\par\nopagebreak\vspace{8pt}
}
\newcounter{popsub}
\newcommand{\courssub}[1]{%
  \stepcounter{popsub}%
  \par\vspace{9pt}\noindent{\popxbold\fontsize{11.5}{13}\selectfont\color{popink}{\color{poporange}\thepopsec.\thepopsub}\hspace{6pt}#1}\par\nopagebreak\vspace{4pt}
}

% ============================================================
% Boîtes pédagogiques — même recette : fond blanc, bordure épaisse colorée,
% ombre dure encre, badge plein. Titre optionnel [#1].
% ============================================================
\tcbset{popbase/.style={breakable,enhanced,colback=white,boxrule=2.3pt,arc=9pt,
  shadow={3pt}{-3pt}{0pt}{popink},left=14pt,right=14pt,top=11pt,bottom=11pt,before skip=11pt,after skip=15pt,
  fontupper=\popsemi\fontsize{10.6}{14.5}\selectfont\color{popink},
  fontlower=\popsemi\fontsize{10.6}{14.5}\selectfont\color{popink}}}
% Coche dessinée (le glyphe ✓ n'existe pas dans les polices du document)
\newcommand{\popcheck}[1]{\tikz[baseline=-0.5ex]\draw[#1,line width=1.6pt,line cap=round,line join=round](-0.13,0.0)--(-0.04,-0.09)--(0.13,0.1);}
\providecommand{\coche}{\popcheck{popgreen}}
\providecommand{\resultat}[1]{\par\smallskip\noindent\fcolorbox{popgreen}{popgreenL}{\parbox{\dimexpr\linewidth-2\fboxsep-2\fboxrule\relax}{\textbf{Résultat.} #1}}\par\smallskip}
\newcommand{\popboxhead}[3]{\poppill{#1}{#2}{white}\ifx\relax#3\relax\else\hspace{6pt}{\popxbold\fontsize{10.6}{12}\selectfont\color{popink}#3}\fi\par\vspace{7pt}}

\newtcolorbox{aretenir}[1][]{popbase,colframe=popgreen,before upper={\popboxhead{À retenir}{popgreen}{#1}}}
\newtcolorbox{exemplebox}[1][]{popbase,colframe=poporange,before upper={\popboxhead{Exemple}{poporange}{#1}}}
\newtcolorbox{definition}[1][]{popbase,colframe=popblue,before upper={\popboxhead{Définition}{popblue}{#1}}}
\newtcolorbox{propriete}[1][]{popbase,colframe=poppurple,before upper={\popboxhead{Propriété}{poppurple}{#1}}}
\newtcolorbox{methode}[1][]{popbase,colframe=popgold,before upper={\popboxhead{Méthode}{popgold}{#1}}}
\newtcolorbox{attention}[1][]{popbase,colframe=poppink,before upper={\popboxhead{Attention}{poppink}{#1}}}
\newtcolorbox{experience}[1][]{popbase,colframe=popblue,colback=popblueL,before upper={\popboxhead{Expérience}{popblue}{#1}}}
% « Le savais-tu ? » : carte violette pleine (contexte, culture, Cameroun)
\newtcolorbox{savaistu}[1][]{popbase,colback=poppurple,colframe=popink,
  fontupper=\popsemi\fontsize{10.4}{14.2}\selectfont\color{white},
  before upper={\poppill{Le savais-tu\,?}{white}{poppurple}\ifx\relax#1\relax\else\hspace{6pt}{\popxbold\color{white}#1}\fi\par\vspace{7pt}}}
% Exercice résolu : énoncé en haut, \tcblower puis la solution sur fond crème
\newcounter{popexo}\newcounter{popexr}
\newtcolorbox{exoresolu}[1][]{popbase,colframe=poporange,colbacklower=popcream,
  segmentation style={popink,line width=1.2pt,dashed},
  before upper={\stepcounter{popexr}\popboxhead{Exercice résolu \thepopexr}{poporange}{#1}},
  before lower={\poppill{Solution}{popink}{popcream}\par\vspace{6pt}}}
% Exercice (non résolu en place) — la correction est dans la partie Corrigés
\newtcolorbox{exercice}[1][]{popbase,colframe=popink,
  before upper={\stepcounter{popexo}\popboxhead{Exercice \thepopexo}{popink}{#1}}}
% Corrigé
\newtcolorbox{corrige}[1][]{popbase,colframe=popgreen,colback=popgreenL,
  before upper={\popboxhead{Corrigé}{popgreen}{#1}}}
% Objectifs / prérequis / bilan
\newtcolorbox{objectifs}{popbase,colframe=popink,colback=poporangeL,
  before upper={\popboxhead{Objectifs de la leçon}{popink}{}}}
\newtcolorbox{prerequis}{popbase,colframe=popink,colback=popblueL,
  before upper={\popboxhead{Prérequis}{popblue}{}}}
\newtcolorbox{bilan}{popbase,colframe=popink,colback=white,boxrule=2.8pt,
  before upper={\popboxhead{Fiche bilan}{poporange}{L'essentiel à connaître pour l'examen}}}
% Figure encadrée avec légende
\newtcolorbox{popfigure}[1][]{enhanced,breakable=false,colback=white,colframe=popline,boxrule=1.4pt,arc=8pt,
  left=10pt,right=10pt,top=10pt,bottom=8pt,before skip=10pt,after skip=12pt,halign=center,
  lower separated=false,halign lower=center,
  fontlower=\popsemi\fontsize{9}{11.5}\selectfont\color{popmuted},
  #1}
\newcommand{\legende}[1]{\par\vspace{6pt}{\popsemi\fontsize{9}{11.5}\selectfont\color{popmuted}#1}}

% ============================================================
% Signature OnBuch+
% ============================================================
\newcommand{\poplogoinline}[1]{{\poplogo\fontsize{#1}{#1}\selectfont\color{popink}OnBuch\color{poporange}+}}
\newcommand{\signatureonbuch}{%
  \begin{tcolorbox}[enhanced,colback=popink,colframe=popink,boxrule=2.2pt,arc=10pt,
      drop shadow=poporange,left=14pt,right=14pt,top=10pt,bottom=10pt,before skip=0pt,after skip=0pt]
    \tikz[baseline=-0.7ex]\node[circle,fill=popgreen,draw=popcream,line width=1.3pt,minimum size=22pt,inner sep=0]{\popcheck{white}};%
    \hspace{9pt}\begin{minipage}[c]{0.62\linewidth}
      {\popxbold\fontsize{12}{13}\selectfont\color{popcream}Signé\ \textbullet\ {\poplogo OnBuch\color{poporange}+}}\par\vspace{2pt}
      {\raggedright\popsemi\fontsize{8}{10.5}\selectfont\color{popline}\MakeUppercase{Équipe pédagogique OnBuch+}\\\MakeUppercase{Conforme au programme officiel MINESEC}\par}
    \end{minipage}\hfill
    {\poplogo\fontsize{22}{22}\selectfont\color{popcream}OnBuch\color{poporange}+}
  \end{tcolorbox}%
}

% ============================================================
% Page de garde
% #1 titre · #2 matière · #3 niveau · #4 module · #5 n° leçon · #6 durée ·
% #7 description / objectifs (texte ou itemize)
% ============================================================
\newcommand{\pagegarde}[7]{%
  \thispagestyle{empty}
  \newgeometry{a4paper,margin=1.7cm,top=1.6cm,bottom=1.4cm}%
  \begin{tikzpicture}[remember picture,overlay]
    \fill[poporange!18] ([xshift=-3.2cm,yshift=-3.6cm]current page.north east) circle (4.6cm);
    \fill[poppurple!14] ([xshift=2.4cm,yshift=7.6cm]current page.south west) circle (3.8cm);
  \end{tikzpicture}%
  {\poplogo\fontsize{30}{30}\selectfont\color{popink}OnBuch\color{poporange}+}\hfill
  \raisebox{0.6ex}{\poppill{Cours signé \textbullet\ Premium}{popink}{popcream}}\par
  \vspace{1pt}\popsquiggle{poppurple}\par
  \vspace{8pt}
  \begin{tcolorbox}[enhanced,colback=poporange,colframe=popink,boxrule=2.8pt,arc=13pt,
      drop shadow=popink,left=18pt,right=18pt,top=12pt,bottom=13pt,after skip=10pt]
    \poppill{#2 \textbullet\ #3}{popink}{popcream}\hspace{5pt}\poppill{#4}{white}{popink}\par\vspace{8pt}
    {\popdisplay\fontsize{14}{14}\selectfont\color{popink}LEÇON #5}\par\vspace{4pt}
    {\raggedright\hyphenpenalty=10000\exhyphenpenalty=10000\popdisplay\fontsize{25}{27}\selectfont\color{popcream}\MakeUppercase{#1}\par}
    \vspace{8pt}
    {\popxbold\fontsize{9.5}{9.5}\selectfont\color{popink}\MakeUppercase{Durée conseillée : #6}}
  \end{tcolorbox}
  \begin{tcolorbox}[enhanced,colback=white,colframe=popink,boxrule=2.2pt,arc=10pt,
      drop shadow=popink,left=16pt,right=16pt,top=10pt,bottom=10pt,after skip=8pt]
    \poppill{Dans cette leçon}{poppurple}{white}\par\vspace{5pt}
    \popsemi\fontsize{9.3}{12.6}\selectfont\color{popink}#7
  \end{tcolorbox}
  \noindent\begin{minipage}[t]{0.487\linewidth}
    \begin{tcolorbox}[enhanced,colback=popgreen,colframe=popink,boxrule=2.2pt,arc=10pt,
        drop shadow=popink,left=11pt,right=11pt,top=10pt,bottom=10pt,height=3.1cm,valign=center]
      \tcbox[colback=white,colframe=popink,boxrule=1.3pt,arc=5pt,boxsep=0pt,nobeforeafter,
        left=3pt,right=3pt,top=3pt,bottom=3pt,valign=center]{\qrcode[height=1.9cm]{https://play.google.com/store/apps/details?id=com.luvvix.onbuch&hl=fr}}%
      \hspace{8pt}\begin{minipage}[c]{0.5\linewidth}
        {\popdisplay\fontsize{11}{12.5}\selectfont\color{white}\MakeUppercase{Télécharge OnBuch}}\par\vspace{4pt}
        {\raggedright\popsemi\fontsize{8.4}{11}\selectfont\color{white}Gratuit \textbullet\ Google Play\\Tuteur IA, annales, concours, résultats\par}
      \end{minipage}
    \end{tcolorbox}
  \end{minipage}\hfill
  \begin{minipage}[t]{0.487\linewidth}
    \begin{tcolorbox}[enhanced,colback=poppurple,colframe=popink,boxrule=2.2pt,arc=10pt,
        drop shadow=popink,left=11pt,right=11pt,top=10pt,bottom=10pt,height=3.1cm,valign=center]
      \tikz[baseline=-0.7ex]\node[circle,fill=white,draw=popink,line width=1.3pt,minimum size=1.6cm,inner sep=0]{\popdisplay\fontsize{24}{24}\selectfont\color{poppurple}+};%
      \hspace{8pt}\begin{minipage}[c]{0.6\linewidth}
        {\popdisplay\fontsize{11}{12.5}\selectfont\color{white}\MakeUppercase{Passe à OnBuch+}}\par\vspace{4pt}
        {\raggedright\popsemi\fontsize{8.4}{11}\selectfont\color{white}Tous les cours signés, Tuteur IA illimité, fiches PDF — onglet OnBuch+ de l'app\par}
      \end{minipage}
    \end{tcolorbox}
  \end{minipage}
  \vspace{8pt}
  \popcontenu
  \vfill
  \signatureonbuch
  \restoregeometry
  \newpage
}

% Rangée « Ce cours contient » (page de garde)
\newcommand{\poptile}[3]{% #1 couleur #2 gros texte #3 légende
  \begin{tcolorbox}[enhanced,colback=#1,colframe=popink,boxrule=2pt,arc=9pt,shadow={2.5pt}{-2.5pt}{0pt}{popink},
      left=6pt,right=6pt,top=9pt,bottom=9pt,halign=center,valign=center,height=1.75cm,width=\linewidth,nobeforeafter]
    {\popdisplay\fontsize{12.5}{13}\selectfont\color{white}\MakeUppercase{#2}}\par\vspace{3pt}
    {\centering\popsemi\fontsize{7.8}{9.5}\selectfont\color{white}#3\par}
  \end{tcolorbox}}
\newcommand{\popcontenu}{%
  \par\noindent{\popxboldls\fontsize{7.5}{7.5}\selectfont\color{popmuted}CE COURS SIGNÉ CONTIENT}\par\vspace{7pt}
  \noindent\begin{minipage}[t]{0.235\linewidth}\poptile{popblue}{Cours}{complet, illustré, pas à pas}\end{minipage}\hfill
  \begin{minipage}[t]{0.235\linewidth}\poptile{poporange}{Méthodes}{et exercices résolus}\end{minipage}\hfill
  \begin{minipage}[t]{0.235\linewidth}\poptile{poppink}{Exercices}{type examen + corrigés}\end{minipage}\hfill
  \begin{minipage}[t]{0.235\linewidth}\poptile{popgreen}{Fiche bilan}{l'essentiel à retenir}\end{minipage}\par}

% Sommaire
\newtcolorbox{sommaire}{enhanced,colback=white,colframe=popink,boxrule=2.2pt,arc=10pt,
  drop shadow=popink,left=18pt,right=18pt,top=13pt,bottom=13pt,before skip=6pt,after skip=20pt,
  before upper={\poppill{Sommaire}{poppurple}{white}\par\vspace{9pt}\popsemi\fontsize{10.6}{14.5}\selectfont\color{popink}}}

% Signature de fin de cours — poussée tout en bas de la dernière page
\newcommand{\findecours}{%
  \par\vfill\nopagebreak
  \begin{center}\popsquiggle{poporange}\end{center}\vspace{4pt}
  \signatureonbuch
}
\AtBeginDocument{\def\llbracket{\mathopen{[\mkern-3.5mu[}}\def\rrbracket{\mathclose{]\mkern-3.5mu]}}} % intervalles d'entiers (glyphe ⟦ absent de la police)
% Glyphes absents de Fira Math : redéfinis à partir de glyphes disponibles
\AtBeginDocument{%
  \def\setminus{\mathbin{\backslash}}\let\smallsetminus\setminus
  \def\vdots{\mathord{\vbox{\baselineskip3.2pt\lineskiplimit0pt\kern4pt\hbox{.}\hbox{.}\hbox{.}}}}%
  \def\ddots{\mathinner{\mkern1mu\raise6.5pt\hbox{.}\mkern2mu\raise3.6pt\hbox{.}\mkern2mu\raise0.7pt\hbox{.}\mkern1mu}}%
  \def\triangle{\mathord{\scalebox{0.9}{$\symup{\Delta}$}}}%
  \def\nabla{\mathord{\raisebox{\depth}{\scalebox{1}[-1]{$\symup{\Delta}$}}}}%
  \def\checkmark{\popcheck{popdark}}%
  \def\star{\mathbin{\ast}}\def\bigstar{\mathord{\ast}}%
  \def\diamond{\mathbin{\scalebox{0.6}{$\Diamond$}}}%
  \def\bigcup{\mathop{\mathchoice{\raisebox{-0.3em}{\scalebox{1.7}{$\cup$}}}{\scalebox{1.25}{$\cup$}}{\cup}{\cup}}\displaylimits}%
  \def\bigcap{\mathop{\mathchoice{\raisebox{-0.3em}{\scalebox{1.7}{$\cap$}}}{\scalebox{1.25}{$\cap$}}{\cap}{\cap}}\displaylimits}%
  \def\lesseqqgtr{\mathrel{\lessgtr}}\def\bigoplus{\mathop{\oplus}\displaylimits}%
}
\providecommand{\ssi}{si et seulement si} % abréviation fréquente
\AtBeginDocument{\providecommand{\wideparen}{\overparen}} % arc de cercle

%% FIGFIX-BEGIN v5 — correctifs globaux des figures (docs/onbuch-generation/tools/figfix)
\usepackage{adjustbox}\usepackage{etoolbox}\tcbuselibrary{raster}
% toute figure TikZ/pgfplots plus large que la ligne est réduite à la largeur disponible
\BeforeBeginEnvironment{tikzpicture}{\begin{adjustbox}{max width=\linewidth}}
\AfterEndEnvironment{tikzpicture}{\end{adjustbox}}
% légendes pgfplots lisibles par-dessus les courbes
\pgfplotsset{every axis/.append style={legend style={fill=white,fill opacity=0.92,text opacity=1,draw=black!15,inner sep=3pt}}}
% étiquettes d'arêtes (syntaxe edge["..."]) avec fond blanc
\tikzset{every edge quotes/.append style={fill=white,inner sep=1.2pt}}
%% FIGFIX-END
\begin{document}

\pagegarde{La violence}{Philosophie}{1re Tech}{Philosophie – classe de Première technique}{N09}{3 séances (fiche de progression harmonisée)}{Cette leçon interroge la violence dans sa définition, son origine et ses multiples formes (physique, psychologique, économique, symbolique, étatique), avant de conduire le procès philosophique de la violence : peut-elle jamais être légitime ? À partir de situations concrètes du quotidien camerounais, l'élève apprend à distinguer violence et force, légitimité et légalité, et à construire un jugement argumenté.
\vspace{4pt}
\begin{multicols}{2}\small
\begin{itemize}
\item À la fin de cette leçon, je sais définir la violence et la distinguer de la force, de la contrainte et de l'agressivité
\item À la fin de cette leçon, je sais expliquer les grandes thèses philosophiques sur l'origine de la violence (naturelle chez Hobbes, culturelle chez Rousseau)
\item À la fin de cette leçon, je sais identifier et classer les différents types de violence (physique, psychologique, verbale, économique, symbolique, politique)
\item À la fin de cette leçon, je sais analyser des situations concrètes du quotidien camerounais en y repérant les formes de violence en jeu
\item À la fin de cette leçon, je sais conduire le procès de la violence en confrontant les arguments de sa condamnation et ceux de sa justification (légitime défense, violence révolutionnaire, monopole étatique)
\item À la fin de cette leçon, je sais rédiger une conclusion argumentée et justifiée sur un sujet de dissertation portant sur la violence
\end{itemize}\end{multicols}}

\begin{sommaire}
\begin{multicols}{2}
\begin{enumerate}
\item Introduction : problématiser la violence
\item Définition de la violence
\item L'origine de la violence : naturelle ou culturelle ?
\item Les types de violence
\item Le procès de la violence (I) : la condamnation
\item Le procès de la violence (II) : peut-on la justifier ? Bilan
\item Atelier méthodologique : appliquer la leçon à des situations concrètes
\item Activité d'intégration
\item Méthodes et exercices résolus
\item Exercices
\item Corrigés des exercices
\item Fiche bilan
\end{enumerate}
\end{multicols}
\end{sommaire}

\begin{objectifs}
\begin{itemize}
\item À la fin de cette leçon, je sais définir la violence et la distinguer de la force, de la contrainte et de l'agressivité
\item À la fin de cette leçon, je sais expliquer les grandes thèses philosophiques sur l'origine de la violence (naturelle chez Hobbes, culturelle chez Rousseau)
\item À la fin de cette leçon, je sais identifier et classer les différents types de violence (physique, psychologique, verbale, économique, symbolique, politique)
\item À la fin de cette leçon, je sais analyser des situations concrètes du quotidien camerounais en y repérant les formes de violence en jeu
\item À la fin de cette leçon, je sais conduire le procès de la violence en confrontant les arguments de sa condamnation et ceux de sa justification (légitime défense, violence révolutionnaire, monopole étatique)
\item À la fin de cette leçon, je sais rédiger une conclusion argumentée et justifiée sur un sujet de dissertation portant sur la violence
\item À la fin de cette leçon, je sais mobiliser les références philosophiques (Hobbes, Rousseau, Weber, Arendt, Gandhi) pour étayer une réflexion personnelle
\end{itemize}
\end{objectifs}

\begin{prerequis}
\begin{itemize}
\item La notion d'État et le vivre-ensemble (Education à la citoyenneté, collège)
\item La distinction entre nature et culture (leçon sur l'homme, début d'année de Première)
\item La notion de liberté et ses limites (leçons précédentes de philosophie)
\item La notion de devoir et de loi morale
\item Les méthodes de l'analyse de sujet et de la définition des termes (acquis de méthodologie philosophique)
\end{itemize}
\end{prerequis}

\newpage

\coursec{Introduction : problématiser la violence}

\courssub{La situation-problème : la violence au quotidien au Cameroun}

Récemment, une bagarre entre élèves dans un lycée de Yaoundé a dégénéré : un élève a été blessé à l'arme blanche. La même semaine, un chauffeur de taxi à Douala a été lynché par une foule qui le soupçonnait de vol, et sur les réseaux sociaux camerounais, des campagnes de dénonciation des violences faites aux femmes se multiplient. La violence semble partout : dans la rue, à l'école, dans les stades, sur les écrans. Elle frappe aux portes de nos classes, s'invite dans nos familles, structure nos rapports sociaux. Face à cette omniprésence, l'élève de Première technique doit dépasser l'émotion légitime pour accéder à l'analyse philosophique rigoureuse, en vue du Probatoire.

\begin{exemplebox}[Le lynchage populaire : un fait social à interroger]
Dans les quartiers de Douala (Bépanda, New Bell, Akwa) comme de Yaoundé (Mvog-Ada, Essos, Nkolndongo), la « justice populaire » ou lynchage est un phénomène récurrent. Un voleur présumé, parfois un simple suspect, est poursuivi, tabassé, parfois brûlé vif par une foule en colère, avant l'arrivée de la police.
\par\medskip
\emph{Analyse philosophique :} Ce n'est pas seulement un fait divers sordide. C'est une \cle{transgression} majeure du monopole de la violence légitime par l'État (Weber). La foule s'arroge le droit de juger et d'exécuter, remplaçant la procédure judiciaire par l'exécution sommaire. Cela pose trois problèmes philosophiques immédiats :
\begin{enumerate}
    \item \textbf{La définition :} Est-ce de la violence ? Oui, physique, collective, extrême. Mais les auteurs se sentent-ils « violents » ou « justiciers » ?
    \item \textbf{L'origine :} Est-ce la « nature » humaine (instinct de meute, peur) ou un produit culturel (défiance envers la justice, impunité perçue, mimétisme) ?
    \item \textbf{La légitimité :} La foule invoque la « justice ». Peut-on justifier la violence au nom de la justice ? C'est le cœur du \cle{procès de la violence}.
\end{enumerate}
\end{exemplebox}

\courssub{Les définitions spontanées : entre « coup », « agression » et « injustice »}

Demandons aux élèves : « Qu'est-ce que la violence ? » Les réponses fusent, ancrées dans le vécu :
\begin{itemize}
    \item « C'est quand on tape quelqu'un. » (\emph{Violence physique, coups})
    \item « C'est quand on vous force à faire ce que vous ne voulez pas. » (\emph{Contrainte, atteinte à la liberté})
    \item « C'est l'injustice : le riche qui écrase le pauvre, le patron qui ne paie pas. » (\emph{Violence structurelle, symbolique})
    \item « C'est les insultes, le harcèlement à l'école. » (\emph{Violence psychologique, verbale})
\end{itemize}

Ces intuitions sont justes mais partielles. Elles révèlent la \cle{polysémie} du terme. Le mot « violence » vient du latin \emph{vis, vim} (force, puissance, énergie) et \emph{violentia} (action de violer, de forcer). Au sens large, c'est l'emploi de la force pour contraindre. Au sens strict, c'est l'abus de la force, la transgression d'une norme (juridique, morale, physique).

\begin{popfigure}
\begin{tikzpicture}[scale=0.9, every node/.style={font=\sffamily\small}]
% Central node
\node[draw=popdark, thick, circle, fill=popblue!20, text width=2.2cm, align=center, minimum height=1.8cm] (violence) at (0,0) {\textbf{VIOLENCE}\\\emph{(vis, violentia)}};

% Satellite nodes - first ring
\node[draw=poporange, thick, rounded corners, fill=poporangeL, text width=2.5cm, align=center] (coup) at (4.5, 2) {Coup\\Blessure\\Atteinte corporelle};
\node[draw=popgreen, thick, rounded corners, fill=popgreenL, text width=2.5cm, align=center] (contrainte) at (4.5, -2) {Contrainte\\Force\\Coercition};
\node[draw=poppurple, thick, rounded corners, fill=poppurpleL, text width=2.5cm, align=center] (injustice) at (-4.5, -2) {Injustice\\Abus de pouvoir\\Violation de droit};
\node[draw=poppink, thick, rounded corners, fill=poppinkL, text width=2.5cm, align=center] (peur) at (-4.5, 2) {Peur\\Terreur\\Intimidation};

% Second ring - specific types
\node[draw=popblue, thin, ellipse, fill=popblue!10, text width=2cm, align=center] (physique) at (7, 3) {Physique};
\node[draw=popblue, thin, ellipse, fill=popblue!10, text width=2cm, align=center] (psych) at (7, 1) {Psychologique};
\node[draw=poporange, thin, ellipse, fill=poporange!10, text width=2cm, align=center] (struct) at (7, -1) {Structurelle};
\node[draw=poporange, thin, ellipse, fill=poporange!10, text width=2cm, align=center] (symbol) at (7, -3) {Symbolique};
\node[draw=popgreen, thin, ellipse, fill=popgreen!10, text width=2cm, align=center] (legitime) at (-7, 3) {Légitime?\\ (État)};
\node[draw=popgreen, thin, ellipse, fill=popgreen!10, text width=2cm, align=center] (illegitime) at (-7, 1) {Illégitime\\ (Crime)};
\node[draw=poppurple, thin, ellipse, fill=poppurple!10, text width=2cm, align=center] (domest) at (-7, -1) {Domestique};
\node[draw=poppurple, thin, ellipse, fill=poppurple!10, text width=2cm, align=center] (collec) at (-7, -3) {Collective\\ (Émeute)};

% Connections
\foreach \src/\dst in {violence/coup, violence/contrainte, violence/injustice, violence/peur} {
    \draw[popline, thick, ->] (\src) -- (\dst);
}
\foreach \src/\dst in {coup/physique, coup/psych, contrainte/struct, contrainte/symbol, injustice/legitime, injustice/illegitime, peur/domest, peur/collec} {
    \draw[popmuted, -] (\src) -- (\dst);
}
\end{tikzpicture}
\legende{Schéma heuristique : la polysémie du concept de violence. Le terme recouvre des réalités hétérogènes (physique, psychologique, structurelle, symbolique) et des statuts normatifs opposés (légitime/illégitime). 1\textsuperscript{er} anneau (couleurs vives) : acceptions spontanées des élèves. 2\textsuperscript{e} anneau (couleurs pâles) : catégories philosophiques à construire dans la leçon.}
\end{popfigure}

\courssub{Les difficultés conceptuelles : toute force est-elle violence ? La violence est-elle toujours visible ?}

Deux écueils majeurs guettent l'analyse naïve.

\begin{attention}[Piège n°1 : Confondre violence, agressivité et colère]
\emph{Erreur fréquente :} « Il est violent » pour dire « Il est agressif » ou « Il est en colère ».
\par\medskip
\textbf{Distinction rigoureuse :}
\begin{itemize}
    \item \textbf{Colère :} Émotion, affect passager (colère de l'élève qui rate son devoir). Pas d'action nécessaire.
    \item \textbf{Agressivité :} Pulsion, disposition psychologique ou biologique à l'attaque (instinct de défense/prédation). Peut ne pas passer à l'acte.
    \item \textbf{Violence :} \textbf{Acte} (ou système d'actes) qui use de la force pour \emph{violer} l'intégrité (corps, psychisme, biens, droits) d'autrui ou de soi-même. Elle suppose une \cle{relation} (un agent, une victime, un contexte normatif).
\end{itemize}
\par\medskip
On peut être agressif sans être violent (se contenir), violent sans être agressif (chirurgien opérant, policier arrêtant un criminel — violence \emph{légitime}), en colère sans être ni l'un ni l'autre. La philosophie examine l'\emph{acte} et sa \emph{justification}, non seulement le mobile psychologique.
\end{attention}

\begin{attention}[Piège n°2 : Réduire la violence au visible (coups, sang)]
La violence \emph{structurelle} (Galtung) ou \emph{symbolique} (Bourdieu) est invisible mais dévastatrice.
\begin{itemize}
    \item \textbf{Exemple camerounais :} L'absence d'eau potable dans un quartier périphérique de Yaoundé alors que le centre-ville en dispose 24h/24. Personne ne frappe personne, mais une \cle{structure sociale} viole le droit à la vie, à la santé, à la dignité. C'est une violence lente, silencieuse, souvent plus meurtrière qu'une bagarre.
    \item \textbf{Violence symbolique :} L'imposition de catégories de pensée (ex. « les filles ne font pas la technique industrielle ») qui dominent sans avoir l'air de contraindre.
\end{itemize}
\emph{Règle méthodologique :} Ne jamais définir la violence par ses seuls instruments (armes, cris) mais par sa \textbf{structure} : rapport de force asymétrique + transgression d'une norme + atteinte à l'intégrité.
\end{attention}

\courssub{Formulation de la problématique}

De ces constats naissent les questions qui guideront toute l'unité. Une bonne problématique ne juxtapose pas des questions, elle les \emph{articule} autour d'une tension.

\begin{methode}[Problématiser un sujet de philosophie : 3 étapes]
\begin{enumerate}
    \item \textbf{Définir les termes} : Distinguer les sens (physique/moral, individuel/collectif, légitime/illégitime). Refuser le sens unique.
    \item \textbf{Dégager le paradoxe / la tension} : La violence est condamnée moralement (« Tu ne tueras point ») mais omniprésente socialement (État, éducation, révolution). Elle est destructrice mais parfois fondatrice (droit, ordre).
    \item \textbf{Formuler la question directrice} : Une question ouverte, qui appelle une réponse argumentée (thèses/antithèses/synthèse), pas un oui/non.
\end{enumerate}
\end{methode}

\paragraph{Application à notre leçon.}
\begin{enumerate}
    \item \textbf{Terme :} Violence = emploi de la force pour contraindre/violer (large) ; abus de force transgressif (strict).
    \item \textbf{Tension :} 
    \begin{itemize}
        \item \emph{Condamnation unanime :} « La violence est le refuge de l'incompétent » (Asimov) ; elle nie l'humanité de la victime.
        \item \emph{Présence nécessaire ? } L'État a le \emph{monopole de la violence physique légitime} (Weber). Une révolution contre la tyrannie use de violence. La légitime défense est un droit.
    \end{itemize}
    \item \textbf{Problématique unique :}
    \begin{center}
    \fbox{\parbox{0.9\linewidth}{\centering
    \textbf{\cle{Qu'est-ce que la violence ? D'où vient-elle (nature ou culture) ? Peut-elle être légitime, c'est-à-dire juste et fondée, ou est-elle toujours un mal à éradiquer ?}}
    }}
    \end{center}
\end{enumerate}

Cette problématique structure les trois grandes parties du programme :
\begin{itemize}
    \item \textbf{I. Définition et origine} (Qu'est-ce que c'est ? Naturelle ou culturelle ?)
    \item \textbf{II. Types de violence} (Physique, psychologique, structurelle, symbolique, politique, légitime/illégitime).
    \item \textbf{III. Procès de la violence} (Condamnation morale/juridique vs justifications politiques/éthiques : légitime défense, révolution, contrainte éducative, monopole étatique).
\end{itemize}

\begin{aretenir}[Ce qu'il faut retenir de l'introduction]
\begin{itemize}
    \item La violence n'est pas un fait brut ; c'est un \textbf{concept polysémique} (coups, contrainte, injustice, peur, structure).
    \item Elle se distingue de l'agressivité (pulsion) et de la colère (émotion) par son caractère \textbf{relationnel, normatif et intentionnel} (même si l'intention est implicite dans la structure).
    \item Elle n'est pas toujours visible : les violences \textbf{structurelles} et \textbf{symboliques} sont majeures dans nos sociétés camerounaises.
    \item La \textbf{problématique} unit trois questions : définition (ontologie), origine (anthropologie), légitimité (éthique/politique).
\end{itemize}
\end{aretenir}

\courssub{Annonce du plan de la leçon}

Nous allons maintenant dérouler le fil de l'analyse en trois temps forts, conformément au programme officiel de Première technique :

\begin{center}
\begin{tabularx}{\linewidth}{|l|X|}\hline
\textbf{Partie} & \textbf{Contenu et enjeux} \\ \hline
\textbf{I. Définition et origine de la violence} & Analyser la structure conceptuelle (force, contrainte, violation). Débat classique : \emph{inné} (instinct de mort, Thanatos chez Freud ; nature humaine belliqueuse chez Hobbes) vs \emph{acquis} (conditionnement social, mimétisme chez Girard ; aliénation chez Marx). \\ \hline
\textbf{II. Les types de violence} & Typologie opératoire : \emph{Physique / Psychologique} ; \emph{Individuelle / Collective} ; \emph{Directe / Structurelle / Symbolique} (Galtung, Bourdieu) ; \emph{Légitime / Illégitime} (Weber, droit positif camerounais). Application aux contextes locaux (violences scolaires, basées sur le genre, politiques). \\ \hline
\textbf{III. Le procès de la violence} & \textbf{(I) La condamnation} : Arguments moraux (Kant : traiter l'humanité comme fin), juridiques (État de droit), politiques (non-violence gandhienne). \textbf{(II) La justification possible ?} : Légitime défense, contrainte pédagogique, violence révolutionnaire (Fanon, Sartre), monopole étatique (Weber). \textbf{Bilan} : Critères de légitimité (proportionnalité, nécessité, finalité juste, autorité compétente). \\ \hline
\end{tabularx}
\end{center}

\begin{exoresolu}[Exercice d'application : qualifier des situations]
\textbf{Énoncé :} Qualifiez philosophiquement (type de violence, légitimité discutable) les trois situations ci-dessous en utilisant les distinctions de l'introduction.
\begin{enumerate}
    \item Un enseignant gifle un élève qui a insulté sa mère.
    \item Un quartier entier est privé d'électricité depuis trois mois alors que les factures sont payées ; les câbles volés ne sont pas remplacés par l'entreprise publique.
    \item Une manifestation étudiante pour la baisse des frais d'inscription est dispersée à coups de gaz lacrymogène et de matraques par les forces de l'ordre ; trois blessés graves.
\end{enumerate}
\tcblower
\textbf{Solution détaillée :}
\begin{enumerate}
    \item \textbf{Violence physique, individuelle, directe, illégitime (a priori).} L'enseignant a un pouvoir d'autorité mais pas le droit de porter atteinte à l'intégrité corporelle (Code de l'enfant, statut de la fonction publique). C'est un abus de pouvoir (violence symbolique + physique). \emph{Nuance :} Certains argueront « correction parentale » (débat culturel), mais le droit camerounais interdit les châtiments corporels à l'école.
    \item \textbf{Violence structurelle, collective, indirecte, illégitime.} Aucune agression physique directe. Mais la structure (entreprise publique, régulation défaillante) viole le droit au service public, à la santé (chaîne du froid médicaments), à la sécurité. Responsabilité diffuse, difficile à sanctionner pénalement, mais violence réelle (Galtung).
    \item \textbf{Violence physique, collective, directe, exercée par l'État (monopole légitime).} \emph{Problème :} Légitimité \emph{de principe} (maintien de l'ordre) mais \emph{illégitimité possible de fait} si disproportionnée (3 blessés graves pour une manifestation assise). Critères : nécessité absolue, proportionnalité, gradation de la force. Si l'ordre n'était pas menacé, c'est une violence illégitime (répression), engageant la responsabilité de l'État.
\end{enumerate}
\end{exoresolu}

\begin{savaistu}[Violence et société camerounaise : quelques repères chiffrés]
\begin{itemize}
    \item \textbf{Violences basées sur le genre (VBG) :} Selon le MINPROFF et l'INS (EDS-MICS 2018), \textbf{43\%} des femmes de 15-49 ans ont subi des violences physiques ou sexuelles depuis 15 ans. Le mari/partenaire est l'auteur principal (60\% des cas).
    \item \textbf{Violences scolaires :} Enquête UNICEF/MINESEC 2022 : \textbf{1 élève sur 3} déclare avoir subi des violences (coups, insultes, harcèlement) dans l'enceinte de l'établissement au cours de l'année.
    \item \textbf{Justice populaire :} Pas de statistique nationale officielle, mais la presse rapporte \textbf{une dizaine de cas médiatisés par an} dans les grandes villes (Douala, Yaoundé, Bafoussam), souvent liés au vol à l'arraché ou aux accusations de sorcellerie.
    \item \textbf{Cadre légal :} Constitution 1996 (Préambule : droits de l'homme), Code pénal (atteintes à l'intégrité), Loi 2016/007 sur le Code de l'enfant (interdiction châtiments corporels), Loi 2016/005 sur la décentralisation (police municipale).
\end{itemize}
\emph{Source : Rapports MINPROFF, INS, UNICEF, presse locale. Ces données ancrent la philosophie dans la réalité du terrain.}
\end{savaistu}

\coursec{Définition de la violence}

L'introduction nous a permis de problématiser la violence : nous avons vu que ce mot, omniprésent dans les journaux, les débats politiques et la vie quotidienne, est souvent employé de manière vague et contradictoire. On parle de violence physique, de violence verbale, de violence symbolique, parfois même de la « violence » d'une tempête ou d'une émotion. Mais si le mot recouvre tout, il ne signifie plus rien. Avant de pouvoir interroger l'origine de la violence, ses types et sa justification, il faut donc d'abord répondre à une question préalable, en apparence simple : \emph{qu'est-ce que la violence ?} Cette section a pour but de construire une définition rigoureuse du concept, en partant de son étymologie, puis en le distinguant de deux notions voisines avec lesquelles on le confond trop souvent : la \cle{force} et l'\cle{agressivité}.

\courssub{L'étymologie du mot « violence »}

\begin{definition}[Étymologie]
Le mot \cle{violence} vient du latin \emph{violentia}, qui signifie « force », « emportement », « caractère impétueux ». Ce terme est lui-même lié à la racine latine \emph{vis}, qui désigne la force, la puissance, l'énergie vitale. Le verbe \emph{violare} (dont dérive « violer ») signifie « traiter avec force », « outrager », « profaner ».
\end{definition}

L'étymologie n'est jamais une simple curiosité de dictionnaire : elle éclaire le sens profond d'un concept. Deux enseignements se dégagent ici.

D'une part, la violence est d'abord pensée comme un \emph{excès} : \emph{violentia} ne désigne pas la force en général, mais la force emportée, débordante, qui franchit une limite. Le violent est celui qui ne se maîtrise plus, dont la puissance échappe à toute mesure. D'autre part, le lien avec \emph{violare} indique que la violence comporte toujours une dimension d'\emph{atteinte} : on viole un serment, une sépulture, une personne. La violence franchit une frontière — celle du corps d'autrui, de sa volonté, de sa dignité — qui aurait dû rester inviolable.

\begin{savaistu}[Une énergie dévoyée]
Le mot « violence » partage sa racine latine \emph{vis} avec « vie » (\emph{vita} dérive de la même famille indo-européenne) et « vigueur ». Cela suggère une intuition profonde : la violence n'est pas une énergie étrangère à la vie, mais une énergie vitale \emph{dévoyée}, une force qui, au lieu de servir à construire, à protéger ou à vivre, se retourne contre autrui pour nuire et détruire. La question philosophique sera donc de comprendre comment cette énergie peut être orientée autrement.
\end{savaistu}

\courssub{Définition générale de la violence}

Partons d'une intuition simple. Lors d'un match de football à Yaoundé, un joueur bouscule son adversaire pour prendre le ballon : personne ne parle de violence. Mais si ce même joueur, furieux, frappe délibérément son adversaire au visage, tout le monde crie à la violence. Quelle différence ? Dans le premier cas, la force est réglée par des règles et ne vise pas à nuire ; dans le second, la force est employée \emph{pour blesser}, contre la volonté de la victime. C'est cette observation qui permet de formuler la définition.

\begin{definition}[La violence]
La \cle{violence} est l'usage de la force destiné à \textbf{contraindre}, \textbf{blesser} ou \textbf{détruire} autrui, contre sa volonté. Elle se caractérise par trois éléments :
\begin{enumerate}
\item un \textbf{usage de la force} (physique, mais aussi psychologique ou verbale) ;
\item une \textbf{intention de nuire} ou de contraindre ;
\item une \textbf{atteinte à la volonté d'autrui}, qui subit au lieu de consentir.
\end{enumerate}
\end{definition}

Analysons chacun de ces trois éléments. Premièrement, il n'y a violence que là où une force s'exerce : frapper, enfermer, humilier, menacer. Mais la force ne suffit pas. Deuxièmement, il faut une visée : le chirurgien qui opère exerce une force sur le corps, mais pour soigner, non pour nuire ; il n'y a pas violence. Troisièmement, et c'est essentiel, la violence s'impose \emph{contre la volonté} de celui qui la subit : elle nie l'autre comme être libre capable de consentir ou de refuser. C'est pourquoi le consentement change tout : deux boxeurs qui montent sur le ring acceptent les coups dans le cadre d'une règle ; la bagarre dans la cour de récréation, elle, est une violence parce qu'elle s'impose à celui qui ne l'a pas choisie.

\begin{exemplebox}[Observer pour définir]
Comparons trois situations observables dans la vie courante :
\begin{itemize}
\item Un père retire de force son enfant de la chaussée au moment où passe une moto : force exercée, mais pour protéger, et dans l'intérêt de l'enfant.
\item Un élève en pousse un autre pour lui prendre son argent de poche : force exercée pour contraindre, contre la volonté de la victime.
\item Un professeur impose le silence pendant une évaluation : contrainte, mais ordonnée par une règle acceptée de tous.
\end{itemize}
Seule la deuxième situation réunit les trois éléments de la définition : force, intention de nuire, négation de la volonté d'autrui. C'est donc la seule qui mérite pleinement le nom de violence.
\end{exemplebox}

\courssub{Violence et force : une distinction fondamentale}

L'erreur la plus répandue consiste à confondre violence et force. Or toute violence est un usage de la force, mais tout usage de la force n'est pas une violence.

\begin{definition}[La force]
La \cle{force} est la capacité d'agir physiquement ou moralement sur les choses et sur les personnes. Elle peut être \textbf{ordonnée par une règle} (la loi, le règlement, le contrat) et mise au service d'une fin légitime : protéger, faire respecter un droit, construire.
\end{definition}

La distinction peut se formuler ainsi : la \textbf{force} est neutre en elle-même ; c'est son \emph{usage} qui la qualifie. La force devient \cle{violence} lorsqu'elle est \textbf{abusive} (excessive, disproportionnée), \textbf{illégitime} (sans fondement dans une règle juste) ou \textbf{destructrice} (visant à nuire pour nuire). Le gendarme qui arrête un voleur en suivant la procédure use de la force ; le même gendarme qui frapperait un suspect menotté commettrait une violence. Le philosophe Blaise Pascal résumait déjà le problème en notant que « \emph{La justice sans la force est impuissante ; la force sans la justice est tyrannique} » (\emph{Pensées}, Lafuma 312) : la force a besoin d'être ordonnée par le droit pour ne pas basculer dans la violence.

\begin{attention}[Erreur fréquente]
Ne jamais affirmer dans une copie que \emph{toute contrainte est une violence}. C'est faux : la loi contraint sans être nécessairement violente. L'obligation de porter le casque à moto, l'interdiction de fumer dans les salles de classe, l'impôt : toutes ces règles contraignent nos actions, mais elles le font au nom d'un intérêt général et selon des procédures acceptées. Une contrainte ne devient violence que si elle est injuste, arbitraire, excessive ou destructrice. Confondre contrainte et violence, c'est rendre impossible toute vie en société, car aucune société ne peut exister sans règles.
\end{attention}

On peut représenter cette distinction sur un axe gradué : plus l'usage de la force s'écarte de la règle et de la mesure, plus il glisse vers la violence destructrice.

\begin{popfigure}
\begin{center}
\begin{tikzpicture}[x=1cm,y=1cm]
% Axe principal
\draw[line width=2.5pt, popgreen!60!popblue] (0,0) -- (5,0);
\draw[line width=2.5pt, poporange] (5,0) -- (8,0);
\draw[line width=2.5pt, poppink] (8,0) -- (12,0);
% Graduations
\foreach \x in {0,2,4,6,8,10,12} \draw[popdark] (\x,-0.15) -- (\x,0.15);
% Étiquettes des zones
\node[above, popgreen!50!black, font=\bfseries] at (2.5,0.35) {Force légitime};
\node[above, poporange!80!black, font=\bfseries] at (6.5,0.35) {Usage abusif};
\node[above, poppink!80!black, font=\bfseries] at (10,0.35) {Violence destructrice};
% Exemples placés
\node[below, align=center, font=\small] at (1.5,-0.25) {Arrestation\\ policière régulière};
\node[below, align=center, font=\small] at (5.5,-0.25) {Gifle\\ (punition excessive)};
\node[below, align=center, font=\small] at (8.8,-0.25) {Lynchage};
\node[below, align=center, font=\small] at (11,-0.25) {Guerre};
% Points repères
\fill[popdark] (1.5,0) circle (2pt);
\fill[popdark] (5.5,0) circle (2pt);
\fill[popdark] (8.8,0) circle (2pt);
\fill[popdark] (11,0) circle (2pt);
% Flèche de glissement
\draw[-{Stealth}, popmuted, dashed] (2.5,1.1) to[bend left=15] node[above, font=\small\itshape]{perte de mesure et de légitimité} (10,1.1);
\end{tikzpicture}
\end{center}
\legende{De la force légitime à la violence destructrice : c'est l'usage de la force (sa légitimité, sa proportion, sa finalité) qui fait basculer l'un dans l'autre.}
\end{popfigure}

Ce schéma montre qu'il n'existe pas de frontière absolue mais un \emph{glissement progressif} : la même force physique (immobiliser quelqu'un, frapper) change de statut moral selon qu'elle est réglée et proportionnée ou arbitraire et destructrice. C'est pourquoi les sociétés cherchent à encadrer l'usage de la force par le droit : l'État, disait le sociologue Max Weber, revendique le monopole de la violence physique \emph{légitime} — autrement dit, il transforme la violence en force ordonnée par la loi.

\courssub{Violence et agressivité : de la pulsion à l'acte}

Deuxième confusion à dissiper : celle de la violence avec l'\cle{agressivité}. Or ces deux notions n'appartiennent pas au même registre.

\begin{definition}[Agressivité et violence]
L'\cle{agressivité} est une \textbf{pulsion}, une tendance intérieure, un état affectif (colère, hostilité, énergie d'affrontement) que l'on retrouve chez l'animal comme chez l'homme. La \cle{violence}, elle, est un \textbf{acte} : le passage à l'acte destructeur ou contraignant contre autrui.
\end{definition}

Cette distinction est capitale pour deux raisons. D'abord, l'agressivité n'est pas en soi condamnable : elle peut être canalisée et devenir positive. Le sportif qui lutte pour gagner, l'élève qui s'acharne sur un exercice difficile, l'entrepreneur qui affronte les obstacles mobilisent tous une énergie agressive, mais sans violence. Le sport de compétition est d'ailleurs un remarquable dispositif social de transformation de l'agressivité : le combat est réglé, limité dans le temps, surveillé par un arbitre, et se termine par une poignée de main. Ensuite, on peut ressentir une vive agressivité sans jamais passer à l'acte violent : l'être humain, précisément parce qu'il est capable de réflexion et de maîtrise de soi, peut retenir la pulsion. À l'inverse, certaines violences sont froides, calculées, sans agressivité apparente — pensons à une dénonciation calomnieuse préparée tranquillement. Pulsion et acte ne coïncident donc pas nécessairement.

\begin{exemplebox}[Pulsion retenue, pulsion déchaînée]
Deux élèves se disputent la même place dans le car qui les conduit au lycée de Ngaoundéré. Tous deux éprouvent la même colère, la même pulsion agressive. Le premier serre les poings, respire et propose de se partager la place : l'agressivité n'est pas devenue violence. Le second frappe son camarade : la pulsion est passée à l'acte. Ce qui distingue l'homme violent de l'homme agressif, c'est précisément ce qui fait l'humanité de l'homme : la capacité de suspendre la pulsion par la volonté et la parole.
\end{exemplebox}

\courssub{La violence comme négation d'autrui}

Pourquoi la violence est-elle moralement condamnable ? Pour le comprendre, il faut approfondir ce qu'elle fait à autrui. La violence ne se contente pas de faire mal : elle \textbf{nie l'autre comme personne}.

La tradition philosophique, notamment depuis Emmanuel Kant, définit la \cle{dignité} humaine : chaque personne a une valeur absolue et doit toujours être traitée \emph{en même temps comme une fin}, jamais simplement comme un moyen. Or que fait la violence ? Elle traite autrui précisément comme un moyen, un obstacle ou une chose :

\begin{itemize}
\item elle réduit autrui au \textbf{silence} : celui qui frappe ou menace empêche l'autre de parler, de répondre, d'argumenter ;
\item elle réduit autrui à l'\textbf{objet} : le corps de la victime devient une chose que l'on manipule, que l'on enferme, que l'on utilise ;
\item elle réduit autrui à la \textbf{chose} : dans les formes extrêmes (esclavage, torture, lynchage), la victime n'est même plus reconnue comme un être humain, mais traitée comme une marchandise ou un animal.
\end{itemize}

La violence est donc, dans son essence, une \textbf{négation de la dignité} : elle refuse à autrui le statut de sujet libre et parlant pour le rabattre au rang de corps subissant. C'est pourquoi même la violence « légère » — l'humiliation, l'insulte — contient déjà ce mouvement de réification : traiter quelqu'un comme s'il ne comptait pas.

\begin{aretenir}[L'essentiel de la définition]
La violence est l'usage de la force pour contraindre, blesser ou détruire autrui contre sa volonté. Elle se distingue de la force (qui peut être légitime et ordonnée) et de l'agressivité (qui est une pulsion, non un acte). Sa gravité tient à ce qu'elle nie autrui comme personne : elle le réduit au silence, à l'objet, à la chose, et marque l'échec du langage.
\end{aretenir}

\courssub{La violence comme échec du langage}

Il reste un dernier trait, peut-être le plus profond. Entre hommes, il existe deux manières de régler un différend : la \textbf{parole} ou la \textbf{force}. Quand deux personnes discutent, négocient, argumentent, elles se reconnaissent mutuellement comme des êtres capables de raisonner : la parole est le règne de la réciprocité. La violence, au contraire, commence exactement là où la parole échoue ou est refusée. On peut formuler cette idée par une formule à retenir : \emph{là où la parole cesse, la violence commence}.

Cela éclaire à la fois la nature de la violence et les moyens de la combattre. Celui qui frappe avoue en réalité une impuissance : il n'a pas su — ou pas voulu — convaincre. La violence est l'argument de celui qui n'a plus d'arguments. Inversement, toutes les institutions humaines qui organisent la parole — le tribunal où l'on plaide au lieu de se venger, la palabre traditionnelle sous l'arbre où les anciens règlent les conflits du village, le débat parlementaire, la médiation scolaire — sont des machines à \emph{substituer la parole à la violence}. Le droit, en particulier, peut être défini comme l'ensemble des procédures par lesquelles la société transforme les vengeances privées en débats publics réglés.

\begin{exoresolu}[Analyse d'une situation concrète]
\textbf{Énoncé.} Dans un quartier de Douala, deux familles se disputent depuis des mois la limite de leurs terrains. Un soir, les fils de l'une des familles détruisent la clôture de l'autre et blessent un voisin venu s'interposer. En utilisant la définition de la violence construite dans cette leçon, montrez en quoi cette situation relève de la violence, et identifiez ce qui aurait pu l'éviter.
\tcblower
\textbf{Solution détaillée.} Appliquons les trois éléments de la définition. (1) \emph{Usage de la force} : la destruction de la clôture et les coups portés au voisin sont des usages physiques de la force. (2) \emph{Intention de contraindre et de nuire} : il s'agit d'imposer par la force une délimitation de terrain et de blesser celui qui s'y oppose. (3) \emph{Atteinte à la volonté d'autrui} : la famille voisine n'a pas consenti ; elle subit. Les trois éléments sont réunis : c'est bien une violence, et non un simple usage légitime de la force (défendre sa propriété par voie de droit serait légitime ; la détruire soi-même par la force ne l'est pas). On observe aussi la négation d'autrui : le voisin est réduit à un obstacle que l'on écarte, et la parole (la négociation, le recours au chef de quartier ou au tribunal) a été court-circuitée. Ce qui aurait pu éviter la violence : la saisine d'une instance de parole — la palabre, la médiation, le juge — c'est-à-dire le remplacement du règlement par la force par un règlement par le droit. La situation illustre parfaitement la formule : la violence commence là où la parole a cessé.
\end{exoresolu}

\begin{attention}[Piège de rédaction]
Dans une dissertation, ne définissez jamais la violence uniquement par la force physique (« taper », « tuer »). Une telle définition est trop étroite : elle exclut les violences psychologiques et verbales (humiliation, menace, insulte), qui contraignent et blessent sans toucher le corps, et elle ne permet pas de distinguer la force légitime de la violence. Une bonne définition doit toujours articuler les trois éléments : force, intention de nuire, négation de la volonté d'autrui.
\end{attention}

Nous disposons désormais d'une définition rigoureuse : la violence est un usage abusif ou destructeur de la force, qui nie autrui comme personne et marque l'échec de la parole. Mais d'où vient-elle ? L'homme est-il violent par nature, ou est-ce la société qui le rend violent ? C'est la question que nous devons maintenant examiner en étudiant l'origine de la violence.

\coursec{L'origine de la violence : naturelle ou culturelle ?}

Nous avons défini la violence comme l'usage de la force pour contraindre, blesser ou détruire autrui, contre sa volonté. Mais une fois cette définition posée, une question plus profonde surgit : \emph{d'où vient la violence ?} L'agresseur qui frappe, le voleur qui agresse, le groupe qui lynche : agissent-ils par nature, c'est-à-dire parce que l'être humain porterait en lui une disposition innée à la violence ? Ou bien agissent-ils sous l'effet de la société, de la misère, de l'injustice, d'une éducation défaillante ? Autrement dit : la violence est-elle \cle{naturelle} (inscrite dans l'homme) ou \cle{culturelle} (produite par les conditions sociales) ?

Cette question n'est pas théorique seulement : elle décide de nos réponses pratiques. Si la violence est naturelle, il faut la contenir par la force et la punition. Si elle est culturelle, il faut transformer la société. C'est exactement le débat qui oppose deux grands philosophes : Thomas \cle{Hobbes} et Jean-Jacques \cle{Rousseau}.

\courssub{La thèse de la violence naturelle : Hobbes}

\textbf{L'intuition de départ.}\par
Observez un marché sans aucun policier, sans règle, sans arbitre : chacun veut la meilleure place, le meilleur prix, et la dispute éclate vite. Hobbes part de cette intuition simple : livrés à eux-mêmes, sans autorité commune, les hommes entrent en conflit, parce qu'ils désirent les mêmes biens et se méfient les uns des autres.

\begin{definition}[L'état de nature]
L'\cle{état de nature} est la situation \emph{hypothétique} des hommes avant toute société, toute loi et toute autorité politique. C'est une fiction philosophique, une expérience de pensée : on imagine ce que serait l'homme « sans État » pour comprendre pourquoi l'État existe et à quoi il sert.
\end{definition}

\textbf{Le raisonnement de Hobbes.}\par
Dans le \emph{Léviathan} (1651), Hobbes décrit l'état de nature de la manière suivante. Les hommes y sont à peu près \emph{égaux} en force et en capacité de nuire : même le plus faible peut tuer le plus fort, par ruse ou en s'associant à d'autres. De cette égalité naît la \emph{diffidence} (la méfiance réciproque) : chacun craint l'autre. Or les hommes désirent les mêmes choses rares (nourriture, terres, honneurs). Trois causes principales de conflit apparaissent alors :
\begin{itemize}
\item la \emph{compétition} : attaquer pour gagner les biens d'autrui ;
\item la \emph{méfiance} : attaquer pour se protéger par anticipation ;
\item la \emph{gloire} : attaquer pour la réputation et le respect.
\end{itemize}
La conséquence est terrible : l'état de nature est une \cle{guerre de tous contre tous} (\emph{bellum omnium contra omnes}). Hobbes résume sa thèse par une formule célèbre : \emph{« l'homme est un loup pour l'homme »} (\emph{homo homini lupus}). Dans une telle situation, « la vie de l'homme est solitaire, misérable, dangereuse, animale et brève ». La violence n'est donc pas un accident : elle est \emph{naturelle}, elle jaillit spontanément dès que les hommes vivent ensemble sans loi.

\textbf{La conséquence politique.}\par
Si la violence est naturelle, on ne peut pas la supprimer : on peut seulement la \emph{contenir}. Pour sortir de la guerre de tous contre tous, les hommes concluent un \cle{contrat} : chacun renonce à son droit illimité de tout faire et transfère sa force à un pouvoir unique et absolu, l'État, que Hobbes nomme le \emph{Léviathan} (monstre biblique gigantesque qui inspire la crainte). Seule la \emph{crainte} d'un pouvoir plus fort que tous peut maintenir la paix. Chez Hobbes, la réponse à la violence est donc : un \textbf{État fort}, des lois sévères, des sanctions dissuasives.

\courssub{La thèse de la violence culturelle : Rousseau}

\textbf{L'intuition de départ.}\par
Un enfant qui naît n'est ni bon ni méchant : il ne connaît ni la propriété, ni l'argent, ni la comparaison sociale. Rousseau part de cette intuition inverse de celle de Hobbes : ce n'est pas l'homme qui est mauvais, c'est la société qui le rend mauvais.

\textbf{Le raisonnement de Rousseau.}\par
Dans le \emph{Discours sur l'origine et les fondements de l'inégalité parmi les hommes} (1755), Rousseau affirme : « L'homme est né bon, et c'est la société qui le corrompt. » Dans l'état de nature, l'homme est un être simple, guidé par deux sentiments naturels : l'\emph{amour de soi} (instinct de conservation) et la \emph{pitié} (répugnance naturelle à voir souffrir un être semblable). Le « sauvage » n'a aucune raison d'attaquer autrui : ses besoins sont simples et la nature y pourvoit.

Comment la violence naît-elle alors ? Rousseau raconte le processus en trois étapes :
\begin{enumerate}
\item \textbf{La propriété.} « Le premier qui, ayant enclos un terrain, s'avisa de dire : \emph{ceci est à moi}, et trouva des gens assez simples pour le croire, fut le vrai fondateur de la société civile. » Avec la propriété privée naissent le vol, la jalousie, la défense armée de ses biens.
\item \textbf{L'inégalité.} La propriété crée des riches et des pauvres ; l'inégalité engendre le ressentiment, la convoitise et la domination des uns sur les autres.
\item \textbf{La comparaison.} Vivant en société, l'homme cesse de vivre pour lui-même : il vit dans le regard des autres. L'\emph{amour-propre} (désir d'être reconnu, d'être supérieur) remplace l'amour de soi. De la comparaison naissent l'orgueil, la haine et la violence.
\end{enumerate}
La violence n'est donc pas dans le cœur de l'homme : elle est un \emph{produit de l'histoire sociale}, un mal \cle{culturel}.

\textbf{La conséquence politique.}\par
Si la violence vient de la société, la solution n'est pas d'abord de punir plus fort, mais de \emph{réformer la société} : réduire les inégalités, éduquer les hommes (c'est le projet de l'\emph{Émile}), fonder un pacte social juste où la volonté générale remplace la domination des puissants.

\begin{attention}[Erreur fréquente]
Ne présentez pas Hobbes et Rousseau comme deux auteurs simplement « opposés » et juxtaposés. Ils répondent à \emph{la même question} — que serait l'homme sans société, et d'où vient la violence ? — et utilisent \emph{la même méthode} : la fiction de l'état de nature. Leur désaccord porte sur le \emph{contenu} de cet état (guerre naturelle chez Hobbes, paix naturelle chez Rousseau) et donc sur le \emph{rôle de l'État} (contenir la violence chez Hobbes, corriger la société chez Rousseau). En dissertation, montrez d'abord ce point commun, puis l'opposition.
\end{attention}

\begin{popfigure}
\begin{tikzpicture}[font=\small]
% Titre colonnes
\node[fill=popblue, text=white, rounded corners=2pt, minimum width=6.6cm, minimum height=0.7cm, anchor=north west] at (0,0) {\textbf{Hobbes : violence naturelle}};
\node[fill=poporange, text=white, rounded corners=2pt, minimum width=6.6cm, minimum height=0.7cm, anchor=north west] at (7,0) {\textbf{Rousseau : violence culturelle}};
% Ligne thèse
\node[draw=popline, fill=popblueL, minimum width=6.6cm, minimum height=2.1cm, anchor=north west, align=left, text width=6.1cm] at (0,-0.8) {\textbf{Thèse :} l'homme est un loup pour l'homme ; l'état de nature est guerre de tous contre tous.};
\node[draw=popline, fill=poporangeL, minimum width=6.6cm, minimum height=2.1cm, anchor=north west, align=left, text width=6.1cm] at (7,-0.8) {\textbf{Thèse :} l'homme naît bon (amour de soi, pitié) ; c'est la société qui le corrompt.};
% Ligne argument
\node[draw=popline, fill=white, minimum width=6.6cm, minimum height=2.5cm, anchor=north west, align=left, text width=6.1cm] at (0,-3.0) {\textbf{Argument :} égalité des forces, méfiance, désir des mêmes biens rares $\Rightarrow$ conflit spontané.};
\node[draw=popline, fill=white, minimum width=6.6cm, minimum height=2.5cm, anchor=north west, align=left, text width=6.1cm] at (7,-3.0) {\textbf{Argument :} propriété, inégalité, comparaison (amour-propre) $\Rightarrow$ violence produite par l'histoire sociale.};
% Ligne conséquence
\node[draw=popline, fill=popblueL, minimum width=6.6cm, minimum height=2.3cm, anchor=north west, align=left, text width=6.1cm] at (0,-5.6) {\textbf{Conséquence politique :} un État fort (le Léviathan), la crainte des lois et des sanctions pour contenir la violence.};
\node[draw=popline, fill=poporangeL, minimum width=6.6cm, minimum height=2.3cm, anchor=north west, align=left, text width=6.1cm] at (7,-5.6) {\textbf{Conséquence politique :} réformer la société : réduire les inégalités, éduquer, fonder un pacte juste.};
\end{tikzpicture}
\legende{Tableau comparatif : Hobbes et Rousseau face à la même question — l'origine de la violence.}
\end{popfigure}

\begin{aretenir}[Références à retenir]
\begin{itemize}
\item Thomas \cle{Hobbes}, \emph{Le Léviathan} (1651) : l'état de nature est guerre de tous contre tous ; seul un État fort contient la violence naturelle.
\item Jean-Jacques \cle{Rousseau}, \emph{Discours sur l'origine de l'inégalité} (1755) : l'homme naît bon ; la propriété, l'inégalité et la comparaison engendrent la violence.
\end{itemize}
\end{aretenir}

\courssub{Discussion critique : que disent les sciences ?}

Philosopher, ce n'est pas seulement citer des auteurs : c'est confronter leurs thèses aux faits. Deux types d'arguments scientifiques éclairent le débat.

\textbf{Les arguments éthologiques (en faveur de Hobbes).}\par
L'\cle{éthologie} est la science du comportement animal. Or l'observation montre que l'agressivité existe chez de nombreuses espèces : combats entre mâles pour la reproduction, défense du territoire, compétition pour la nourriture. Chez les chimpanzés, nos plus proches parents, on observe même des raids violents organisés contre des groupes voisins. Cela suggère que l'homme possède, héritées de son évolution, des \emph{dispositions} à l'agressivité : peur, colère, instinct de défense du territoire et du groupe. La violence aurait donc bien une racine \emph{naturelle}.

\textbf{Les arguments sociologiques (en faveur de Rousseau).}\par
Mais l'anthropologie et la sociologie apportent un contre-argument décisif : toutes les sociétés humaines ne sont \emph{pas également violentes}. Les taux d'homicide varient énormément d'un pays à l'autre et d'une époque à l'autre. Certaines sociétés traditionnelles règlent leurs conflits par la palabre, la médiation et la compensation plutôt que par les armes. Si la violence était purement naturelle et constante, elle devrait être partout la même ; or elle varie avec la pauvreté, les inégalités, l'organisation politique, l'éducation. La violence dépend donc aussi, et fortement, du \emph{culturel}.

\begin{propriete}[Position synthétique]
La violence puise dans des \cle{dispositions naturelles} (peur, colère, agressivité héritée de l'animalité), mais ses \emph{formes}, sa \emph{fréquence} et son \emph{intensité} dépendent des \cle{conditions sociales} : inégalités, injustice, éducation, efficacité des institutions. La nature fournit la poudre ; la société fournit — ou retire — l'étincelle.
\end{propriete}

Cette synthèse dépasse le faux dilemme « tout naturel ou tout culturel » : elle reconnaît à Hobbes que l'homme n'est pas un ange, et à Rousseau que la société porte une lourde responsabilité dans le déclenchement des violences.

\courssub{Application : la violence urbaine au Cameroun}

Appliquons cette analyse à une situation concrète. Dans les grandes villes camerounaises (Douala, Yaoundé, Bamenda), on observe des phénomènes de violence urbaine : agressions, vols à l'arraché, affrontements entre bandes de jeunes, parfois justice populaire. Comment les expliquer ?

\begin{itemize}
\item \textbf{La pauvreté} : quand les besoins élémentaires (nourriture, logement) ne sont pas satisfaits, la compétition pour les ressources rares s'aiguise — exactement le mécanisme décrit par Hobbes.
\item \textbf{Le chômage des jeunes} : des milliers de jeunes diplômés sans emploi éprouvent un sentiment d'inutilité et d'humiliation ; la frustration peut se convertir en agressivité.
\item \textbf{Les inégalités} : le contraste visible entre quartiers riches et quartiers défavorisés nourrit le ressentiment — c'est le mécanisme de la comparaison et de l'amour-propre blessé décrit par Rousseau.
\item \textbf{Le sentiment d'injustice} : quand les institutions semblent faibles ou partiales, chacun est tenté de « se faire justice soi-même », ce qui ramène à une situation proche de l'état de nature.
\end{itemize}

\begin{exemplebox}[Le débat camerounais sur la délinquance juvénile]
Face à la délinquance juvénile, deux discours s'affrontent dans les débats publics camerounais. Les uns demandent de \emph{punir plus sévèrement} : plus de policiers, des peines plus lourdes, plus de fermeté — c'est la logique de \textbf{Hobbes} (contenir par la crainte). Les autres demandent de \emph{réformer la société} : créer des emplois, des centres de formation, des terrains de sport, lutter contre l'exclusion — c'est la logique de \textbf{Rousseau} (supprimer les causes sociales). La position synthétique invite à combiner les deux : une justice ferme mais équitable \emph{et} des politiques sociales qui attaquent les racines du mal.
\end{exemplebox}

\begin{popfigure}
\begin{center}
\begin{tikzpicture}
\begin{axis}[
    ybar,
    width=0.95\linewidth,
    height=6.5cm,
    bar width=0.9cm,
    ymin=0, ymax=50,
    ylabel={Pourcentage de réponses (\%)},
    symbolic x coords={Pauvreté,Injustice,Éducation,Nature humaine},
    xtick=data,
    x tick label style={font=\small, align=center},
    nodes near coords,
    nodes near coords style={font=\small\bfseries},
    axis line style={draw=popline},
    tick style={draw=popline},
    ymajorgrids=true,
    grid style={popline!40, dashed},
]
\addplot[fill=popblue, draw=popdark] coordinates {(Pauvreté,38) (Injustice,27) (Éducation,21) (Nature humaine,14)};
\end{axis}
\end{tikzpicture}
\end{center}
\legende{Diagramme en barres : perception des causes de la violence chez des jeunes (données fictives à titre pédagogique). La majorité des réponses désigne des causes \emph{sociales} (pauvreté, injustice, éducation), ce qui conforte la thèse de Rousseau et la position synthétique.}
\end{popfigure}

Ce diagramme fictif illustre un point essentiel : interrogés, les jeunes eux-mêmes attribuent massivement la violence à des causes \emph{sociales} plutôt qu'à la « nature humaine ». Cela ne prouve pas que Rousseau a entièrement raison, mais cela montre que les conditions sociales jouent un rôle décisif dans le déclenchement des violences.

\begin{exoresolu}[Analyse d'une situation concrète]
\textbf{Énoncé.} Dans un quartier de Yaoundé, un jeune de 19 ans, sans emploi depuis la fin de ses études, participe à un vol à main armée. Interpellé, il déclare : « Personne ne m'a donné ma chance ; j'ai pris ce qu'on me refusait. » Analysez cette situation à la lumière de Hobbes, puis de Rousseau, puis proposez une conclusion synthétique.
\tcblower
\textbf{Solution détaillée.}
\begin{enumerate}
\item \emph{Analyse hobbesienne.} Hobbes dirait que ce jeune illustre la tendance naturelle de l'homme à prendre par la force ce qu'il désire dès que la contrainte fait défaut. La réponse adaptée est le renforcement de l'autorité de l'État : la loi doit être appliquée pour que la crainte de la sanction dissuade les agresseurs et protège les victimes.
\item \emph{Analyse rousseauiste.} Rousseau dirait que ce jeune n'est pas né violent : c'est la société — le chômage, l'inégalité, le sentiment d'injustice — qui l'a corrompu et poussé au crime. Sa déclaration (« on me refusait ») exprime l'amour-propre blessé par la comparaison sociale. La vraie solution est de réformer la société : emploi, formation, justice équitable.
\item \emph{Conclusion synthétique.} La disposition à la violence existe en tout homme (Hobbes), mais c'est la misère et l'injustice qui l'ont ici transformée en acte (Rousseau). Il faut donc à la fois une réponse judiciaire ferme (le vol reste un acte condamnable) et une action sociale de fond (insertion des jeunes) pour prévenir la récidive.
\end{enumerate}
\end{exoresolu}

\begin{exercice}[Pour s'entraîner]
Un débat est organisé dans votre classe sur le thème : « Faut-il plus de policiers ou plus d'emplois pour réduire la violence en ville ? » Rédigez un paragraphe argumenté (10 à 15 lignes) dans lequel vous : 1) présentez la thèse de Hobbes ; 2) présentez la thèse de Rousseau ; 3) défendez une position personnelle synthétique en vous appuyant sur un exemple camerounais.
\end{exercice}

\begin{corrige}[Exercice : éléments de réponse attendus]
Une bonne copie doit : (1) rappeler que pour Hobbes, la violence étant naturelle, seule la crainte d'un État fort la contient — d'où la demande de plus de forces de l'ordre ; (2) rappeler que pour Rousseau, la violence étant produite par les inégalités et l'exclusion, il faut d'abord des emplois et de la justice sociale ; (3) conclure par une synthèse argumentée : la fermeté judiciaire est nécessaire à court terme, mais seules des politiques sociales (emploi, éducation, réduction des inégalités) traitent les racines de la violence. L'exemple camerounais (chômage des jeunes, quartiers défavorisés, débats sur la délinquance) doit être précis et relié explicitement à l'argumentation.
\end{corrige}

Ainsi, la question de l'origine de la violence ne se résout pas par un choix brutal entre nature et culture : l'homme hérite de dispositions agressives, mais c'est la société qui décide si ces dispositions restent endormies ou éclatent en violences. Reste à savoir si la violence, quelle que soit son origine, peut parfois être justifiée : ce sera l'objet du procès de la violence.

\coursec{Les types de violence}

Après avoir interrogé l'origine de la violence — jaillit-elle de notre nature biologique ou est-elle le fruit de nos constructions culturelles ? — nous devons désormais en dresser l'inventaire. Car parler de « la violence » au singulier est un piège : la violence est \cle{polymorphe}. Elle change de visage, d'intensité et de portée selon qu'elle frappe les corps, les esprits, les conditions de vie ou les structures sociales. La distinguer, la nommer, la classer, c'est déjà commencer à la combattre. C'est l'objet de cette section : passer de la question « d'où vient-elle ? » à la question « quels sont ses visages ? ».

\courssub{La violence physique : l'atteinte à l'intégrité corporelle}

C'est la forme la plus immédiate, la plus visible, celle qui saute aux yeux et dont le vocabulaire est le plus partagé : coups, blessures, mutilations, meurtre, torture, guerre, châtiments corporels. Elle se définit par l'usage de la force matérielle contre le corps d'autrui pour lui causer de la douleur, des dommages ou la mort.

\begin{definition}[Violence physique]
Emploi de la force corporelle ou d'un instrument matériel contre une personne, intentionnellement, dans le but de lui infliger une souffrance, une blessure, une incapacité ou la mort. Elle laisse des traces objectives (ecchymoses, fractures, cicatrices) ou aboutit à la suppression de la vie.
\end{definition}

Au Cameroun, la violence physique se manifeste tragiquement dans les conflits armés (régions du Nord-Ouest et du Sud-Ouest, Extrême-Nord face aux attaques de Boko Haram), mais aussi dans le quotidien : violences conjugales, bagarres, châtiments corporels encore pratiqués dans certaines familles ou écoles malgré leur condamnation, brutalités lors d'interpellations. La guerre, comme le note le théoricien militaire Carl von Clausewitz, est « la continuation de la politique par d'autres moyens » : elle est la violence physique organisée, portée à son paroxysme.

\courssub{La violence psychologique et verbale : l'atteinte à la dignité et à l'esprit}

Moins spectaculaire, souvent invisible, elle n'en est pas moins dévastatrice. Elle vise non le corps mais la subjectivité : l'estime de soi, la réputation, la santé mentale, la liberté de pensée. Insultes, humiliations, menaces, intimidation, chantage affectif, dénigrement systématique, isolement forcé, manipulation faisant douter la victime de sa propre perception de la réalité.

\begin{exemplebox}[Le harcèlement scolaire et le cyberharcèlement entre élèves]
Dans un lycée de Douala, une élève de Première, appelons-la Marie, subit depuis des mois des moqueries sur son physique dans la cour de récréation. Le harcèlement a basculé dans le \emph{cyberharcèlement} : un groupe d'élèves a créé un groupe WhatsApp où circulent des montages photo humiliants, des rumeurs sur sa vie privée, des insultes. Marie ne dort plus, ses notes s'effondrent, elle sombre dans la dépression. Aucun coup n'a été porté, pourtant la violence est extrême. L'outil numérique (téléphone, réseaux sociaux) agit comme un \emph{amplificateur} : la violence suit la victime jusque dans sa chambre, à toute heure, devant un public potentiellement illimité. C'est une violence \emph{psychologique} (destruction de l'image de soi) et \emph{verbale} (insultes écrites), aggravée par l'architecture des plateformes (anonymat, viralité, modération insuffisante).
\end{exemplebox}

\begin{attention}[Piège à éviter : l'invisibilité des violences « douces »]
L'erreur la plus fréquente, dans les copies comme dans la vie, est de croire que « ce n'est pas de la violence, il n'y a pas eu de coups ». \textbf{La violence psychologique tue aussi.} Elle tue l'avenir (décrochage scolaire), elle tue l'être (dépression, suicide), elle tue le lien social (méfiance généralisée). Au Probatoire technique, une dissertation qui réduit la violence à sa dimension physique est une dissertation partiellement hors sujet. Il faut \cle{systématiquement} articuler le visible et l'invisible.
\end{attention}

\courssub{La violence économique et sociale : l'exploitation et l'exclusion}

Elle ne frappe pas directement le corps ni l'esprit par la parole, mais par la structure des rapports de production et de distribution. C'est la violence de la misère imposée, de l'exploitation du travail, de l'exclusion de l'accès aux ressources vitales (santé, éducation, logement, crédit). Elle touche les plus vulnérables : travailleurs précaires, enfants, personnes âgées, déplacés internes.

\begin{exemplebox}[Le travail des enfants dans les carrières de sable]
Au Cameroun, des enfants non scolarisés ou déscolarisés extraient du sable à la pelle dans le lit des fleuves ou dans des carrières, sous un soleil de plomb, pour le revendre aux transporteurs du bâtiment. Ils gagnent quelques centaines de francs CFA par jour, inhalent la poussière, risquent la noyade, portent des charges au-delà de leurs forces. Aucun adulte ne les frappe (pas de violence physique directe), aucun ne les insulte (pas de violence verbale). Pourtant, c'est une \textbf{violence économique} (vol de l'enfance, de la scolarité et de la santé pour un revenu de survie) et \textbf{sociale} (abandon, indifférence, reproduction de la pauvreté). L'agresseur n'est pas un individu identifiable mais un système économique qui profite de cette main-d'œuvre bon marché.
\end{exemplebox}

\courssub{La violence symbolique : la domination douce et intériorisée (Bourdieu)}

C'est sans doute la plus insidieuse, car elle obtient le consentement de la victime. Le sociologue Pierre Bourdieu la définit comme une violence qui s'exerce « avec la complicité de ceux qui la subissent ». Elle passe par les catégories de perception, les normes esthétiques, les jugements de goût, le langage, les diplômes. Elle impose les catégories du groupe dominant comme universelles, naturelles, légitimes.

\begin{definition}[Violence symbolique (Pierre Bourdieu)]
Forme de domination qui s'exerce non par la contrainte physique ni par la loi explicite, mais par l'imposition de catégories de pensée, de perception et d'évaluation (langage, culture, goût, diplômes) qui sont celles du groupe dominant. Elle est \emph{invisible} car elle passe par les structures mentales ; elle est \emph{efficace} car la victime l'intériorise et se perçoit elle-même comme « naturellement » inférieure, indigne ou mal élevée.
\end{definition}

\begin{savaistu}[Johan Galtung : le triangle de la violence]
Le sociologue norvégien Johan Galtung, fondateur de l'irénologie (science de la paix), propose une grille de lecture majeure, très utile en dissertation. Il distingue trois niveaux imbriqués :
\begin{itemize}
    \item \textbf{Violence directe :} celle qu'on voit (coups, bombes, insultes). Auteur identifiable.
    \item \textbf{Violence structurelle :} celle qui est \emph{inscrite dans les structures} (inégalités d'accès à la santé, à l'éducation, à la justice). Pas d'agresseur unique, mais des victimes identifiables (mortalité évitable, analphabétisme).
    \item \textbf{Violence culturelle :} celle qui \emph{légitime} les deux précédentes (idéologies racistes, sexistes, discours de mépris). Elle rend la violence directe ou structurelle « normale », « juste », « nécessaire ».
\end{itemize}
Cette distinction est un \cle{outil analytique puissant} pour ne pas rester à la surface des faits divers.
\end{savaistu}

\courssub{La violence politique et étatique : répression, guerre, terrorisme et le monopole weberien}

La politique est l'arène par excellence de la violence organisée : répression des manifestations, guerres civiles, guerres entre États, terrorisme, coups d'État, génocides. Le sociologue allemand Max Weber a posé la définition canonique de l'État moderne par son rapport à la violence.

\begin{definition}[Monopole de la violence physique légitime (Max Weber)]
L'État est la communauté humaine qui, dans les limites d'un territoire déterminé, revendique avec succès pour elle-même le \emph{monopole de la violence physique légitime}. Cela signifie deux choses : (1) seul l'État (police, armée, justice) a le droit d'user de la force physique (arrestation, emprisonnement, guerre défensive) \emph{dans le respect du droit} ; (2) tout usage de la force par un particulier ou un groupe hors du cadre légal est \emph{illégitime} (crime, rébellion, milice).
\end{definition}

Ce monopole n'est pas un blanc-seing. Il est \emph{conditionné} par le droit (État de droit) et par le contrôle des citoyens. Quand l'État torture, tue hors de toute procédure judiciaire ou fait disparaître des opposants, il bascule dans la \emph{violence d'État illégitime}. À l'inverse, la tradition philosophique et juridique reconnaît un \emph{droit de résistance à l'oppression} (inscrit à l'article 2 de la Déclaration des droits de l'homme et du citoyen de 1789) : la question de la légitimité de la violence revendicative sera examinée dans la section consacrée au procès de la violence.

\courssub{La violence structurelle : l'injustice sans agresseur identifiable}

Proche de la violence économique mais plus large, la violence structurelle (concept affiné par Galtung) désigne des souffrances et des morts évitables qui résultent non d'une action violente ponctuelle, mais du fonctionnement « normal » des institutions : lois, marchés, bureaucraties, normes culturelles.

\begin{propriete}[Caractères de la violence structurelle]
\begin{itemize}
    \item \textbf{Invisible :} pas de coup, pas de sang immédiat, pas de cri.
    \item \textbf{Systémique :} elle est produite par la structure sociale elle-même (répartition inégale des richesses, du pouvoir, des soins).
    \item \textbf{Sans auteur unique :} on ne peut pas désigner un « coupable » précis pour la mort d'un enfant de paludisme faute de centre de santé à proximité ; les responsabilités sont diluées (sous-financement, mauvaise gouvernance, inégalités).
    \item \textbf{Massive :} elle tue bien plus, à l'échelle du monde, que la violence directe (famines, maladies évitables, accidents du travail non prévenus).
\end{itemize}
\end{propriete}

\courssub{Classement croisé : croiser les critères pour une analyse fine}

Pour ne pas enfermer une situation dans une seule case, les philosophes et sociologues croisent trois paires de critères. Ce tableau est un \cle{outil méthodologique indispensable} pour l'analyse de documents ou la dissertation.

\begin{center}
\begin{tabularx}{\linewidth}{|l|X|X|}\hline
\rowcolor{popblueL}\textbf{Critère} & \textbf{Pôle A} & \textbf{Pôle B} \\ \hline
\textbf{Directe / Indirecte} (Galtung) & Agression intentionnelle, auteur identifiable, lien causal visible (coup, bombe, insulte). & Violence structurelle ou culturelle : effet retardé, cause diffuse, pas d'auteur unique (inégalité scolaire, norme sexiste). \\ \hline
\textbf{Individuelle / Collective} & Auteur unique ou petit groupe ; victime unique ou petit groupe (agression, harcèlement). & Groupes organisés, foules, États, classes sociales ; victimes massives (guerre, génocide, exploitation, discrimination d'État). \\ \hline
\textbf{Privée / Publique} & Sphère domestique, intime, hors du regard de la loi (violences conjugales, maltraitance parentale). & Sphère politique, espace public, institutionnelle (répression, guerre, lois discriminatoires, violence symbolique médiatique). \\ \hline
\end{tabularx}
\end{center}

\emph{Remarque :} une même situation peut être à la fois collective et privée (violences conjugales banalisées dans une société patriarcale), directe et structurelle (expulsion violente ordonnée par une décision administrative). La richesse de l'analyse réside dans le \textbf{recoupement} des critères.

\begin{popfigure}
\begin{tikzpicture}[
    grow=right,
    level distance=4.6cm,
    level 1/.style={sibling distance=2.6cm},
    level 2/.style={sibling distance=1.6cm, level distance=4.2cm},
    every node/.style={font=\small, align=center},
    edge from parent/.style={draw=popdark, thick, -{Stealth[length=2mm]}}
]
\node[draw=popblue, thick, fill=popblueL, rounded corners, minimum width=2.2cm, align=center]{\textbf{LA}\\\textbf{VIOLENCE}}
    child { node[draw=popgold, thick, fill=popgoldL, rounded corners, align=center]{\textbf{Politique / Étatique}\\répression, guerre,\\monopole (Weber)} }
    child { node[draw=poppink, thick, fill=poppinkL, rounded corners, align=center]{\textbf{Symbolique}\\normes, mépris,\\stigmatisation (Bourdieu)} }
    child { node[draw=poppurple, thick, fill=poppurpleL, rounded corners, align=center]{\textbf{Économique / Sociale}\\exploitation, misère,\\exclusion} }
    child { node[draw=popgreen, thick, fill=popgreenL, rounded corners, align=center]{\textbf{Psychologique / Verbale}\\insultes, menaces,\\harcèlement, cyberharcèlement} }
    child { node[draw=poporange, thick, fill=poporangeL, rounded corners, align=center]{\textbf{Physique}\\coups, blessures,\\guerre, châtiments} };
\end{tikzpicture}
\legende{Classification des grands types de violence. Chaque type peut se combiner avec les autres dans une situation concrète.}
\end{popfigure}

\begin{popfigure}
\begin{tikzpicture}[every node/.style={align=center}]
% Cercle extérieur : violence culturelle
\draw[popgreen, thick, fill=popgreen!10] (0,0) circle (4.6);
\node[font=\normalsize\bfseries, text=popgreen!60!black] at (0,3.9) {VIOLENCE CULTURELLE};
\node[font=\footnotesize, text width=6.5cm] at (0,3.1) {Idéologies et discours qui légitiment\\(racisme, sexisme, mépris social)};
% Cercle intermédiaire : violence structurelle
\draw[poporange, thick, fill=poporange!12] (0,-0.4) circle (3.1);
\node[font=\normalsize\bfseries, text=poporange!70!black] at (0,2.0) {VIOLENCE STRUCTURELLE};
\node[font=\footnotesize, text width=4.6cm] at (0,1.2) {Inégalités inscrites dans les institutions\\(santé, éducation, justice, marché)\\sans agresseur identifiable};
% Cercle central : violence directe
\draw[popblue, thick, fill=popblue!15] (0,-1.2) circle (1.6);
\node[font=\normalsize\bfseries, text=popblue!80!black] at (0,-0.9) {VIOLENCE DIRECTE};
\node[font=\footnotesize, text width=2.6cm] at (0,-1.8) {Agression visible,\\auteur identifiable};
% Flèches d'explication
\draw[popdark, -{Stealth[length=2mm]}, thick] (4.6,1.5) -- (5.6,1.5) node[right, font=\footnotesize, text width=3.6cm] {La culture légitime};
\draw[popdark, -{Stealth[length=2mm]}, thick] (3.1,-2.2) -- (4.4,-3.0) node[right, font=\footnotesize, text width=3.6cm] {La structure produit les conditions};
\draw[popdark, -{Stealth[length=2mm]}, thick] (-1.6,-2.4) -- (-3.2,-3.2) node[left, font=\footnotesize, text width=3.4cm] {L'acte visible n'est que la pointe émergée};
\end{tikzpicture}
\legende{Le triangle de la violence selon Johan Galtung, représenté en cercles imbriqués : la violence culturelle (extérieur) fournit les justifications ; la violence structurelle (milieu) crée les conditions objectives ; la violence directe (centre) est l'acte visible.}
\end{popfigure}

\courssub{Atelier méthodologique : analyser une situation concrète}

Face à un sujet de dissertation, un texte à commenter ou un fait de société, il ne suffit pas de lister les types. Il faut \emph{qualifier} la violence en jeu. Voici une démarche en cinq étapes, à appliquer systématiquement.

\begin{methode}[Analyser et classifier une situation de violence]
\begin{enumerate}
    \item \textbf{Identifier la (les) victime(s) :} qui subit ? Individu ? Groupe ? Classe sociale ? Dans son corps ? Son esprit ? Son avenir ?
    \item \textbf{Identifier l'auteur ou l'agent :} individu identifiable ? Groupe ? Institution ? Structure anonyme ? « Système » ?
    \item \textbf{Identifier le moyen ou le canal :} force physique ? Parole, écrit, image, outil numérique ? Règle juridique ou économique ? Norme culturelle ou symbole ? Absence de soin ou d'école ?
    \item \textbf{Identifier le but ou la fonction (manifeste ou latente) :} domination ? Profit ? Défense ? Intimidation ? Reproduction de l'ordre social ?
    \item \textbf{Classer et croiser :} physique, psychologique, économique, symbolique ou politique ? Directe ou indirecte ? Individuelle ou collective ? Privée ou publique ? Niveau de Galtung (directe, structurelle, culturelle) ?
\end{enumerate}
\end{methode}

\begin{exoresolu}[Application : l'expulsion forcée d'un quartier précaire à Yaoundé]
\textbf{Énoncé :} À 5 h du matin, des bulldozers escortés par des forces de l'ordre détruisent les habitations de fortune d'un quartier précaire pour « assainir la ville ». Les habitants, réveillés en sursaut, n'ont pas le temps de sauver leurs affaires. Ceux qui tentent de s'opposer sont repoussés avec brutalité. Une décision administrative a été signée la veille. Aucun relogement n'est prévu.
\tcblower
\textbf{Analyse par la méthode :}
\begin{enumerate}
    \item \textbf{Victimes :} une collectivité d'habitants pauvres (violence collective), atteinte dans les corps (brutalités), les biens (destructions), la dignité (humiliation publique) et l'avenir (errance).
    \item \textbf{Auteurs :} l'État (autorité administrative, forces de l'ordre) — auteur collectif, institutionnel, identifiable.
    \item \textbf{Moyens :} force physique (bulldozers, brutalités) + acte juridique (décision administrative) + discours légitimant (« insalubrité », « embellissement »).
    \item \textbf{Buts :} manifeste : salubrité publique, aménagement urbain ; latent : récupération foncière, élimination visible de la pauvreté, démonstration de force.
    \item \textbf{Classification croisée :}
    \begin{itemize}
        \item \textbf{Physique} (brutalités, destructions) — \emph{directe} ;
        \item \textbf{Psychologique} (traumatisme, terreur, humiliation) — \emph{directe} ;
        \item \textbf{Économique et sociale} (spoliation de l'habitat et des maigres biens, absence de logement social) — \emph{structurelle} ;
        \item \textbf{Symbolique} (stigmatisation des pauvres comme « sales » ou « illégitimes ») — \emph{culturelle} au sens de Galtung ;
        \item \textbf{Politique et étatique} (usage du monopole de la violence détourné de sa fin protectrice) — \emph{collective et publique}.
    \end{itemize}
\end{enumerate}
\textbf{Conclusion synthétique :} c'est un cas d'école de \textbf{violence d'État} où la violence directe (l'expulsion musclée) est l'instrument de la violence structurelle (exclusion foncière), légitimée par la violence symbolique (discours hygiéniste). Son illégitimité tient à la disproportion des moyens, à l'absence de procédure contradictoire et à l'absence de solution de relogement, contraires au droit au logement et au principe de proportionnalité.
\end{exoresolu}

\begin{aretenir}[Synthèse : la violence, un prisme à multiples facettes]
\begin{itemize}
    \item \textbf{Ne pas réduire la violence au physique.} La violence psychologique (harcèlement, cyberharcèlement), économique (exploitation, misère), symbolique (mépris, stigmatisation — Bourdieu), structurelle (inégalités institutionnelles — Galtung) et politique (répression, monopole weberien) est tout aussi réelle et souvent plus meurtrière à long terme.
    \item \textbf{Croiser les grilles de lecture.} Une situation réelle mobilise presque toujours plusieurs types : directe + structurelle + symbolique ; individuelle + collective ; privée + publique.
    \item \textbf{Le monopole de la violence physique légitime (Weber) définit l'État moderne}, mais il est borné par le droit : son dépassement constitue une violence d'État illégitime.
    \item \textbf{La grille de Galtung (directe / structurelle / culturelle)} est l'outil de référence pour passer de l'anecdote au structurel en dissertation.
    \item \textbf{Méthode d'analyse en cinq étapes :} victime — auteur — moyen — but — classification croisée. C'est la démarche attendue au Probatoire technique pour « appliquer les notions à des situations concrètes ».
\end{itemize}
\end{aretenir}

\coursec{Le procès de la violence (I) : la condamnation}

Après avoir identifié les \cle{types de violence} (physique, psychologique, structurelle, symbolique), nous devons maintenant la juger. La violence s'impose-t-elle comme une fatalité ou comme un échec ? Pourquoi la philosophie, de Kant à Arendt en passant par Gandhi, la condamne-t-elle presque unanimement ? C'est ce que nous allons examiner en instruisant le procès de la violence, non pour la bannir sans comprendre, mais pour saisir pourquoi elle détruit ce qu'elle prétend défendre.

\courssub{La violence comme négation de la dignité humaine : l'argument kantien}

\begin{definition}[L'homme comme fin en soi]
Selon \textbf{Immanuel Kant} (\emph{Fondements de la métaphysique des mœurs}, 1785), tout être rationnel existe comme une \cle{fin en soi} et non comme un simple moyen. La \cle{dignité} est cette valeur intrinsèque, incomparable, qui interdit de traiter autrui comme un objet ou un instrument.
\end{definition}

\emph{Intuition :} lorsque je frappe quelqu'un pour qu'il obéisse, je ne m'adresse pas à sa raison, je contourne sa volonté. Je le réduis à un corps qui réagit à la douleur, non à un sujet qui comprend et choisit.

\begin{propriete}[La violence nie l'humanité dans la personne]
La violence \textbf{traite l'autre comme une chose} (un moyen) et non comme une \cle{personne} (une fin). Elle substitue la contrainte physique ou psychologique au dialogue rationnel. Pour Kant, même si la violence visait un « bon » résultat (ex. : frapper un enfant « pour son bien »), elle reste moralement illégitime car elle piétine l'autonomie morale de la victime.
\end{propriete}

\begin{exemplebox}[Exemple concret : le travail forcé vs le contrat de travail]
Un employeur qui paie un salaire respecte le travailleur comme une fin : il reconnaît sa liberté de consentir. Un employeur qui utilise la menace ou la contrainte physique (travail forcé, esclavage moderne) traite l'humain comme un outil. Au Cameroun, la lutte contre le \emph{travail des enfants} dans les plantations de cacao ou les carrières repose sur ce principe : l'enfant n'est pas une main-d'œuvre bon marché, mais un sujet de droits en devenir.
\end{exemplebox}

\begin{attention}[Piège à éviter]
Ne confondez pas \emph{contrainte légitime} (loi, police, sanction judiciaire) et \emph{violence arbitraire}. L'État de droit monopolise la contrainte physique pour \emph{protéger} la dignité (empêcher le meurtre, le vol), non pour l'anéantir. La distinction tient au \textbf{respect de la procédure} et à la \textbf{proportionnalité}.
\end{attention}

\courssub{La violence comme échec de la parole et de la raison : Hannah Arendt}

\begin{aretenir}[Hannah Arendt, \emph{Du mensonge à la violence} (1970)]
\begin{itemize}
\item \textbf{La violence est muette :} elle ne persuade pas, elle contraint. « La violence peut détruire le pouvoir, mais elle ne peut jamais le créer. »
\item \textbf{Le pouvoir} naît du \cle{consentement} libre des hommes agissant de concert (parole, persuasion, accord).
\item \textbf{La violence} repose sur des \cle{instruments} (armes, matraques, argent, technologie) qui multiplient la force naturelle.
\item \emph{Conséquence :} celui qui n'a plus que la violence pour se faire obéir a déjà perdu le pouvoir.
\end{itemize}
\end{aretenir}

\emph{Démarche d'investigation :} observez un conflit familial, scolaire ou politique.
\begin{enumerate}
\item Quand le dialogue fonctionne, l'accord repose sur des \textbf{raisons} acceptées librement $\rightarrow$ c'est du \cle{pouvoir}.
\item Quand l'un impose sa volonté par la menace, le cri, le coup ou l'exclusion $\rightarrow$ c'est de la \cle{violence}.
\item Résultat : la victime obéit \emph{de force}, mais sa \emph{conviction} n'a pas changé. Le conflit est \emph{gelé}, non \emph{résolu}.
\end{enumerate}

\begin{exemplebox}[Exemple : la grève et la négociation]
Des ouvriers revendiquent une augmentation.
\begin{itemize}
\item \textbf{Voie du pouvoir :} ils débattent, font grève (action concertée), négocient, signent un accord. L'employeur consent $\rightarrow$ pouvoir partagé.
\item \textbf{Voie de la violence :} l'employeur fait appel à des milices pour casser la grève ; ou les ouvriers sabotent les machines. La contrainte l'emporte sur la parole $\rightarrow$ violence. La relation est brisée, la rancœur s'installe.
\end{itemize}
\end{exemplebox}

\courssub{La spirale de la violence : l'engrenage destructeur}

\begin{propriete}[La violence engendre la violence]
La violence provoque \cle{souffrance} et \cle{humiliation}. La victime (ou son groupe) ressent le besoin de \cle{vengeance} ou de \cle{réparation} par la force. La riposte est souvent \emph{disproportionnée} $\rightarrow$ \cle{violence aggravée}. C'est un cercle vicieux auto-entretenu.
\end{propriete}

\begin{popfigure}
\begin{tikzpicture}[
    scale=0.9,
    every node/.style={font=\small},
    arrow/.style={->, >=Stealth, thick, popdark},
    stage/.style={draw=popdark, thick, rounded corners=6pt, fill=popcream, minimum width=3.2cm, minimum height=1.6cm, align=center},
    plus/.style={font=\Large, poporange, rotate=45}
]
% Nœuds en spirale
\node[stage, align=center] (v1) at (0,0){Violence\\initiale};
\node[stage, align=center] (s1) at (4.5,0){Souffrance\\humiliation};
\node[stage, align=center] (v2) at (4.5,-3.5){Vengeance\\représailles};
\node[stage, align=center] (s2) at (0,-3.5){Violence\\aggravée};
\node[stage, align=center] (v3) at (0,-7){Conflit\\généralisé};

% Flèches principales
\draw[arrow] (v1.east) -- node[above, popblue] {entraîne} (s1.west);
\draw[arrow] (s1.south) -- node[right, popblue] {provoque} (v2.north);
\draw[arrow] (v2.west) -- node[below, popblue] {devient} (s2.east);
\draw[arrow] (s2.south) -- node[left, popblue] {génère} (v3.north);

% Flèche de retour (bouclage)
\draw[arrow, poporange, dashed, bend left=20] (v3.west) to node[left, poporange] {reboucle} ++(-2,0) |- (v1.west);

% Indicateurs d'escalade
\node[plus] at (2.25, -0.4) {+};
\node[plus] at (4.9, -1.75) {+};
\node[plus] at (2.25, -3.9) {+};
\node[plus] at (-0.4, -5.25) {+};

% Légende interne
\node[align=center, text width=5cm, popmuted, font=\footnotesize] at (7.5, -3.5) {
    \textbf{Escalade :} chaque tour de spirale\\
    augmente l'intensité et\\
    élargit le cercle des victimes.
};
\end{tikzpicture}
\legende{Schéma de la spirale de la violence : violence $\to$ souffrance $\to$ vengeance $\to$ violence aggravée (escalade sans fin).}
\end{popfigure}

\begin{exemplebox}[Exemple camerounais : conflit entre éleveurs et agriculteurs]
Dans l'Extrême-Nord (ex. : département du Logone-et-Chari), un troupeau détruit un champ $\rightarrow$ \textbf{violence initiale} (dommage matériel). L'agriculteur blesse un éleveur $\rightarrow$ \textbf{souffrance/vengeance}. Le clan de l'éleveur incendie le village $\rightarrow$ \textbf{violence aggravée}. Les forces de sécurité interviennent, des innocents sont touchés $\rightarrow$ \textbf{conflit généralisé}. La spirale ne s'arrête que par la \textbf{médiation des chefs traditionnels} (lamibé, comités de sages) qui rétablissent la \emph{parole}.
\end{exemplebox}

\begin{exoresolu}[Analyse d'un cycle de violence scolaire]
\textbf{Énoncé :} dans un lycée, un élève (A) humilie publiquement un camarade (B) sur les réseaux sociaux. B répond en cassant le téléphone de A. A fait appel à son grand frère qui tabasse B. L'administration exclut les deux. Les parents portent plainte. La tension monte entre deux quartiers. Identifiez les étapes de la spirale et proposez une issue non-violente.
\tcblower
\textbf{Solution détaillée :}
\begin{enumerate}
\item \textbf{Violence symbolique initiale :} cyberharcèlement (atteinte à la dignité, en public).
\item \textbf{Souffrance/humiliation :} B se sent piégé, sans recours par la parole.
\item \textbf{Vengeance (violence physique) :} B casse le téléphone $\rightarrow$ riposte disproportionnée (bien matériel contre honneur).
\item \textbf{Violence aggravée (par intermédiaire) :} le grand frère de A tabasse B $\rightarrow$ escalade, sortie du cadre scolaire.
\item \textbf{Institutionnalisation du conflit :} exclusion, plainte en justice $\rightarrow$ la violence a détruit le lien éducatif.
\item \textbf{Issue non-violente possible :} médiation par les pairs, conseil de discipline à visée réparatrice, excuses publiques, réparation symbolique (service à la communauté).
\end{enumerate}
\end{exoresolu}

\courssub{La voie de la non-violence : Gandhi et la force de l'âme}

\begin{definition}[Non-violence (\emph{ahimsa})]
La \cle{non-violence} n'est \textbf{pas} la passivité, la lâcheté ou la soumission. C'est un \textbf{refus actif de nuire} associé à une \textbf{résistance déterminée à l'injustice} par des moyens pacifiques (désobéissance civile, grève, jeûne, marche, boycott). Gandhi l'appelle \emph{Satyagraha} : « force de la vérité » ou « force de l'âme ».
\end{definition}

\begin{aretenir}[Gandhi (1869--1948) : la non-violence comme arme des forts]
\begin{itemize}
\item « La non-violence est l'arme des forts. » Elle exige plus de courage que la violence : on accepte de souffrir sans infliger de souffrance.
\item Elle \textbf{désarme l'adversaire moralement} : le bourreau face à une victime qui ne répond pas par la force perd sa justification.
\item Elle \textbf{préserve l'humanité des deux camps} : l'opprimé ne devient pas oppresseur ; l'oppresseur peut retrouver son humanité.
\item \emph{Exemple historique :} la marche du sel (1930) et l'indépendance de l'Inde (1947), obtenue sans guerre civile généralisée contre le colonisateur.
\end{itemize}
\end{aretenir}

\begin{attention}[Erreur fréquente : non-violence $\neq$ résignation]
\begin{itemize}
\item \textbf{Résignation :} subir l'injustice sans réagir, par peur ou fatalisme $\rightarrow$ \emph{passivité}.
\item \textbf{Non-violence :} \emph{agir} puissamment pour transformer la situation, sans détruire l'adversaire $\rightarrow$ \emph{action politique radicale}.
\item \textbf{Martin Luther King} (inspiré de Gandhi) : « La non-violence est une arme puissante et juste, qui coupe sans blesser et ennoblit celui qui la manie. »
\end{itemize}
\end{attention}

\begin{exemplebox}[Exemple camerounais : les « villes mortes » et la désobéissance civile]
Dans les régions du Nord-Ouest et du Sud-Ouest (crise anglophone, depuis 2016), les populations ont observé des \emph{villes mortes} (fermeture des marchés, des écoles, des transports) sur appel de la société civile. C'est une \textbf{action non-violente} : refus de coopérer avec un ordre jugé injuste, sans armes. Elle vise à \textbf{contraindre par la pression économique et morale}, non par la destruction. Son efficacité politique est débattue, mais sa nature \emph{non-violente} est claire.
\end{exemplebox}

\courssub{Distinction capitale chez Arendt : violence et pouvoir}

\begin{propriete}[Violence vs Pouvoir (Hannah Arendt)]
\begin{center}
\begin{tabularx}{\linewidth}{|l|X|X|}\hline
\rowcolor{popblueL} \textbf{Critère} & \textbf{Pouvoir (\emph{Power})} & \textbf{Violence (\emph{Violence})} \\ \hline
\cle{Fondement} & Consentement libre, accord, légitimité & Contrainte, instruments, force brute \\ \hline
\cle{Moyen} & Parole, persuasion, organisation collective & Armes, argent, technologie, bureaucratie \\ \hline
\cle{Effet} & Crée de la stabilité, des institutions durables & Détruit le pouvoir, instaure la peur, instable \\ \hline
\cle{Portée} & Illimitée (plus on est d'accord, plus il grandit) & Limitée par les instruments disponibles \\ \hline
\cle{Rapport au nombre} & Nécessite le nombre (action concertée) & Peut être exercée par un seul (avec des armes) \\ \hline
\end{tabularx}
\end{center}
\end{propriete}

\begin{popfigure}
\begin{tikzpicture}[
    scale=0.85,
    every node/.style={font=\small},
    box/.style={draw=popdark, thick, rounded corners=4pt, minimum height=1.2cm, align=center},
    pbox/.style={box, fill=popgreenL, minimum width=4.5cm},
    vbox/.style={box, fill=poporangeL, minimum width=4.5cm},
    arrow/.style={->, >=Stealth, thick, popdark},
    brace/.style={decorate, decoration={brace, amplitude=6pt}, thick, popdark}
]
% Colonnes
\node[pbox, align=center] (p1) at (0,0){Fondement : \textbf{Consentement}\\Accord libre des agents};
\node[pbox, align=center] (p2) at (0,-1.5){Moyen : \textbf{Parole, persuasion}\\Action concertée};
\node[pbox, align=center] (p3) at (0,-3.0){Effet : \textbf{Création}\\Institutions, lois, confiance};
\node[pbox, align=center] (p4) at (0,-4.5){Portée : \textbf{Illimitée}\\Grandit avec le nombre};

\node[vbox, align=center] (v1) at (6,0){Fondement : \textbf{Contrainte}\\Instruments, force};
\node[vbox, align=center] (v2) at (6,-1.5){Moyen : \textbf{Armes, argent, tech.}\\Multiplication de la force};
\node[vbox, align=center] (v3) at (6,-3.0){Effet : \textbf{Destruction}\\Peur, instabilité, ruines};
\node[vbox, align=center] (v4) at (6,-4.5){Portée : \textbf{Limitée}\\Dépend des stocks d'armes};

% Titres
\node[font=\bfseries\large, popgreen] at (0,1.0) {POUVOIR};
\node[font=\bfseries\large, poporange] at (6,1.0) {VIOLENCE};

% Flèches d'opposition
\draw[arrow, dashed, poppink] (p1.east) -- node[above, poppink] {\textbf{OPPOSÉ}} (v1.west);
\draw[arrow, dashed, poppink] (p2.east) -- (v2.west);
\draw[arrow, dashed, poppink] (p3.east) -- (v3.west);
\draw[arrow, dashed, poppink] (p4.east) -- (v4.west);

% Accolades
\draw[brace] (-2.8, 0.6) -- (-2.8, -5.1) node[midway, left=8pt, rotate=90, popgreen] {\textbf{Processus générateur}};
\draw[brace] (8.8, 0.6) -- (8.8, -5.1) node[midway, right=8pt, rotate=-90, poporange] {\textbf{Processus destructeur}};
\end{tikzpicture}
\legende{Tableau de distinction : Pouvoir vs Violence selon Hannah Arendt. Le pouvoir naît du consentement et crée ; la violence repose sur les instruments et détruit.}
\end{popfigure}

\begin{methode}[Comment distinguer pouvoir et violence dans une situation concrète]
\begin{enumerate}
\item \textbf{Identifiez la source de l'obéissance :} les gens obéissent-ils parce qu'ils sont \emph{convaincus} (pouvoir) ou parce qu'ils \emph{craignent} une sanction (violence) ?
\item \textbf{Observez les moyens :} y a-t-il \emph{débat, vote, négociation} (parole) ou \emph{armes, matraques, coupures de salaire, censure} (instruments) ?
\item \textbf{Évaluez l'effet durable :} la relation est-elle \emph{stabilisée, institutionnalisée} (pouvoir) ou \emph{tendue, surveillée, prête à exploser} (violence) ?
\item \textbf{Test du nombre :} si le « chef » perd ses instruments, l'obéissance persiste-t-elle ? Si oui $\rightarrow$ pouvoir. Si non $\rightarrow$ violence.
\end{enumerate}
\end{methode}

\courssub{Application : le règlement pacifique des conflits au Cameroun}

\begin{exemplebox}[Les médiations traditionnelles et les comités de sages]
Au Cameroun, la \cle{résolution pacifique des conflits} repose sur un double héritage :
\begin{itemize}
\item \textbf{Traditionnel :} les \emph{chefs traditionnels}, les \emph{lamibé} (Nord), les \emph{fons} (Grassfields), les \emph{comités de sages} (villages, quartiers) pratiquent la \emph{médiation}, l'\emph{arbitrage}, la \emph{réparation symbolique} (excuses publiques, cadeaux, partage de la noix de cola, libations). \textbf{Principe :} restaurer l'harmonie sociale plutôt que punir.
\item \textbf{Moderne :} la \emph{Commission Nationale des Droits de l'Homme et des Libertés (CNDHL)}, les \emph{maisons de la justice et du droit}, la \emph{justice de proximité}, les juridictions encadrées par la loi.
\end{itemize}
\end{exemplebox}

\begin{exoresolu}[Cas pratique : conflit intercommunautaire à Kousseri (Extrême-Nord, 2021)]
\textbf{Énoncé :} en décembre 2021, des affrontements entre pêcheurs mousgoums et éleveurs arabes choas pour l'accès à l'eau et aux pâturages du Logone ont fait des dizaines de morts, des villages brûlés et des milliers de déplacés. Analysez ce conflit avec les concepts du cours et dégagez les voies d'une sortie non-violente.
\tcblower
\textbf{Analyse par les concepts du cours :}
\begin{enumerate}
\item \textbf{Violence structurelle :} raréfaction de la ressource (changement climatique, retrait du lac Tchad) + absence de cadre concerté de gestion $\rightarrow$ tension latente.
\item \textbf{Déclencheur :} incident mineur (contestation d'un canal de pêche) $\rightarrow$ \textbf{violence initiale} (rixe, blessure).
\item \textbf{Spirale :} représailles intercommunautaires $\rightarrow$ \textbf{violence aggravée} (tueries, incendies, exil vers le Tchad voisin).
\item \textbf{Échec de la parole :} pas d'espace de dialogue préventif, méfiance, rumeurs.
\end{enumerate}
\textbf{Issue non-violente :}
\begin{itemize}
\item \textbf{Médiation de l'État :} autorités administratives (gouverneur, préfet) et forces de sécurité pour faire cesser les affrontements.
\item \textbf{Médiation traditionnelle :} sultans, lamidos, chefs mousgoums $\rightarrow$ \textbf{parole}, serments, délimitation de couloirs de transhumance, calendrier d'accès à l'eau.
\item \textbf{Justice réparatrice :} indemnisation des familles, reconstruction des cases, cérémonies de réconciliation (partage de repas).
\item \textbf{Institutionnalisation :} création de \textbf{comités mixtes de gestion des ressources} (pouvoir partagé, consentement).
\end{itemize}
\textbf{Leçon :} la violence a détruit ; seule la \textbf{parole institutionnalisée} (pouvoir) peut reconstruire durablement.
\end{exoresolu}

\begin{aretenir}[Bilan du procès : pourquoi condamner la violence ?]
\begin{enumerate}
\item \textbf{Moralement :} elle nie la \cle{dignité} (Kant) — l'homme devient chose.
\item \textbf{Politiquement :} elle détruit le \cle{pouvoir} (Arendt) — la contrainte remplace le consentement.
\item \textbf{Dynamiquement :} elle enclenche une \cle{spirale} sans fin — vengeance, escalade, guerre de tous contre tous.
\item \textbf{Éthiquement :} la \cle{non-violence} (Gandhi) offre une alternative \emph{active}, \emph{courageuse}, \emph{constructrice} qui préserve l'humanité des deux camps.
\item \textbf{Concrètement (Cameroun) :} la \emph{médiation}, le \emph{dialogue}, la \emph{justice réparatrice} (traditionnelle et moderne) sont les seules voies durables pour « vivre ensemble ».
\end{enumerate}
\end{aretenir}

\begin{savaistu}[La « palabre » africaine : une philosophie de la non-violence ?]
La \emph{palabre} (du portugais \emph{palavra}, « parole ») est l'assemblée où l'on « parle jusqu'à l'accord ». Chez les Bamiléké, les Bassa, les Douala, les Peuls, etc., le conflit se règle \emph{par la parole}, en public, avec des tiers (sages, notables). On ne vote pas : on \emph{cherche le consensus}. C'est l'exercice concret du \textbf{pouvoir arendtien} : le consentement libre par la parole. Le recours à la violence (coup, arme) est l'aveu d'échec de la palabre. Aujourd'hui, les dispositifs de médiation au Cameroun tentent d'inscrire cette sagesse dans le droit positif.
\end{savaistu}

\coursec{Le procès de la violence (II) : peut-on la justifier ? Bilan}

Dans la section précédente, nous avons instruit le procès de la violence et entendu l'accusation : la violence détruit la parole, nie la dignité de la personne humaine, engendre la spirale de la vengeance et s'avère souvent impuissante à fonder un ordre durable. Le verdict semblait sans appel. Pourtant, le procès n'est pas clos. Il reste à entendre la défense. Car la vie réelle oppose à notre condamnation de principe des situations troublantes : la victime qui repousse son agresseur, le policier qui emploie la force pour protéger les citoyens, le peuple qui se soulève contre un tyran. Toute violence est-elle également condamnable ? Ou bien existe-t-il des cas où la violence, sans devenir bonne, pourrait être \emph{justifiée} ? C'est à cette question redoutable que cette dernière section est consacrée, avant de tirer le bilan de toute la leçon.

\courssub{L'hypothèse de la légitime défense}

Commençons par l'intuition la plus simple. Un soir, à Douala, un passant est agressé par un individu armé qui veut lui arracher son sac. Le passant se défend, frappe l'agresseur et parvient à fuir. Qui oserait condamner ce passant au même titre que l'agresseur ? Personne. Notre sens moral spontané distingue ici la violence qui \emph{attaque} de la violence qui \emph{repousse}. C'est cette distinction que le droit consacre sous le nom de \cle{légitime défense}.

\begin{definition}[La légitime défense]
La \cle{légitime défense} est la riposte \textbf{nécessaire}, \textbf{immédiate} et \textbf{proportionnée} à une agression \textbf{injuste} et \textbf{en cours}. Le droit pénal camerounais (Code pénal, art. 84) reconnaît que n'est pas punissable l'auteur d'un acte de défense lorsque celui-ci répond à une agression injustifiée envers lui-même ou autrui, à condition que la riposte soit simultanée à l'attaque et proportionnée à sa gravité.
\end{definition}

Analysons rigoureusement les trois conditions, car c'est elles qui séparent la légitime défense de la vengeance déguisée :

\begin{enumerate}
\item \textbf{Une agression injuste et actuelle} : il faut qu'une attaque soit réellement en train de se produire. On ne peut pas frapper « par précaution » quelqu'un qui n'a rien fait, ni se venger le lendemain d'une agression passée : la riposte tardive redevient une violence privée, c'est-à-dire un crime.
\item \textbf{La nécessité} : la force ne doit être employée que s'il n'existe pas d'autre moyen d'échapper à l'agression (fuite, appel au secours). La légitime défense est un \emph{dernier recours}, jamais un choix de commodité.
\item \textbf{La proportionnalité} : la riposte ne doit pas excéder ce qu'exige la protection. Répondre à une gifle par un coup de couteau mortel, c'est franchir la ligne : la défense devient à son tour agression.
\end{enumerate}

L'enjeu philosophique est profond : la légitime défense ne dit pas que la violence est bonne ; elle dit que, face à une violence injuste qui nie le droit, le droit lui-même autorise la victime à survivre. Comme le notait déjà la tradition du droit naturel, le premier devoir d'un être vivant est de se conserver : nul ne peut être tenu de se laisser détruire. La légitime défense est donc moins une \emph{justification} de la violence qu'une \emph{excuse} : elle rend la riposte non punissable, sans la transformer en bien moral.

\courssub{La violence de l'État : force publique ou violence ?}

Passons à un second cas, plus délicat. Chaque jour, au Cameroun, la police et la gendarmerie emploient la force : elles arrêtent des malfaiteurs, dispersent des attroupements dangereux, parfois tirent. Est-ce de la violence ? Et si oui, est-elle justifiable ?

Rappelons la définition célèbre du sociologue Max \textbf{Weber} (\emph{Le savant et le politique}, 1919) : l'État est la communauté humaine qui revendique, avec succès, le \emph{monopole de la violence physique légitime} sur un territoire donné. Cette formule contient toute la difficulté. Weber ne dit pas que l'État supprime la violence : il dit que l'État la \emph{concentre} entre ses mains et prétend être le seul à pouvoir l'employer légitimement. La force publique est donc bien une violence — elle contraint des corps — mais une violence qui se veut \emph{ordonnée au droit}.

Sous quelles conditions cette prétention est-elle recevable ? Trois critères permettent de distinguer la \cle{force légitime} de la violence étatique arbitraire :

\begin{itemize}
\item \textbf{La légalité} : la force publique ne peut agir que dans le cadre des lois. Un policier qui frappe un citoyen paisible, ou qui punit lui-même au lieu de déférer le suspect à la justice, sort du droit et retombe dans la violence criminelle.
\item \textbf{La proportionnalité} : la contrainte employée doit être mesurée à la menace. Tirer sur un voleur de mangues en fuite serait une disproportion monstrueuse.
\item \textbf{La finalité} : la force publique doit viser la protection des personnes et le maintien de l'ordre juridique, non l'intérêt privé d'un gouvernant ou la terreur. Dès qu'elle sert à opprimer plutôt qu'à protéger, elle perd sa légitimité.
\end{itemize}

On comprend alors la position de Hobbes : si l'État existe, c'est précisément pour \emph{retirer} la violence des mains des particuliers et mettre fin à la « guerre de tous contre tous ». La force publique est le prix — toujours à surveiller — de la paix civile. Mais le philosophe ajoute aussitôt une vigilance critique : le monopole de la violence n'est légitime que s'il sert l'ordre commun. Un État qui torture, qui emprisonne sans jugement, qui tire sur ses propres citoyens désarmés, trahit sa mission et rejoint, par un paradoxe tragique, le camp de la violence qu'il était chargé de contenir.

\courssub{La violence révolutionnaire : le droit de résistance à l'oppression}

Troisième et dernier cas, le plus explosif : si l'État lui-même devient tyrannique, le peuple peut-il légitimement renverser le pouvoir par la force ?

L'argument en faveur du \cle{droit de résistance à l'oppression} est puissant. Si la légitimité de la force publique repose sur sa finalité — protéger les droits —, alors un pouvoir qui viole systématiquement ces droits détruit lui-même le fondement de son autorité. Locke, dans le \emph{Traité du gouvernement civil} (1690), l'affirme clairement : lorsque le gouvernement trahit la confiance du peuple et se livre à la force sans droit, il entre en « état de guerre » contre ses citoyens, et ceux-ci retrouvent leur liberté naturelle de lui résister. La Révolution française de 1789, l'abolition des régimes coloniaux, les luttes contre les dictatures du XX\textsuperscript{e} siècle s'inscrivent dans cette lignée.

\begin{savaistu}[Un droit naturel inscrit dans un texte fondateur]
La Déclaration des droits de l'homme et du citoyen de 1789 reconnaît la \textbf{résistance à l'oppression} comme l'un des « droits naturels et imprescriptibles de l'homme » (article 2), au même titre que la liberté, la propriété et la sûreté. Ce texte, né de la Révolution française, a inspiré la plupart des constitutions modernes, y compris africaines : il affirme qu'un peuple n'est jamais tenu d'obéir à un pouvoir qui le tyrannise.
\end{savaistu}

Mais la justification révolutionnaire se heurte à de redoutables objections, que Hannah \textbf{Arendt} (\emph{Du mensonge à la violence}, 1970) formule avec lucidité. D'abord, la violence ne crée pas le pouvoir, elle le détruit : un régime renversé par les armes laisse un vide que la violence seule ne sait pas remplir — d'où la tentation, pour les vainqueurs, de gouverner par la terreur à leur tour. L'histoire le montre cruellement : bien des révolutions ont débouché sur de nouvelles dictatures, parfois pires que l'ancienne. Ensuite, qui décide qu'un pouvoir est « tyrannique » ? Chaque groupe mécontent peut s'autoproclamer opprimé et légitimer sa révolte : le droit de résistance risque alors de devenir le droit du plus fort ou du plus violent. Enfin, la violence révolutionnaire exige presque toujours des victimes innocentes : peut-on sacrifier des vies présentes au nom d'un avenir hypothétique ?

\courssub{Critique de ces justifications : les trois pièges}

Face à ces trois candidats à la justification — légitime défense, force publique, résistance à l'oppression —, la philosophie doit exercer sa vigilance. Trois dangers guettent toute justification de la violence.

\begin{enumerate}
\item \textbf{Le risque de l'arbitraire} : chacun peut se déclarer « en état de légitime défense » ou « opprimé ». Sans juges impartiaux et sans critères stricts, la justification devient un blanc-seing que n'importe quel agresseur peut invoquer. C'est pourquoi le droit encadre si étroitement ces notions.
\item \textbf{Le risque de l'escalade} : la violence justifiée appelle une contre-violence qui se prétendra à son tour justifiée. La riposte proportionnée d'aujourd'hui devient l'agression à venger demain : c'est la mécanique tragique des vendettas et des guerres civiles.
\item \textbf{Le risque de la banalisation} : à force d'admettre des exceptions, on finit par habiter la violence, par la trouver normale, voire nécessaire en toute circonstance. Or, rappelle Arendt, la violence reste \emph{muette} : elle ne prouve rien, ne convainc personne, elle impose. La banaliser, c'est renoncer à la parole, seul instrument proprement humain du règlement des conflits.
\end{enumerate}

\begin{attention}[Erreur fréquente en dissertation]
Conclure par « la violence, c'est mal » sans nuancer ni justifier philosophiquement est une double faute. D'une part, c'est une \emph{petitio principii} : on affirme ce qu'il fallait démontrer. D'autre part, c'est ignorer le problème réel : le candidat sérieux doit montrer \emph{pourquoi} la violence est condamnable (destruction de la parole, négation de la dignité) \emph{et} affronter honnêtement les cas limites (légitime défense, tyrannie) au lieu de les balayer d'un slogan moral.
\end{attention}

\begin{popfigure}
\begin{center}
\begin{tikzpicture}[scale=1]
% Poutre centrale
\draw[line width=2pt, popdark] (-3.2,0) -- (3.2,0);
\draw[line width=2pt, popdark] (0,0) -- (0,-0.6);
\draw[popdark, fill=popdark] (-0.5,-1.6) -- (0.5,-1.6) -- (0,-0.6) -- cycle;
% Plateau gauche (pour)
\draw[line width=1.2pt, popgreen] (-3.2,0) -- (-3.9,-1.1);
\draw[line width=1.2pt, popgreen] (-3.2,0) -- (-2.5,-1.1);
\draw[line width=1.5pt, popgreen, fill=popgreenL] (-4.3,-1.1) arc (180:360:1.1);
\node[align=center, font=\small\bfseries, text=popgreen!60!black] at (-3.2,-2.5) {ARGUMENTS POUR\\ \footnotesize légitime défense\\ \footnotesize force publique de l'État\\ \footnotesize résistance à la tyrannie};
% Plateau droit (contre)
\draw[line width=1.2pt, poporange] (3.2,0) -- (2.5,-1.1);
\draw[line width=1.2pt, poporange] (3.2,0) -- (3.9,-1.1);
\draw[line width=1.5pt, poporange, fill=poporangeL] (2.1,-1.1) arc (180:360:1.1);
\node[align=center, font=\small\bfseries, text=poporange!70!black] at (3.2,-2.5) {ARGUMENTS CONTRE\\ \footnotesize dignité de la personne\\ \footnotesize spirale de l'escalade\\ \footnotesize échec de la parole};
% Aiguille
\draw[line width=1.5pt, popink] (0,0) -- (0,0.9);
\fill[popink] (0,0.9) circle (2pt);
\node[font=\small\itshape, text=popink] at (0,1.25) {la balance hésite};
\end{tikzpicture}
\end{center}
\legende{La balance du procès : les arguments en faveur d'une violence justifiée (légitime défense, force publique, résistance à l'oppression) pèsent contre les arguments de la condamnation (dignité, escalade, échec de la parole). Le déséquilibre final dépend des conditions : nécessité, proportionnalité, dernier recours.}
\end{popfigure}

\courssub{Bilan dialectique : le problème tragique du moindre mal}

Comment trancher ? La démarche dialectique consiste à dépasser l'opposition simple entre la thèse (toute violence est condamnable) et l'antithèse (certaines violences sont justifiables) vers une position synthétique et nuancée.

La synthèse s'énonce ainsi : \textbf{la violence est en principe condamnable, parce qu'elle détruit la parole et nie la dignité de la personne ; mais certaines situations extrêmes posent le problème tragique du moindre mal}. Face à une agression meurtrière, face à une tyrannie sanglante, le refus absolu de toute force peut revenir à livrer les innocents à leurs bourreaux. Dans ces cas limites, la violence défensive n'est jamais un \emph{bien} — elle reste un mal — mais elle peut être le \emph{moindre mal}, c'est-à-dire le seul moyen d'empêcher un mal plus grand. Le tragique tient précisément à ceci : quelle que soit la décision, on ne sort pas indemne ; même justifiée, la violence laisse une blessure, une culpabilité, un deuil.

C'est pourquoi la justification ne doit jamais être automatique : elle doit passer par un examen rigoureux, cas par cas, que l'on peut figurer sous la forme d'un arbre de décision.

\begin{popfigure}
\begin{center}
\begin{tikzpicture}[
  every node/.style={font=\small},
  quest/.style={draw=popblue, fill=popblueL, rounded corners=3pt, align=center, inner sep=4pt, text width=3.4cm},
  rep/.style={font=\small\itshape, text=popmuted},
  concl/.style={draw, rounded corners=3pt, align=center, inner sep=4pt, text width=3.2cm, font=\small\bfseries},
  oui/.style={concl, draw=popgreen, fill=popgreenL, text=popgreen!50!black},
  non/.style={concl, draw=poporange, fill=poporangeL, text=poporange!70!black},
  arr/.style={-{Stealth[length=2.5mm]}, thick, popdark}
]
\node[quest] (q1) at (0,0) {Y a-t-il une agression \textbf{injuste} et \textbf{immédiate} ?};
\node[quest] (q2) at (4.6,0) {Tous les moyens \textbf{non violents} sont-ils épuisés ? (dernier recours)};
\node[quest] (q3) at (9.2,0) {La riposte est-elle \textbf{proportionnée} à la menace ?};
\node[oui] (o1) at (9.2,-2.6) {Violence \textbf{justifiable} (excusable), jamais « bonne »};
\node[non] (n1) at (0,-2.6) {Violence \textbf{injustifiable} : c'est une agression};
\node[non] (n2) at (4.6,-2.6) {Violence \textbf{injustifiable} : préférer la parole, la fuite, le droit};
\node[non] (n3) at (12.6,-2.6) {Violence \textbf{injustifiable} : la riposte devient agression};
\draw[arr] (q1) -- node[rep, above]{oui} (q2);
\draw[arr] (q2) -- node[rep, above]{oui} (q3);
\draw[arr] (q3) -- node[rep, right]{oui} (o1);
\draw[arr] (q1) -- node[rep, right]{non} (n1);
\draw[arr] (q2) -- node[rep, right]{non} (n2);
\draw[arr] (q3.east) -- node[rep, right]{non} (n3);
\end{tikzpicture}
\end{center}
\legende{Arbre de décision : « une violence est-elle justifiable ? ». Une riposte n'est excusable que si elle passe les trois tests : agression injuste et immédiate, épuisement des voies non violentes (dernier recours), proportionnalité. Un seul « non » suffit à la condamner.}
\end{popfigure}

\begin{exoresolu}[Appliquer l'arbre de décision à une situation concrète]
\textbf{Énoncé.} À Bamenda, un commerçant surpris par un cambrioleur désarmé qui s'enfuit avec sa caisse lui tire dans le dos et le blesse grièvement. Le commerçant invoque la légitime défense. En appliquant les trois critères de l'arbre de décision, dites si cette violence est justifiable.
\tcblower
\textbf{Solution détaillée.}
\begin{enumerate}
\item \emph{Agression injuste et immédiate ?} Le vol est bien une agression injuste, mais au moment du tir, le cambrioleur \textbf{fuyait} : le danger immédiat pour la vie du commerçant n'existait plus. Le premier critère n'est que partiellement rempli.
\item \emph{Dernier recours ?} D'autres moyens existaient : crier, alerter le voisinage, appeler la police, qui pouvait poursuivre le fuyard. La force n'était pas nécessaire.
\item \emph{Proportionnalité ?} Tirer une balle sur un homme désarmé pour protéger de l'argent est une disproportion flagrante : on échange un bien matériel contre une vie humaine, ce qui renverse la hiérarchie des valeurs.
\end{enumerate}
\textbf{Conclusion} : deux « non » sur trois. La violence du commerçant n'est pas une légitime défense mais une agression, punissable par le droit pénal. Ce cas montre l'utilité pratique de la philosophie : elle fournit des critères pour juger, là où l'émotion spontanée (« il n'avait qu'à pas voler ») condamnerait à tort.
\end{exoresolu}

\courssub{Conclusion de la leçon}

Nous pouvons maintenant répondre à la question qui a guidé toute la leçon : la violence peut-elle être justifiée ? L'analyse nous a conduits à une réponse en deux temps. \emph{En principe}, la violence est condamnable : elle est l'échec de la parole, la négation de l'autre comme personne, et elle engendre mécaniquement la spirale de l'escalade. \emph{Dans des situations extrêmes} — agression immédiate, tyrannie sans recours —, une force strictement nécessaire, proportionnée et ordonnée à la protection du droit peut être excusée comme le moindre mal, sans jamais devenir un bien. La philosophie n'absout donc pas la violence : elle nous apprend à la \emph{juger}, c'est-à-dire à refuser à la fois la naïveté qui croit pouvoir l'ignorer et le cynisme qui prétend tout justifier.

L'enjeu de cette leçon dépasse le seul problème de la violence : il touche à la question de savoir comment des hommes peuvent vivre ensemble sans se détruire. Si la violence est l'échec de la parole, alors la paix exige des institutions qui rendent la parole efficace : des lois, des tribunaux, un État de droit. Mais qu'est-ce qu'une loi juste ? Et que vaut un État qui emploie la force ? C'est à ces questions que nous conduira la prochaine leçon, consacrée à la \cle{justice} et à l'\cle{État}.

\begin{methode}[Rédiger une conclusion de dissertation]
Une conclusion réussie se construit en trois gestes :
\begin{enumerate}
\item \textbf{Répondre nettement à la question posée} : reprenez le sujet et donnez votre réponse nuancée en une ou deux phrases (ici : « la violence est en principe condamnable, mais certaines situations extrêmes posent le problème du moindre mal »). Interdiction de rester vague ou de découvrir une idée nouvelle.
\item \textbf{Rappeler l'enjeu} : montrez pourquoi la question méritait d'être posée (ici : vivre ensemble sans se détruire, fonder la paix).
\item \textbf{Ouvrir sur une question nouvelle} : prolongez par un problème voisin qui découle logiquement de votre réponse (ici : la justice, l'État de droit), sans le traiter.
\end{enumerate}
\end{methode}

\begin{exemplebox}[Ce qui fait la différence entre deux copies]
En dissertation, une copie gagnante mobilise \textbf{au moins 3 références précises} — par exemple Hobbes (l'état de nature comme guerre de tous contre tous), Weber (le monopole de la violence légitime) et Arendt (la violence comme échec de la parole) — et \textbf{au moins 2 exemples concrets} (une agression dans la rue à Douala jugée selon les critères de la légitime défense ; la Révolution française et l'article 2 de la Déclaration de 1789). Une copie qui ne cite personne et ne donne aucun exemple reste une opinion, pas une démonstration philosophique.
\end{exemplebox}

\begin{aretenir}[L'essentiel de la section]
\begin{itemize}
\item La \cle{légitime défense} (Code pénal camerounais, art. 84) : riposte nécessaire, immédiate et proportionnée à une agression injuste en cours.
\item La force publique de l'État n'est légitime que sous trois conditions : \textbf{légalité}, \textbf{proportionnalité}, \textbf{finalité} (protection de l'ordre et des personnes).
\item Le \cle{droit de résistance à l'oppression} (Locke ; Déclaration de 1789, art. 2) permet de renverser une tyrannie, mais la violence révolutionnaire risque l'arbitraire, l'escalade et la banalisation.
\item Bilan : la violence est en principe condamnable ; seules des situations extrêmes posent le problème tragique du \emph{moindre mal}. La philosophie n'absout pas la violence : elle apprend à la juger.
\end{itemize}
\end{aretenir}

\begin{exercice}[Pour s'entraîner]
Sujet de dissertation : « Peut-on parfois faire le bien par la violence ? » Construisez un plan dialectique en trois parties (thèse, antithèse, synthèse) en mobilisant au moins trois auteurs étudiés dans la leçon et deux exemples camerounais ou historiques. Rédigez ensuite l'introduction et la conclusion complètes.
\end{exercice}

\coursec{Atelier méthodologique : appliquer la leçon à des situations concrètes}

Après avoir mené le \cle{procès de la violence} — en examinant tour à tour ses condamnations morales et juridiques, puis les rares tentatives de justification (légitime défense, révolution, contrainte étatique) — nous devons maintenant vérifier que les concepts acquis ne restent pas des abstractions. La philosophie, au programme du MINESEC en classe de Première technique, vise précisément à \textbf{appliquer les notions à des situations concrètes de la vie courante} et à \textbf{rédiger et justifier une démarche, une réponse ou une conclusion}. C'est l'objet de cet atelier : transformer la culture philosophique en compétence argumentative, celle-là même que le Probatoire et le Baccalauréat technique évaluent.

\courssub{Méthode : traiter une étude de cas philosophique}

\begin{methode}[Démarche en quatre étapes pour analyser un cas concret]
\begin{enumerate}
\item \textbf{Décrire les faits sans les juger} : identifier les acteurs, l'action, le contexte, les conséquences visibles. Distinguer ce qui est \emph{donné} (observable) de ce qui est \emph{interprété}.
\item \textbf{Qualifier à l'aide des notions du cours} : quel(s) type(s) de violence ? (physique, psychologique, structurelle, symbolique, légitime/illégitime). Quelle origine invoquée ? (naturelle, culturelle, institutionnelle).
\item \textbf{Mobiliser une thèse d'auteur} : choisir le référent pertinent (Weber sur le monopole de la violence légitime, Kant sur la dignité, Arendt sur la distinction pouvoir/violence, Gandhi sur la non-violence, etc.) et expliciter \emph{pourquoi} il éclaire ce cas.
\item \textbf{Produire un jugement nuancé} : conclure en pesant les arguments, en distinguant le plan juridique, moral et politique, sans réduire la complexité à un « pour/contre » binaire.
\end{enumerate}
\end{methode}

\begin{attention}[Piège fréquent : raconter au lieu d'analyser]
L'erreur majeure, en dissertation comme en commentaire, est de \textbf{substituer le récit à l'analyse}. Dire « un élève se fait racketter, c'est triste, il a peur » n'est pas une réponse philosophique. La réponse attendue : « Cette situation relève de la \cle{violence physique} (contrainte corporelle) et \cle{psychologique} (intimidation), exercée par un groupe sur un individu isolé ; elle illustre la définition de la violence comme \emph{atteinte à l'intégrité et à la liberté d'autrui} et appelle, selon Gandhi, une résistance non-violente qui désarme l'agresseur sans le détruire. » \textbf{Chaque affirmation doit être étayée par une notion ou une référence du cours.}
\end{attention}

\courssub{Étude de cas 1 : un élève racketté à la sortie du lycée}

\begin{exemplebox}[Situation]
Chaque soir, à la sortie d'un lycée technique de Yaoundé, un groupe de trois élèves plus âgés attend un camarade de Première. Ils lui demandent « la taxe » : 500 FCFA pour « protéger son vélo ». S'il refuse, ils le bousculent et menacent de le « casser » le lendemain. L'élève paie, rentre chez lui en silence, n'ose pas en parler à ses parents ni au surveillant général.
\end{exemplebox}

\courssub{Analyse guidée}

\begin{enumerate}
\item \textbf{Identification des types de violence} :
\begin{itemize}
\item \cle{Violence physique} : bousculades, menaces de coups (« te casser ») — atteinte à l'intégrité corporelle.
\item \cle{Violence psychologique} (ou morale) : intimidation, peur quotidienne, silence forcé — atteinte à la liberté psychique.
\item \cle{Violence structurelle} : rapport de domination asymétrique (âge, nombre, force) qui s'institutionnalise en « taxe » quasi coutumière et prive la victime de la satisfaction d'un besoin fondamental, la sécurité.
\item \cle{Violence symbolique} : l'imposition d'un vocabulaire (« taxe », « protection ») qui maquille l'arbitraire en échange légitime et fait intérioriser la soumission.
\end{itemize}

\item \textbf{Réponses non-violentes : pistes concrètes} :
\begin{itemize}
\item \textbf{Rompre la loi du silence} : alerter un adulte de confiance (surveillant général, conseiller d'orientation, parent) — la violence prospère dans le secret.
\item \textbf{Collectiviser la riposte} : signaler collectivement (délégués de classe, association des élèves) pour briser l'isolement de la victime.
\item \textbf{Mobiliser le règlement intérieur} : la procédure disciplinaire (conseil de discipline) est une \cle{procédure légitime}, non violente, qui substitue la sanction rationnelle à la vengeance privée.
\item \textbf{Éducation par les pairs} : ateliers de médiation entre élèves, apprentissage de l'expression non-violente des besoins (« j'ai besoin de sécurité ») sans agresser autrui.
\end{itemize}
\end{enumerate}

\begin{exoresolu}[Application de la méthode]
\textbf{Énoncé} : En vous appuyant sur la distinction entre pouvoir et violence (Arendt) et sur le concept de violence structurelle, analysez la situation du racket à la sortie du lycée.
\tcblower
\textbf{Solution détaillée} :
\begin{enumerate}
\item \textbf{Faits} : un groupe dominant impose une redevance sous la menace à un élève isolé ; répétition quotidienne ; silence de la victime.
\item \textbf{Notions} :
\begin{itemize}
\item \emph{Violence} (Arendt) : elle est \emph{instrumentale} (elle a besoin d'outils, ici la force du nombre et la menace) et elle détruit le pouvoir — ici, le pouvoir de l'élève de disposer de lui-même et de ses biens.
\item \emph{Pouvoir} (Arendt) : capacité d'\emph{agir de concert} — absent chez la victime isolée, présent chez le groupe agresseur qui agit ensemble.
\item \emph{Violence structurelle} : des rapports sociaux inégalitaires (hiérarchie d'âge, impunité) empêchent la victime de satisfaire ses besoins fondamentaux (sécurité, scolarité sereine), sans qu'un agresseur unique soit toujours identifiable.
\end{itemize}
\item \textbf{Thèses} : Arendt montre que la violence et le pouvoir sont des contraires : « là où l'un règne absolument, l'autre est absent » — le racket règne précisément parce que le pouvoir collectif des élèves et de l'institution est défaillant. L'analyse de la violence structurelle invite à transformer les \emph{structures} (encadrement, règlement, médiation) et non seulement à punir les individus.
\item \textbf{Jugement} : la réponse non-violente la plus efficace n'est pas la passivité mais l'\emph{action collective institutionnalisée} (signalement, médiation, sanction éducative) qui restaure le \cle{pouvoir} d'agir ensemble là où la violence l'avait détruit.
\end{enumerate}
\end{exoresolu}

\courssub{Étude de cas 2 : manifestation réprimée par les forces de l'ordre}

\begin{exemplebox}[Situation]
Dans une grande ville du Cameroun, des élèves et étudiants manifestent pacifiquement contre la hausse des frais d'inscription. Ils occupent un rond-point, assis, en chantant l'hymne national. Les forces de maintien de l'ordre interviennent : gaz lacrymogènes, matraques. Bilan : plusieurs blessés et de nombreuses interpellations. Le communiqué officiel parle de « rétablissement de l'ordre public » et de « casseurs infiltrés ».
\end{exemplebox}

\courssub{Analyse guidée : force légitime ou violence illégitime ? (Weber)}

\begin{definition}[Max Weber, \emph{Le savant et le politique} (1919)]
L'\cle{État} se définit par le \textbf{monopole de la violence physique légitime} sur un territoire donné. La \cle{force légitime} est celle qui s'exerce \emph{dans le cadre de la loi}, pour faire respecter la loi, avec \emph{proportionnalité} et \emph{nécessité}. La \cle{violence illégitime} est tout usage de la contrainte physique \emph{hors cadre légal}, \emph{disproportionné} ou \emph{arbitraire}.
\end{definition}

\begin{enumerate}
\item \textbf{Qualification de l'action des manifestants} :
\begin{itemize}
\item Assis, chantant, non armés : \cle{action non-violente} (désobéissance civile, au sens de Gandhi et de Martin Luther King).
\item Occupation du rond-point : \cle{entrave à la circulation} — infraction éventuelle, mais non \emph{violente} au sens strict (pas d'atteinte à l'intégrité corporelle d'autrui).
\end{itemize}

\item \textbf{Qualification de la riposte des forces de l'ordre} :
\begin{itemize}
\item \emph{Légitime} SI : sommations préalables, proportionnalité (disperser, non blesser), ciblage des seuls fauteurs de troubles identifiés.
\item \emph{Illégitime} SI : usage de munitions dangereuses contre des personnes assises (disproportionné), matraquages au sol (traitement dégradant), arrestations arbitraires.
\end{itemize}

\item \textbf{Grille weberienne d'évaluation} :
\begin{center}
\begin{tabularx}{\linewidth}{|l|X|X|}\hline
\rowcolor{popblueL}\textbf{Critère weberien} & \textbf{Force légitime (idéal)} & \textbf{Violence illégitime (dérive possible)} \\ \hline
Monopole & État seul autorisé à contraindre & Groupes para-publics, bavures non sanctionnées \\ \hline
Légalité & Loi, procédures écrites & Ordres oraux flous, impunité \\ \hline
Proportionnalité & Moyens adaptés au but & Moyens létaux contre des manifestants assis \\ \hline
Nécessité & Dernier recours après négociation & Intervention immédiate sans sommation \\ \hline
Reddition de comptes & Enquête indépendante, sanctions & Classement sans suite, mutisme \\ \hline
\end{tabularx}
\end{center}
\end{enumerate}

\begin{aretenir}[Distinction cruciale pour l'examen]
\textbf{Force légitime} = contrainte \emph{encadrée par le droit}, au service d'un ordre \emph{juste}. \textbf{Violence illégitime} = contrainte \emph{arbitraire}, \emph{disproportionnée}, qui \emph{nie le droit} qu'elle prétend servir. \textbf{La légitimité ne va pas de soi : elle se prouve par la procédure, la proportionnalité et le contrôle.}
\end{aretenir}

\courssub{Étude de cas 3 : un mari qui « corrige » sa femme}

\begin{exemplebox}[Situation]
Dans un quartier d'une ville camerounaise, un homme frappe régulièrement son épouse pour « la corriger » : elle aurait « mal servi le repas », « répondu sans respect », « sorti sans permission ». Il déclare : « C'est notre culture, l'homme est le chef, il doit éduquer sa femme. » La famille de la femme conseille la patience : « Le mariage, c'est l'endurance. » Elle ne porte pas plainte.
\end{exemplebox}

\courssub{Analyse guidée : violence domestique et négation de la dignité (Kant)}

\begin{definition}[Immanuel Kant, \emph{Fondements de la métaphysique des mœurs} (1785)]
La \cle{dignité} est la valeur \emph{intrinsèque} et \emph{incomparable} de tout être raisonnable, qui ne doit jamais être traité \emph{uniquement comme un moyen}, mais toujours \emph{en même temps comme une fin}. \textbf{L'humanité, dans la personne de chacun, est une fin en soi.}
\end{definition}

\begin{enumerate}
\item \textbf{Types de violence présents} :
\begin{itemize}
\item \cle{Violence physique} : coups, blessures.
\item \cle{Violence psychologique} : terreur quotidienne, surveillance, isolement.
\item \cle{Violence symbolique} : langage qui naturalise la domination (« culture », « éducation », « chef »).
\item \cle{Violence structurelle} : dépendance économique, poids des normes sociales et de la famille élargie qui enferment la victime.
\end{itemize}

\item \textbf{Analyse kantienne} :
\begin{itemize}
\item Le mari traite sa femme comme un \textbf{moyen} (instrument de service, de soumission, de satisfaction de son autorité) et \textbf{nie sa qualité de fin en soi} (sa capacité à se fixer ses propres fins, à consentir ou à refuser).
\item La « correction » n'est pas une \emph{éducation} (qui vise l'autonomie) mais un \emph{dressage} (qui vise la soumission). Elle détruit la \cle{personnalité morale} : la capacité d'agir par devoir, et non par peur.
\item L'argument « culturel » est un \textbf{sophisme} : il confond la \emph{coutume} (ce qui se fait) et la \emph{loi morale universelle} (ce qui doit être). Pour Kant, \emph{aucune tradition ne peut légitimer la violation de l'impératif moral}.
\item Le silence de l'entourage (famille, voisins) constitue une \cle{complicité par omission} : ne pas secourir une personne en danger est une faute morale, et le droit positif camerounais réprime les violences au sein du couple.
\end{itemize}
\end{enumerate}

\begin{savaistu}[Cameroun : cadre juridique et réalités]
Le Code pénal camerounais (loi n°2016/007 du 12 juillet 2016) réprime expressément les \textbf{violences exercées sur le conjoint ou le partenaire} (notamment ses articles 277-2 et suivants). Le Cameroun a par ailleurs ratifié la \textbf{Convention sur l'élimination de toutes les formes de discrimination à l'égard des femmes} (CEDAW, 1994), qui oblige l'État à prévenir, punir et éradiquer les violences faites aux femmes. Pourtant, les enquêtes démographiques et de santé (EDS) montrent qu'une part importante des femmes en union déclare avoir subi des violences physiques, sexuelles ou émotionnelles de la part de leur conjoint. \textbf{L'écart entre le droit écrit et le droit vécu} illustre la persistance de la violence structurelle et symbolique.
\end{savaistu}

\courssub{Rédaction guidée : construire un paragraphe argumenté}

\begin{methode}[Structure d'un paragraphe type « thèse – argument – référence – exemple – mini-conclusion »]
\begin{enumerate}
\item \textbf{Thèse (phrase d'ouverture)} : réponse directe à la question, affirmée clairement.
\item \textbf{Argument (explication)} : développement logique, définition des concepts, enchaînement des raisons.
\item \textbf{Référence (autorité philosophique)} : citation ou paraphrase précise d'un auteur du cours (Kant, Weber, Arendt, Gandhi…).
\item \textbf{Exemple concret (camerounais si possible)} : illustration factuelle, localisée, qui \emph{éprouve} l'argument.
\item \textbf{Mini-conclusion (transition)} : résumé de l'apport du paragraphe et ouverture sur le suivant.
\end{enumerate}
\end{methode}

\begin{popfigure}
\begin{tikzpicture}[
  node distance=1.2cm and 1.5cm,
  box/.style={rectangle, draw=popblue, thick, fill=popblue!10, rounded corners, minimum width=2.8cm, minimum height=1.1cm, align=center, font=\small},
  arrow/.style={-Stealth, thick, popdark},
  labelbox/.style={rectangle, draw=poporange, fill=poporange!15, rounded corners, minimum width=1.6cm, minimum height=0.7cm, align=center, font=\scriptsize}
]
\node[box, align=center] (these){\textbf{THÈSE}\\\emph{« La violence domestique nie la dignité humaine. »}};
\node[box, right=of these, align=center] (arg){\textbf{ARGUMENT}\\\emph{Traiter autrui seulement comme moyen, non comme fin.}};
\node[box, right=of arg, align=center] (ref){\textbf{RÉFÉRENCE}\\\emph{Kant : « L'humanité est une fin en soi. »}};
\node[box, right=of ref, align=center] (ex){\textbf{EXEMPLE}\\\emph{Femme battue, silence de l'entourage.}};
\node[box, right=of ex, align=center] (ccl){\textbf{MINI-CCL}\\\emph{Donc : violation morale et juridique.}};

\draw[arrow] (these) -- (arg) node[midway, above, labelbox] {explique};
\draw[arrow] (arg) -- (ref) node[midway, above, labelbox] {fonde};
\draw[arrow] (ref) -- (ex) node[midway, above, labelbox] {illustre};
\draw[arrow] (ex) -- (ccl) node[midway, above, labelbox] {prouve};

\node[draw=popgreen, fill=popgreen!10, rounded corners, fit=(these)(ccl), inner sep=0.4cm, label={[font=\small\bfseries, popgreen]above:STRUCTURE D'UN PARAGRAPHE ARGUMENTÉ (TYPE EXAMEN)}] {};
\end{tikzpicture}
\legende{Schéma de la structure standard d'un paragraphe argumenté en dissertation ou en commentaire. Chaque flèche correspond à un connecteur logique (car, ainsi, par exemple, donc).}
\end{popfigure}

\begin{exoresolu}[Rédaction d'un paragraphe complet]
\textbf{Sujet} : « La violence peut-elle être juste ? » — Paragraphe sur la violence domestique.
\tcblower
\textbf{Production} :
\begin{quote}
\textbf{[Thèse]} La violence conjugale, fût-elle présentée comme « correction culturelle », est \emph{radicalement injuste} car elle nie la dignité de la victime. \textbf{[Argument]} En infligeant une souffrance physique et psychologique pour imposer une soumission, le mari traite son épouse comme un \emph{moyen} au service de son autorité et non comme une \emph{fin en soi}, c'est-à-dire comme un être raisonnable capable d'autonomie morale. \textbf{[Référence]} Kant affirme dans les \emph{Fondements de la métaphysique des mœurs} : « Agis de telle sorte que tu traites l'humanité, aussi bien dans ta personne que dans celle de tout autre, toujours en même temps comme une fin, et jamais simplement comme un moyen. » \textbf{[Exemple camerounais]} Dans bien des foyers camerounais, des épouses subissent des coups pour des « fautes » domestiques ; l'entourage valide souvent par le silence (« le mariage, c'est l'endurance »), transformant une coutume en complicité. \textbf{[Mini-conclusion]} Ainsi, ni la tradition, ni l'intimité du foyer ne sauraient légitimer cette violence : elle est une \emph{violation du droit moral universel} et, depuis la réforme de 2016, une \emph{infraction pénale} au Cameroun.
\end{quote}
\end{exoresolu}

\begin{attention}[Erreurs à ne pas commettre dans ce paragraphe]
\begin{itemize}
\item \textbf{Remplacer l'argument par l'exemple} : « Une femme se fait battre, c'est injuste. » $\rightarrow$ Pas d'analyse conceptuelle (dignité, moyen/fin), pas de référence.
\item \textbf{Citer Kant sans l'appliquer} : « Kant dit qu'il faut traiter l'homme comme une fin. » $\rightarrow$ Où est le lien avec le mari ? Avec les coups ? Avec la « culture » ?
\item \textbf{Oublier la mini-conclusion} : le paragraphe s'arrête sur l'exemple sans tirer la leçon philosophique ni préparer la suite.
\item \textbf{Mélanger les registres} : dire « c'est illégal donc c'est injuste » — la légalité ne suffit pas à fonder la justice (des lois injustes sont possibles). Il faut l'argument \emph{moral} (Kant) \emph{et} le rappel juridique.
\end{itemize}
\end{attention}

\courssub{Entraînement : définir les termes du sujet « La violence peut-elle être juste ? »}

\begin{methode}[Définition rigoureuse des termes — fiche de révision]
\begin{enumerate}
\item \textbf{Violence} (étym. lat. \emph{vis} : force, contrainte) : \emph{contrainte physique ou psychologique exercée sur autrui contre sa volonté, portant atteinte à son intégrité.} Distinguer : violence \emph{physique} / \emph{psychologique} / \emph{structurelle} / \emph{symbolique} ; violence \emph{individuelle} / \emph{collective} / \emph{étatique} ; violence \emph{légitime} (Weber) / \emph{illégitime}.
\item \textbf{Peut-elle} : double modalité — \emph{possibilité factuelle} (arrive-t-il qu'une violence soit juste ?) et \emph{légitimité normative} (existe-t-il des critères de justice applicables à la violence ?).
\item \textbf{Juste} (lat. \emph{justus} : conforme au droit) : \emph{conforme à la justice, au droit, à l'équité.} Distinguer : \emph{juste légal} (conforme à la loi positive) / \emph{juste moral} (conforme à la loi morale universelle, Kant) / \emph{juste politique} (conforme au bien commun).
\item \textbf{Articulation du sujet} : la question ne demande pas « La violence est-elle juste ? » (non, elle ne l'est pas par nature) mais « \emph{Peut-elle l'être} ? » — c'est-à-dire : \emph{existe-t-il des cas où l'usage de la contrainte est justifié par la justice elle-même ?} (légitime défense, action de la police, résistance à la tyrannie).
\end{enumerate}
\end{methode}

\begin{exemplebox}[Plan type de dissertation (Probatoire / Bac technique) — 3 parties équilibrées]
\textbf{Sujet :} « Toute violence est-elle injustifiable ? »

\textbf{Introduction} : accroche (ex. : « La violence est le dernier refuge de l'incompétence » — Isaac Asimov), définition des termes, problématique (la violence, par essence destructrice, peut-elle trouver une justification morale ou juridique ?), annonce du plan.

\textbf{I. La violence est \emph{par nature} injustifiable (condamnation morale)}
\begin{itemize}
\item A. Violence = atteinte à l'intégrité d'autrui, destruction du pouvoir d'agir ensemble (Arendt), négation de la dignité (Kant).
\item B. Une fin juste ne peut être atteinte par des moyens injustes (Gandhi : la violence engendre la violence).
\item C. Même l'État doit limiter sa force (État de droit, proportionnalité).
\end{itemize}

\textbf{II. Pourtant, certaines situations \emph{exigent} la contrainte (justifications possibles)}
\begin{itemize}
\item A. Légitime défense : violence \emph{réactive}, proportionnée, pour sauver une vie (droit naturel, reconnue par le Code pénal camerounais).
\item B. Monopole de la violence légitime (Weber) : police, justice — contrainte \emph{encadrée par la loi} pour protéger la liberté de tous.
\item C. Résistance à l'oppression : violence \emph{libératrice} contre une violence \emph{structurelle} tyrannique, lorsque toute voie légale est fermée.
\end{itemize}

\textbf{III. La justice ne \emph{justifie} la contrainte qu'à des conditions strictes (critères de discernement)}
\begin{itemize}
\item A. Nécessité (dernier recours), proportionnalité (moyens adaptés), finalité (restaurer le droit, non détruire).
\item B. Distinction pouvoir/violence (Arendt) : le \emph{pouvoir} (agir ensemble selon des règles) est juste ; la \emph{violence} (instrumentalité destructive) ne l'est jamais pleinement.
\item C. Idéal de non-violence active (Gandhi, Martin Luther King) : la justice vraie vise à \emph{supprimer} la violence, non à la réguler.
\end{itemize}

\textbf{Conclusion} : bilan nuancé — la violence n'est \emph{jamais bonne} ; elle est \emph{parfois inévitable} (légitime défense, police juste) mais \emph{toujours à regretter et à limiter}. Ouverture : construire une société juste, c'est \emph{réduire au minimum} le recours à la contrainte physique.
\end{exemplebox}

\courssub{Correction collective : critères d'évaluation au Baccalauréat technique}

\begin{propriete}[Grille d'évaluation type dissertation de philosophie (sur 20 points)]
\begin{center}
\begin{tabularx}{\linewidth}{|l|X|c|}\hline
\rowcolor{popblueL}\textbf{Critère} & \textbf{Descriptif (niveau « Assez bien / Bien »)} & \textbf{Points} \\ \hline
\cle{Pertinence} & Sujet analysé, termes définis, problématique explicite, plan annoncé et respecté. Pas de hors-sujet. & 4 \\ \hline
\cle{Cohérence} & Enchaînement logique des parties et des paragraphes, transitions claires, thèse directrice tenue du début à la fin. & 4 \\ \hline
\cle{Connaissances / Références} & Auteurs du cours cités \emph{à bon escient} (Kant, Weber, Arendt, Gandhi…), concepts maniés avec rigueur. & 5 \\ \hline
\cle{Argumentation} & Arguments explicités (pas d'affirmations gratuites), exemples concrets (camerounais si possible) \emph{analysés}, objections anticipées et traitées. & 4 \\ \hline
\cle{Expression} & Français correct (syntaxe, orthographe, vocabulaire philosophique), style académique, connecteurs logiques (donc, cependant, en effet, par conséquent). & 3 \\ \hline
\end{tabularx}
\end{center}
\end{propriete}

\begin{exercice}[Entraînement chronométré — 1 h 30]
\textbf{Sujet :} « La violence peut-elle être juste ? »

\textbf{Consignes} :
\begin{enumerate}
\item Rédigez l'\textbf{introduction complète} (accroche, définitions, problématique, plan) — 20 min.
\item Rédigez \textbf{deux paragraphes argumentés} (un pour la partie I, un pour la partie II) en respectant la structure \emph{Thèse – Argument – Référence – Exemple – Mini-conclusion} — 50 min.
\item Relisez-vous avec la \textbf{grille ci-dessus} : cochez chaque critère — 20 min.
\end{enumerate}
\end{exercice}

\begin{corrige}[Exercice n°1 — Éléments attendus pour l'introduction]
\begin{itemize}
\item \textbf{Accroche} : citation (Arendt, Gandhi, Asimov) ou fait marquant (manifestation réprimée, racket, violences conjugales).
\item \textbf{Définitions} : violence (contrainte exercée contre la volonté d'autrui), juste (conforme au droit et à la morale), distinction violence/force légitime.
\item \textbf{Problématique} : « Si la violence nie par essence la liberté et la dignité, comment comprendre qu'elle soit parfois invoquée au nom de la justice (légitime défense, police, résistance) ? La justice peut-elle \emph{user} de violence sans se contredire ? »
\item \textbf{Plan} : I. La violence est par nature injustifiable (Kant, Arendt, Gandhi). II. Des cas où la contrainte s'impose (légitime défense, monopole étatique weberien, résistance à la tyrannie). III. Critères pour ne pas confondre force légitime et violence illégitime (nécessité, proportionnalité, finalité, idéal non-violent).
\end{itemize}
\end{corrige}

\begin{aretenir}[Check-list finale avant l'épreuve]
\begin{itemize}
\item \textbf{Je définis} chaque terme du sujet dans l'introduction.
\item \textbf{Je cite} au moins 3 auteurs du cours \emph{avec leurs concepts précis} (pas seulement le nom).
\item \textbf{Je donne} au moins 2 exemples camerounais \emph{analysés} (pas seulement mentionnés).
\item \textbf{Je structure} chaque paragraphe : Thèse – Argument – Référence – Exemple – Mini-conclusion.
\item \textbf{Je fais} des transitions explicites entre les parties (« Cependant », « Dès lors », « C'est pourquoi »).
\item \textbf{Je conclus} par un bilan nuancé et une ouverture (ex. : éducation à la non-violence, médiation, accès effectif à la justice).
\item \textbf{Je relis} : orthographe, accords, connecteurs, absence de langage familier (« c'est nul », « grave ») — le vocabulaire philosophique est exigé.
\end{itemize}
\end{aretenir}

\begin{savaistu}[Philosopher au Cameroun aujourd'hui]
La philosophie en classe de Première technique ne forme pas seulement à l'examen : elle outille pour \textbf{penser les violences ordinaires} (racket, violences conjugales, répression, harcèlement scolaire) et \textbf{y répondre par le droit, la parole et l'organisation collective}. Chaque élève qui sait distinguer \emph{force légitime} et \emph{violence illégitime}, qui sait nommer la \emph{violence symbolique} derrière une « coutume », qui sait construire un paragraphe \emph{Thèse – Argument – Référence – Exemple}, devient un citoyen plus lucide, plus libre, plus capable de \emph{transformer} le réel que de le subir. C'est tout l'enjeu de cet atelier — et de votre année de Première.
\end{savaistu}

\newpage

\coursec{Activité d'intégration}

\courssub{Objectifs de l'activité}
Mobiliser l'ensemble des savoirs de la leçon (\cle{définition}, \cle{origine}, \cle{types}, \cle{procès} de la violence) pour construire et justifier, oralement et par écrit, un jugement argumenté sur un cas concret complexe. L'élève doit être capable de :
\begin{itemize}
    \item identifier et qualifier les formes de violence présentes dans une situation réelle (physique, psychologique, structurelle, symbolique) ;
    \item articuler les thèses philosophiques étudiées (Arendt, Kant, Weber, Gandhi) pour étayer une position (condamnation ou justification conditionnelle) ;
    \item rédiger un \cle{verdict motivé} respectant la structure dialectique de la conclusion de dissertation (thèse / antithèse / dépassement) ;
    \item coopérer dans un dispositif de \cle{tribunal philosophique} en respectant les rôles (accusation, défense, jury) et la contradiction argumentée.
\end{itemize}

\courssub{Situation}
\textbf{Contexte :} Douala, quartier New-Bell, mars 2025. Un fait divers choque l'opinion publique et alimente les débats sur les réseaux sociaux et dans les établissements scolaires.

\textbf{Dossier documentaire fourni aux trois groupes :}

\begin{exemplebox}[Document 1 : Article de presse (extrait fictif inspiré de faits réels) — « Lynchage à New-Bell : la foule fait justice »]
\emph{« Samedi 15 mars, vers 14h, un jeune homme d'une vingtaine d'années, identifié comme K. M., a été intercepté par des riverains du carrefour Shell alors qu'il tentait de s'enfuir après avoir arraché le sac à main d'une commerçante. La victime, Mme N., criant au voleur, a alerté les passants. En quelques minutes, une foule compacte s'est formée. Dépassant la simple interpellation, plusieurs individus ont frappé le suspect à coups de pierres, de bâtons et de ceintures. La police, alertée, est arrivée 25 minutes plus tard, trouvant le jeune homme inanimé, le visage tuméfié. Transporté à l'hôpital Laquintinie, son pronostic vital était engagé dimanche matin. Témoignage d'un riverain, M. E. : « On en a marre de l'insécurité. La police ne fait rien, les voleurs sortent de prison le lendemain. Nous, on fait le travail. » Le Procureur de la République a ouvert une enquête pour « violences volontaires en réunion ». »}
\end{exemplebox}

\begin{exemplebox}[Document 2 : Max Weber, \emph{Le Savant et le Politique} (1919), extrait]
\emph{« L'État est une communauté humaine qui, dans les limites d'un territoire déterminé — il faut souligner ce trait : un territoire — revendique avec succès le monopole de la violence physique légitime pour tous ceux qui se trouvent dans ce territoire. [...] Le monopole de la violence physique légitime est le trait distinctif de l'État moderne. »}
\end{exemplebox}

\begin{exemplebox}[Document 3 : Gandhi, écrits sur la non-violence (1920-1940), extraits]
\emph{« La non-violence n'est pas une arme des faibles ; c'est l'arme des forts. [...] Elle ne cherche pas à détruire l'adversaire, mais à le convertir. [...] La violence, même lorsqu'elle semble faire le bien, fait un bien temporaire ; le mal qu'elle fait est permanent. »}
\end{exemplebox}

\begin{exemplebox}[Document 4 : Données statistiques — Violences en milieu scolaire au Cameroun (enquête fictive d'entraînement, inspirée des rapports MINESEC / UNICEF, échantillon 12 000 élèves)]
\begin{center}
\begin{tabularx}{\linewidth}{|l|X|X|}\hline
\rowcolor{popblueL}
\textbf{Type de violence} & \textbf{Prévalence (\% d'élèves déclarant en avoir été victime ou témoin)} & \textbf{Auteur principal identifié} \\ \hline
Physique (coups, blessures) & 38\% & Camarades (62\%), Personnel éducatif (28\%) \\ \hline
Verbale / Psychologique (insultes, menaces, harcèlement, cyberharcèlement) & 54\% & Camarades (70\%), Personnel (20\%) \\ \hline
Sexuelle (attouchements, propos obscènes, viol) & 12\% & Personnel (45\%), Camarades (35\%), Externe (20\%) \\ \hline
Structurelle (exclusion, corruption notes/argent, travail forcé) & 27\% & Administration / Personnel (80\%) \\ \hline
\end{tabularx}
\end{center}
\legende{Source : tableau d'entraînement construit pour la leçon, inspiré des enquêtes MINESEC-UNICEF sur les violences en milieu scolaire. Les pourcentages ne sont pas exclusifs (poly-victimation) : un même élève peut déclarer plusieurs types de violence.}
\end{exemplebox}

\begin{attention}[Lecture des données du Document 4]
Les pourcentages de la colonne « Prévalence » (38\%, 54\%, 12\%, 27\%) totalisent plus de 100\% : c'est normal, car un élève peut être victime de plusieurs formes de violence à la fois. De même, les répartitions par auteur (par exemple 62\% + 28\% = 90\% pour la violence physique) ne totalisent pas 100\% car il existe d'autres auteurs (personnes extérieures, inconnus) non détaillés dans le tableau. Ne jamais additionner ces colonnes pour conclure.
\end{attention}

\courssub{Consignes}
\textbf{Organisation :} La classe est divisée en trois groupes de taille équivalente. Chaque groupe dispose de 45 minutes pour préparer sa prestation à partir du dossier.

\begin{enumerate}
    \item \textbf{Groupe « Accusation » (Condamnation radicale) :} À partir de \textbf{Kant} (l'homme comme fin, interdiction de se faire justice soi-même) et \textbf{Hannah Arendt} (\emph{Du mensonge à la violence} : la violence est muette, destructrice du pouvoir politique, jamais légitime), vous devez démontrer que \textbf{aucune violence ne saurait être légitime}, ni celle de la foule (Doc 1), ni celle de l'État si elle excède le droit (Doc 2), ni celle subie à l'école (Doc 4). Préparez une plaidoirie écrite de 300 mots et une restitution orale de 5 min.
    \item \textbf{Groupe « Défense » (Justification conditionnelle) :} À partir de \textbf{Weber} (monopole étatique, Doc 2), du droit positif (\cle{légitime défense}, reconnue par le Code pénal camerounais sous les conditions de nécessité et de proportionnalité), et de la \cle{résistance à l'oppression} (Déclaration des Droits de l'Homme et du Citoyen de 1789, art. 2 ; Fanon sur la violence libératrice), vous devez montrer que \textbf{certaines violences sont justifiées, voire nécessaires} : violence policière légale, légitime défense de la commerçante, résistance à la violence structurelle scolaire (Doc 4). Préparez une plaidoirie écrite de 300 mots et une restitution orale de 5 min.
    \item \textbf{Groupe « Jury » (Synthèse et Verdict) :} Vous ne prenez pas parti a priori. Votre tâche :
    \begin{itemize}
        \item Classer les arguments des deux camps dans un \textbf{tableau dialectique} (Thèse / Antithèse / Dépassement).
        \item Vérifier la pertinence des références (les citations correspondent-elles aux thèses réelles des auteurs ?).
        \item Rédiger un \textbf{verdict motivé d'une demi-page (environ 250 mots)} respectant la méthode de la conclusion de dissertation : rappel du problème, résumé dialectique, dépassement (position propre nuancée), ouverture.
    \end{itemize}
    Préparez la restitution orale du verdict (5 min) et l'affichage du tableau au tableau.
    \item \textbf{Restitution collective (30 min) :} Chaque groupe présente sa production (5 min). Le Jury affiche son tableau dialectique et lit son verdict. L'enseignant anime la synthèse finale : correction des approximations, rappel des distinctions conceptuelles (violence/force, légitimité/légalité, violence physique/structurelle).
\end{enumerate}

\courssub{Ressources à mobiliser (Rappel du cours)}
\begin{itemize}
    \item \textbf{Définition :} Violence = contrainte physique ou morale qui porte atteinte à l'intégrité d'un être (étymologie latine \emph{vis} : force dévastatrice). Distinction \cle{violence} (subie, destructrice) / \cle{force} (capacité d'agir, pouvoir).
    \item \textbf{Origine :} \cle{Naturelle} (instinct agressif, Lorenz, pulsion de mort chez Freud) vs \cle{Culturelle} (apprentissage social, désir mimétique chez Girard, construction historique).
    \item \textbf{Types :} \cle{Physique}, \cle{Psychologique} (verbale, harcèlement), \cle{Structurelle} (Galtung : inégalités inscrites dans les structures sociales), \cle{Symbolique} (Bourdieu : imposition de catégories de pensée dominantes).
    \item \textbf{Procès (Condamnation) :} Kant (impératif catégorique : ne jamais traiter l'humanité comme simple moyen), Arendt (violence \emph{instrumentale}, destructrice du \emph{pouvoir} qui naît de la parole et de la concertation), Gandhi (non-violence active, \emph{Satyagraha}).
    \item \textbf{Procès (Justification) :} Weber (monopole de la violence \emph{légitime} = définition de l'État), Légitime défense (nécessité, proportionnalité, immédiateté), Résistance à l'oppression (droit naturel ; Fanon : la violence de l'opprimé comme restauration de sa dignité).
\end{itemize}

\begin{exoresolu}[Corrigé commenté et verdict type]
\textbf{1. Analyse des documents (Travail préalable commun)}
\begin{itemize}
    \item \textbf{Doc 1 :} Violence \cle{physique} collective, \cle{privée} (justice populaire), \cle{illégale} (lynchage), motivée par le sentiment d'impunité (défaillance étatique perçue). Victime : un suspect dont la présomption d'innocence est bafouée.
    \item \textbf{Doc 2 :} Weber définit l'État moderne par son \cle{monopole de la violence physique légitime}. La violence n'est légitime que si elle émane de l'État \emph{et} respecte le droit. La foule (Doc 1) usurpe ce monopole.
    \item \textbf{Doc 3 :} Gandhi refuse la violence comme moyen, même pour une fin juste. La non-violence est une \cle{force morale} supérieure, qui vise la conversion de l'adversaire et non sa destruction.
    \item \textbf{Doc 4 :} Révèle la \cle{violence structurelle} et \cle{symbolique} à l'école (corruption, abus d'autorité). 38\% de violence physique, 54\% de violence psychologique. L'école, lieu de socialisation, devient aussi un lieu de reproduction de la violence.
\end{itemize}

\textbf{2. Production du Groupe « Accusation » (Extrait de plaidoirie)}
\emph{« Madame la Présidente, le lynchage de New-Bell (Doc 1) est l'illustration paroxystique de la \cle{violence privée} que Kant condamne au nom de l'impératif catégorique : « Agis de telle sorte que tu traites l'humanité [...] toujours en même temps comme une fin, et jamais simplement comme un moyen ». La foule a traité le suspect comme une chose, niant sa dignité d'homme et sa présomption d'innocence. Hannah Arendt nous enseigne que la violence est « muette » : elle ne persuade pas, elle écrase. Elle détruit le \emph{pouvoir} (le « pouvoir d'agir de concert ») pour le remplacer par la contrainte brute. Même l'État, s'il use de violence hors du droit (torture, bavures), perd sa légitimité (Weber). Quant à l'école (Doc 4), les 54\% de violences psychologiques prouvent que la violence structurelle broie l'avenir des élèves. La non-violence gandhienne n'est pas passivité : c'est la seule voie qui brise le cycle « violence $\rightarrow$ contre-violence ». Condamnons toute violence sans exception. »}

\textbf{3. Production du Groupe « Défense » (Extrait de plaidoirie)}
\emph{« Madame la Présidente, l'abstraction moralisante de l'Accusation ignore la réalité du terrain. Weber (Doc 2) pose le principe de réalité : l'État \emph{doit} détenir le monopole de la violence \emph{légitime} pour protéger les citoyens. Sans police armée, la commerçante du Doc 1 n'a que la fuite ou la légitime défense, reconnue par le droit pénal à condition que la riposte soit nécessaire, immédiate et proportionnée à l'agression. La violence de l'État n'est pas « violence » au sens destructeur : c'est de la \cle{force publique} encadrée par le droit. Ensuite, face à la violence structurelle scolaire (Doc 4 : 12\% de violences sexuelles, dont 45\% imputées au personnel !), la soumission n'est pas une vertu. La résistance à l'oppression est un droit sacré (DDHC 1789, art. 2). Fanon a montré que la violence libératrice restaure la dignité de l'opprimé. Gandhi lui-même distinguait la lâcheté (ne pas réagir) de la non-violence active (résister sans tuer). Nous défendons la violence \emph{légitime, proportionnée, encadrée par le droit ou la nécessité morale}. »}

\textbf{4. Travail du Groupe « Jury » : Tableau dialectique}
\begin{center}
\begin{tabularx}{\linewidth}{|X|X|X|}\hline
\rowcolor{popblueL}
\textbf{THÈSE : La violence est radicalement illégitime} & \textbf{ANTITHÈSE : La violence peut être légitime (sous conditions)} & \textbf{DÉPASSEMENT : La légitimité par le Droit et la non-violence active} \\ \hline
\begin{itemize}
    \item Kant : l'homme = fin en soi. Violence = instrumentalisation.
    \item Arendt : la violence détruit le pouvoir (parole et action concertée). Elle est « muette ».
    \item Gandhi : violence = mal permanent. Non-violence = force des forts (\emph{Satyagraha}).
    \item Doc 1 : lynchage = justice privée, arbitraire, cruelle.
    \item Doc 4 : violence scolaire = échec éducatif, reproduction sociale.
\end{itemize} & \begin{itemize}
    \item Weber : État = monopole de la violence \emph{légitime} (légalité + territoire).
    \item Légitime défense (droit) : nécessité, proportionnalité, immédiateté.
    \item Résistance à l'oppression (DDHC 1789, Fanon) : désobéissance civile / violence libératrice.
    \item Doc 1 : commerçante en légitime défense ? Police nécessaire.
    \item Doc 4 : résistance aux abus (corruption, viol) = légitime.
\end{itemize} & \begin{itemize}
    \item \textbf{Distinction cruciale :} \cle{violence} (arbitraire, destructrice) vs \cle{force publique} (monopole étatique, encadrée par le \emph{Droit}, contrôlée).
    \item La \cle{légitimité} ne va pas de soi : elle suppose \textbf{procédure}, \textbf{proportionnalité}, \textbf{contrôle démocratique}.
    \item La \cle{non-violence} (Gandhi) n'est pas l'absence de conflit, mais un \textbf{mode d'action politique} (grève, désobéissance, sit-in) qui vise la conversion, non la destruction.
    \item \textbf{Solution pour Doc 1 :} police formée, justice rapide $\rightarrow$ retire le « besoin » de justice populaire.
    \item \textbf{Solution pour Doc 4 :} éducation à la citoyenneté, sanctions disciplinaires et judiciaires fermes, protection des victimes.
\end{itemize} \\ \hline
\end{tabularx}
\end{center}

\textbf{5. Verdict motivé (Rédigé par le Jury — Modèle attendu)}

« Le tribunal philosophique, après avoir entendu l'Accusation et la Défense, rend le verdict suivant.

\textbf{Problème :} La violence est-elle parfois légitime, ou doit-elle être condamnée sans réserve ?

\textbf{Résumé dialectique :} La thèse (Accusation), s'appuyant sur Kant et Arendt, établit que la violence, par nature, nie la dignité humaine (traiter l'homme comme moyen) et détruit le tissu politique (le pouvoir naît de la parole, non de la contrainte). Le lynchage de New-Bell (Doc 1) et les violences scolaires (Doc 4) en sont la preuve tragique : violence privée arbitraire et violence structurelle institutionnelle. L'antithèse (Défense), appuyée sur Weber et le droit positif, démontre qu'une violence \emph{encadrée par la loi} (monopole étatique, légitime défense) est la condition de la sécurité et de la protection des droits. Sans force publique légitime, c'est la loi du plus fort (lynchage) qui règne.

\textbf{Dépassement :} La légitimité ne réside pas dans la violence elle-même, mais dans sa \textbf{soumission au Droit}. La distinction opérée par Weber est décisive : l'État moderne se définit par le \emph{monopole de la violence physique légitime}. « Légitime » signifie : prévue par la loi, exécutée par des agents formés, soumise au contrôle juridictionnel, proportionnée. Le lynchage (Doc 1) est illégitime car il usurpe ce monopole et ignore la procédure. La violence scolaire (Doc 4) est illégitime car elle viole la loi et la mission éducative. Inversement, l'intervention policière pour arrêter le voleur, ou la sanction judiciaire du harceleur scolaire, sont des \emph{forces publiques légitimes}, non des violences au sens destructeur.

\textbf{Position nuancée :} Nous ne justifions pas la violence, nous légitimons la \textbf{contrainte légale}. La non-violence gandhienne reste l'horizon éthique supérieur : elle oblige à inventer des modes de résistance (désobéissance civile, médiation, éducation) qui brisent la spirale mimétique sans détruire l'adversaire.

\textbf{Ouverture :} Comment former, au Cameroun, une police et une école qui incarnent la \emph{force de la loi} plutôt que la \emph{loi de la force} ? C'est le défi citoyen de notre génération. »
\end{exoresolu}

\courssub{Bilan de l'activité}
Cette activité a permis de vérifier l'appropriation des savoirs de l'unité à travers une mise en situation authentique :
\begin{itemize}
    \item \textbf{Maîtrise conceptuelle :} Les élèves distinguent désormais \cle{violence} (arbitraire, privée, destructrice) et \cle{force publique légitime} (monopole étatique, encadrée par le droit), ainsi que les quatre types (physique, psychologique, structurelle, symbolique) repérés dans le Doc 4.
    \item \textbf{Argumentation philosophique :} Ils mobilisent correctement les auteurs au programme (Kant, Arendt, Weber, Gandhi) sans les déformer, en articulant citations et concepts (impératif catégorique, monopole légitime, \emph{Satyagraha}, pouvoir vs violence).
    \item \textbf{Méthodologie dissertation :} Le verdict du Jury valide la structure \emph{Thèse / Antithèse / Dépassement / Ouverture}, exigée au Probatoire et au Baccalauréat technique. La rédaction du tableau dialectique au tableau a structuré la pensée collective.
    \item \textbf{Ancrage camerounais :} Le traitement du lynchage (phénomène récurrent à Douala et Yaoundé) et des violences scolaires (données inspirées des enquêtes MINESEC) ancre la philosophie dans le réel, développant l'esprit critique citoyen.
    \item \textbf{Compétences transverses :} Expression orale (plaidoirie, verdict), travail collaboratif (rôles définis), éthique du débat (écoute, contradiction respectueuse).
\end{itemize}
\emph{Point de vigilance pour l'enseignant :} Veiller à ce que le groupe « Défense » ne glisse pas vers une apologie de la violence brute, mais défende strictement la \emph{contrainte légale} (Weber) ou la \emph{résistance non-violente active} (Gandhi, Fanon nuancé). Reprendre en synthèse la distinction \emph{légalité / légitimité}.

\newpage

\coursec{Méthodes et exercices résolus}

\courssub{Fiche méthode 1 : Appliquer les notions de la leçon à une situation concrète de la vie courante}

L'un des savoir-faire exigés au Probatoire et au Baccalauréat technique consiste à appliquer les notions de l'unité à des situations concrètes de la vie courante. Cette compétence suppose une démarche rigoureuse, que nous détaillons ici.

\begin{methode}[Analyser une situation concrète avec les notions de la leçon]
Face à un fait de la vie courante (une bagarre dans un quartier de Douala, une dispute sur un terrain à Bafoussam, une répression policière, une moquerie répétée sur les réseaux sociaux), la démarche à suivre est rigoureuse :
\begin{enumerate}
\item \textbf{Décrire la situation avec précision} : qui agit ? sur qui ? par quels moyens (coups, paroles, loi, humiliation) ? dans quel but ?
\item \textbf{Identifier la notion mobilisée} : s'agit-il de violence ? De quel type (physique, psychologique, verbale, économique, légitime ou illégitime) ?
\item \textbf{Appliquer la définition} : vérifier que la situation remplit les critères de la violence (usage de la force, contrainte, atteinte à l'intégrité ou à la liberté d'autrui).
\item \textbf{Interroger l'origine} : l'acte relève-t-il d'une pulsion naturelle (agressivité) ou d'un apprentissage culturel (éducation, milieu social) ?
\item \textbf{Porter un jugement argumenté} : la violence est-elle condamnable ? Peut-elle être justifiée (légitime défense, violence de l'État) ? Justifier la réponse par un auteur ou un argument.
\item \textbf{Conclure} par une phrase claire qui répond à la question posée.
\end{enumerate}
\textbf{Réflexe} : toujours passer du concret à la notion, puis de la notion au jugement. Ne jamais rester au simple récit.
\end{methode}

\begin{exoresolu}[Analyse d'une situation concrète]
\textbf{Énoncé.} À Yaoundé, un jeune vendeur à la sauvette se fait confisquer sa marchandise par des agents municipaux. Blessé dans sa dignité, il frappe l'un des agents. Un témoin déclare : « C'est lui le violent ». Analysez cette situation en appliquant les notions de la leçon (définition et types de violence).
\tcblower
\textbf{Solution détaillée.}
\begin{enumerate}
\item \textbf{Description.} Deux actes sont en présence : (a) la confiscation de la marchandise par les agents, accompagnée d'une humiliation ; (b) le coup porté par le vendeur à l'agent.
\item \textbf{Notion mobilisée.} La violence se définit comme l'usage de la force, physique ou symbolique, pour contraindre autrui et porter atteinte à son intégrité, à sa liberté ou à sa dignité.
\item \textbf{Application de la définition.} Le coup porté par le vendeur est une \cle{violence physique} directe : il remplit tous les critères (force, atteinte à l'intégrité). Mais la confiscation brutale et l'humiliation constituent une \cle{violence institutionnelle et psychologique} : elle s'exerce au nom d'une autorité et atteint la dignité et les moyens d'existence du vendeur (violence économique).
\item \textbf{Origine.} La réaction du vendeur peut s'expliquer par une agressivité naturelle face à la frustration, mais aussi par un contexte social (précarité, rapports de force habituels) : l'origine de la violence est à la fois naturelle et culturelle.
\item \textbf{Jugement argumenté.} Le témoin ne voit que la violence visible et illégitime (le coup). Or, comme le montre la distinction classique entre violence illégitime et violence légitime, l'État détient le monopole de la violence légitime (Max Weber) : les agents agissent au nom de la loi. Mais la légalité n'efface pas la brutalité : une violence peut être légale et pourtant excessive ou humiliante.
\item \textbf{Conclusion.} La situation montre que la violence ne se réduit pas aux coups : il existe des violences invisibles (institutionnelles, économiques, psychologiques). Le témoin a tort de n'en voir qu'une forme ; une analyse philosophique complète distingue les types de violence en présence avant de juger.
\end{enumerate}
\end{exoresolu}

\courssub{Fiche méthode 2 : Définir une notion sans erreur}

Toute dissertation commence par la définition des termes du sujet. Définir correctement la violence est donc une étape décisive : une définition fausse ou vague compromet l'ensemble du devoir.

\begin{methode}[Construire une définition philosophique correcte]
Définir la violence est souvent la première étape d'une dissertation. Méthode :
\begin{enumerate}
\item \textbf{Partir du sens commun} : la violence, c'est « frapper », « faire mal ».
\item \textbf{Critiquer le sens commun} : montrer qu'il est trop étroit (il oublie les violences sans coups : insultes, exclusion, misère) ou trop large (toute force n'est pas violence : le chirurgien use de force sans violence).
\item \textbf{Dégager les critères essentiels} : usage de la force ou de la contrainte ; intention de nuire ou de dominer ; atteinte à l'intégrité physique, morale ou à la liberté d'autrui.
\item \textbf{Formuler la définition} en une phrase complète et générale.
\item \textbf{Distinguer les notions voisines} : force, puissance, contrainte légitime, conflit.
\end{enumerate}
\textbf{Formule à retenir} : « La violence est l'usage de la force, physique ou symbolique, pour contraindre autrui, en portant atteinte à son intégrité, à sa liberté ou à sa dignité. »
\end{methode}

\begin{exoresolu}[Exercice de définition]
\textbf{Énoncé.} Un élève écrit : « La violence, c'est la force. » Cette définition est-elle satisfaisante ? Corrigez-la en suivant la méthode de définition.
\tcblower
\textbf{Solution détaillée.}
\begin{enumerate}
\item \textbf{Sens commun.} L'élève identifie violence et force : frapper, c'est employer sa force.
\item \textbf{Critique.} Cette définition est fautive par excès : toute force n'est pas violence. Le père qui porte son enfant, le médecin qui opère, l'enseignant qui corrige une copie usent d'une forme de force ou d'autorité sans violence. Inversement, elle est fautive par défaut : une insulte, une rumeur, une exclusion peuvent être violentes sans aucune force physique.
\item \textbf{Critères essentiels.} La violence suppose : (a) une contrainte exercée sur autrui ; (b) contre sa volonté ; (c) avec atteinte à son intégrité, sa liberté ou sa dignité.
\item \textbf{Définition corrigée.} « La violence est l'usage illégitime ou destructeur de la force, physique ou symbolique, visant à contraindre autrui et portant atteinte à son intégrité, à sa liberté ou à sa dignité. »
\item \textbf{Distinction.} La \cle{force} est un simple pouvoir d'agir ; elle devient \cle{violence} lorsqu'elle s'exerce contre autrui pour le nier comme sujet. La force peut être au service du droit ; la violence, elle, détruit la relation.
\item \textbf{Conclusion.} La définition de l'élève est insuffisante car elle confond la violence avec son simple moyen. Une bonne définition philosophique précise la finalité (contraindre, nuire) et l'effet (atteinte à autrui).
\end{enumerate}
\end{exoresolu}

\begin{aretenir}[Définition et distinctions essentielles]
La \cle{violence} est l'usage de la force, physique ou symbolique, pour contraindre autrui contre sa volonté, en portant atteinte à son intégrité, à sa liberté ou à sa dignité. On distingue : la \cle{force} (simple pouvoir d'agir, neutre en soi), la \cle{violence illégitime} (agression privée, arbitraire) et la \cle{force légitime} (violence réglée par le droit, comme celle de l'État selon Max Weber). Les types de violence sont : physique, psychologique (ou morale), verbale, économique et institutionnelle.
\end{aretenir}

\courssub{Fiche méthode 3 : Rédiger et justifier une réponse argumentée}

Savoir définir ne suffit pas : l'examen exige aussi de rédiger et de justifier une réponse ou une conclusion. C'est l'objet de cette troisième fiche.

\begin{methode}[Justifier une réponse ou une conclusion en dissertation]
Pour rédiger et justifier une réponse à une question de philosophie (ex. « La violence est-elle toujours condamnable ? ») :
\begin{enumerate}
\item \textbf{Annoncer clairement sa thèse} en une phrase (« Nous soutiendrons que... »).
\item \textbf{Formuler un argument} : un raisonnement qui établit la thèse (par exemple : la violence nie la liberté d'autrui, or la liberté est la dignité même de l'homme).
\item \textbf{Appuyer l'argument sur un auteur} cité brièvement et exactement (Hobbes, Rousseau, Weber, Arendt...), sans le laisser parler à votre place.
\item \textbf{Illustrer par un exemple concret}, de préférence camerounais ou proche de l'expérience des élèves.
\item \textbf{Envisager l'objection} : que répondrait l'adversaire ? (ex. « Mais la légitime défense ? »).
\item \textbf{Réfuter ou nuancer} l'objection, puis \textbf{conclure} en réaffirmant la thèse justifiée.
\end{enumerate}
\textbf{Réflexe} : une affirmation sans argument ni exemple vaut zéro point de justification. Chaque idée doit être prouvée.
\end{methode}

\begin{exoresolu}[Rédaction d'un paragraphe argumenté]
\textbf{Énoncé.} Rédigez un paragraphe argumenté (thèse, argument, auteur, exemple, objection, réponse) répondant à la question : « La violence de l'État est-elle une violence comme les autres ? »
\tcblower
\textbf{Solution détaillée (paragraphe rédigé).}

\emph{Thèse.} La violence de l'État n'est pas une violence comme les autres, car elle est en principe une violence légitime, au service du droit.

\emph{Argument.} Ce qui rend une violence condamnable, c'est qu'elle s'exerce arbitrairement, pour l'intérêt privé de celui qui frappe. Or la violence de l'État est réglée par la loi, s'applique à tous de la même manière et vise la protection de tous : elle transforme la force brute en force ordonnée.

\emph{Appui sur un auteur.} Max Weber définit l'État comme l'instance qui revendique le monopole de la violence physique légitime : seul l'État peut user de la force (police, justice, armée), et c'est précisément ce monopole qui met fin à la violence privée de chacun contre chacun.

\emph{Exemple.} Au Cameroun, lorsque la gendarmerie arrête un bandit qui rançonnait les voyageurs sur la route Yaoundé--Bafoussam, elle use de la force ; mais cette force protège les usagers et s'exerce au nom de la loi, non pour un profit personnel.

\emph{Objection.} On objectera que la violence d'État peut être abusive : tabassage en détention, répression excessive d'une manifestation.

\emph{Réponse.} C'est exact : mais l'abus ne supprime pas la distinction, il la confirme. Une violence d'État qui viole la loi cesse d'être légitime et redevient une violence comme les autres, condamnable au même titre. C'est donc le respect du droit qui fait la différence, non l'identité de celui qui frappe.

\emph{Conclusion.} La violence de l'État se distingue de la violence privée par sa légalité et sa finalité commune ; mais elle n'est légitime que dans les limites du droit, faute de quoi elle retombe dans la violence ordinaire.
\end{exoresolu}

\courssub{Fiche méthode 4 : Problématiser une question d'examen}

Avant même de rédiger, il faut problématiser : c'est la première étape de l'introduction de dissertation, et c'est souvent elle qui détermine la note.

\begin{methode}[Construire une problématique sur la violence]
Une bonne introduction de dissertation repose sur une problématique, c'est-à-dire une tension entre deux réponses possibles.
\begin{enumerate}
\item \textbf{Dégager le sens commun} du mot du sujet (ex. « la violence est mal »).
\item \textbf{Trouver la difficulté} : un cas où le sens commun vacille (la légitime défense, la révolte contre l'injustice, la violence éducative).
\item \textbf{Formuler la tension} sous forme d'alternative : « Si toute violence est condamnable, comment justifier... ? Mais si certaines violences sont justes, où tracer la limite ? »
\item \textbf{Poser la question directrice} en une phrase interrogative qui annonce le plan.
\end{enumerate}
\textbf{Formule à retenir} : problématique = tension entre deux thèses + question qui les dépasse.
\end{methode}

\begin{exoresolu}[Construction d'une problématique]
\textbf{Énoncé.} Sujet : « Faut-il condamner toute violence ? » Construisez la problématique de ce sujet en quatre étapes.
\tcblower
\textbf{Solution détaillée.}
\begin{enumerate}
\item \textbf{Sens commun.} Spontanément, nous condamnons la violence : elle fait souffrir, détruit, humilie. Dire « oui » à la question semble aller de soi.
\item \textbf{Difficulté.} Pourtant, le sens commun lui-même hésite : il excuse celui qui se défend, il admire parfois la révolte contre l'oppresseur, il accepte la force de la police. Condamner \emph{toute} violence reviendrait à mettre sur le même plan l'agresseur et l'agressé qui se défend.
\item \textbf{Tension.} D'un côté, la violence paraît toujours mauvaise parce qu'elle nie autrui ; de l'autre, refuser toute violence, c'est laisser le champ libre aux violents et renoncer à la justice. La violence est-elle un mal absolu, ou un mal qui peut parfois servir le bien ?
\item \textbf{Question directrice.} « Si la violence est la négation de l'autre, faut-il la condamner sans exception ? Mais une condamnation absolue ne risque-t-elle pas de désarmer la victime devant l'agresseur ? Faut-il alors distinguer les violences pour n'en condamner que certaines ? »
\end{enumerate}
\textbf{Conclusion.} Cette problématique annonce naturellement un plan : (I) pourquoi la violence est condamnable ; (II) les cas où elle semble justifiable ; (III) le bilan : distinguer violence illégitime et violence légitime.
\end{exoresolu}

\courssub{Fiche méthode 5 : Mobiliser correctement un auteur}

La problématique posée, le développement exige l'appui des grands auteurs étudiés dans la leçon. Encore faut-il les mobiliser correctement.

\begin{methode}[Citer un auteur dans une copie]
L'appui sur les grands auteurs est exigé au Bac, mais il obéit à des règles :
\begin{enumerate}
\item \textbf{Attribuer exactement} la thèse à l'auteur : ne pas prêter à Rousseau la thèse de Hobbes.
\item \textbf{Citer brièvement} : nom + thèse en une ou deux phrases, pas de longue biographie.
\item \textbf{Expliquer la thèse} avec ses propres mots : que veut dire l'auteur ?
\item \textbf{Articuler avec son propre raisonnement} : l'auteur confirme, précise ou nuance votre argument ; il ne le remplace pas.
\item \textbf{Éventuellement discuter} l'auteur (accord, désaccord argumenté).
\end{enumerate}
\textbf{Réflexes sur cette leçon} : Hobbes = l'homme est naturellement porté à la violence (« l'homme est un loup pour l'homme ») ; Rousseau = l'homme est naturellement bon, c'est la société qui le rend violent ; Weber = l'État détient le monopole de la violence légitime.
\end{methode}

\begin{exoresolu}[Mobilisation des auteurs sur l'origine de la violence]
\textbf{Énoncé.} « La violence vient-elle de la nature humaine ou de la société ? » Répondez en mobilisant Hobbes et Rousseau, puis en justifiant votre propre conclusion.
\tcblower
\textbf{Solution détaillée.}
\begin{enumerate}
\item \textbf{Thèse de Hobbes.} Pour Hobbes, à l'état de nature, l'homme est mus par ses passions (peur, désir de dominer) : « l'homme est un loup pour l'homme ». La violence est donc \cle{naturelle} ; seul l'État, par la crainte de la loi, peut la contenir.
\item \textbf{Explication.} Hobbes veut dire que sans autorité commune, chacun redoute chacun, et cette méfiance rend la guerre de tous contre tous permanente.
\item \textbf{Thèse de Rousseau.} Pour Rousseau, au contraire, l'homme naturel est bon et porté à la pitié ; c'est la \cle{société} — la propriété, l'inégalité, la comparaison aux autres — qui le rend méchant et violent.
\item \textbf{Explication.} Rousseau veut dire que la violence est un produit \cle{culturel} : elle s'apprend dans les rapports sociaux, elle n'est pas inscrite dans notre nature.
\item \textbf{Articulation et jugement personnel.} Les deux thèses s'opposent mais s'éclairent : l'expérience montre qu'il existe une agressivité spontanée (l'enfant qui frappe pour un jouet), ce qui donne raison à Hobbes ; mais l'éducation et le milieu orientent, amplifient ou apaisent cette agressivité (un quartier où règne l'impunité produit plus de violence), ce qui donne raison à Rousseau.
\item \textbf{Conclusion justifiée.} La violence a donc une double origine : une racine naturelle (l'agressivité) que la culture et les institutions peuvent soit déchaîner, soit discipliner. C'est pourquoi l'éducation et la justice sont les vrais remèdes à la violence.
\end{enumerate}
\end{exoresolu}

\courssub{Fiche méthode 6 : Conclure et justifier une conclusion}

Dernière étape du devoir : la conclusion. Trop souvent négligée, elle est pourtant le moment où le candidat tranche la question et justifie sa position finale.

\begin{methode}[Rédiger une conclusion de dissertation ou de réponse]
La conclusion n'est pas un résumé mécanique : elle tranche la question posée.
\begin{enumerate}
\item \textbf{Rappeler la question} de départ en une phrase.
\item \textbf{Bilan des étapes} : rappeler brièvement les deux moments de la réflexion (condamnation / justification possible).
\item \textbf{Trancher clairement} : dire ce que l'on soutient finalement, sans rester dans le « oui et non ».
\item \textbf{Justifier la position finale} par l'argument décisif (ici : la distinction entre violence illégitime et violence légitime).
\item \textbf{Ouvrir éventuellement} sur une question voisine (les moyens non violents de régler les conflits).
\end{enumerate}
\textbf{Erreur à éviter} : introduire une idée nouvelle dans la conclusion.
\end{methode}

\begin{exoresolu}[Rédaction d'une conclusion]
\textbf{Énoncé.} Rédigez la conclusion d'une dissertation sur le sujet : « La violence peut-elle être juste ? »
\tcblower
\textbf{Solution détaillée (conclusion rédigée).}

Nous nous demandions si la violence peut être juste. Notre réflexion a d'abord montré que la violence est par essence condamnable : en usant de la force pour contraindre, elle nie la liberté et la dignité d'autrui, elle détruit la parole et la relation entre les hommes. Pourtant, l'examen de la légitime défense et de la violence de l'État a montré qu'une force réglée par le droit et tournée vers la protection des personnes ne se confond pas avec la violence arbitraire. Nous concluons donc que la violence, strictement définie comme force illégitime et destructrice, ne peut jamais être juste ; ce que l'on appelle parfois « violence juste » est en réalité une \cle{force légitime}, limitée par la loi et par la nécessité. La vraie question n'est donc pas de justifier la violence, mais de savoir comment la justice peut se passer d'elle : c'est tout l'enjeu de l'éducation, du dialogue et de l'État de droit.
\end{exoresolu}

\begin{attention}[Les 5 erreurs qui coûtent des points au Bac]
\begin{enumerate}
\item \textbf{Confondre force et violence.} Écrire « la violence, c'est la force » est une définition fausse : la force devient violence seulement lorsqu'elle contraint autrui contre sa volonté et porte atteinte à son intégrité ou à sa dignité. Toujours distinguer les deux notions dès l'introduction.
\item \textbf{Attribuer les thèses aux mauvais auteurs.} Hobbes = violence naturelle (« l'homme est un loup pour l'homme ») ; Rousseau = violence produite par la société ; Weber = monopole de la violence légitime par l'État. Inverser Hobbes et Rousseau est l'erreur la plus fréquente et la plus sanctionnée.
\item \textbf{Raconter au lieu d'analyser.} Dans l'application à une situation concrète, se contenter de raconter la scène (la bagarre, la confiscation) sans nommer le type de violence (physique, psychologique, économique, légitime) et sans porter de jugement argumenté : cela ne vaut qu'une fraction des points.
\item \textbf{Répondre « oui et non » sans trancher.} Une conclusion qui ne prend pas position (« la violence est parfois bien, parfois mal ») est jugée non justifiée. Il faut trancher en s'appuyant sur une distinction précise (violence illégitime / force légitime) et la justifier.
\item \textbf{Citer un auteur sans l'expliquer ni l'articuler.} Jeter un nom d'auteur dans la copie sans dire sa thèse ni le relier à son propre argument ne rapporte rien : l'examinateur attend une thèse expliquée au service d'un raisonnement personnel.
\end{enumerate}
\end{attention}

\newpage

\coursec{Exercices}

\courssub{Je vérifie mes connaissances}

\begin{exercice}[QCM : les notions de base]
Pour chaque question, une seule réponse est correcte. Recopie la lettre correspondante.
\begin{enumerate}
\item La violence se définit d'abord comme :
\begin{enumerate}
\item[a)] une force naturelle toujours légitime ;
\item[b)] l'usage de la force pour contraindre, blesser ou détruire ;
\item[c)] un simple désaccord verbal ;
\item[d)] une émotion passagère.
\end{enumerate}
\item Selon Hobbes, l'homme à l'état de nature est :
\begin{enumerate}
\item[a)] naturellement bon et pacifique ;
\item[b)] naturellement porté à la violence, « loup pour l'homme » ;
\item[c)] indifférent aux autres ;
\item[d)] incapable de violence.
\end{enumerate}
\item Selon Rousseau, la violence vient :
\begin{enumerate}
\item[a)] de la nature humaine elle-même ;
\item[b)] de la société et de l'inégalité qu'elle engendre ;
\item[c)] des animaux ;
\item[d)] du hasard.
\end{enumerate}
\item Pour Hannah Arendt, pouvoir et violence sont :
\begin{enumerate}
\item[a)] deux mots synonymes ;
\item[b)] deux réalités distinctes : le pouvoir repose sur le consentement, la violence sur les instruments de la force ;
\item[c)] deux formes de langage ;
\item[d)] deux vertus morales.
\end{enumerate}
\end{enumerate}
\end{exercice}

\begin{exercice}[Vrai ou faux : justifier]
Dis si chaque affirmation est vraie ou fausse, puis justifie ta réponse en une ou deux phrases.
\begin{enumerate}
\item La violence ne peut être que physique.
\item La violence verbale (insultes, menaces, rumeurs) n'est pas une véritable violence.
\item La violence dite « légitime défense » est reconnue par le droit.
\item Un État qui use de la force est toujours violent au sens condamnable du terme.
\end{enumerate}
\end{exercice}

\begin{exercice}[Définitions à compléter]
Recopie et complète les définitions suivantes avec les mots de la liste : \emph{force, contraindre, structurelle, psychologique, institutionnelle}.
\begin{enumerate}
\item La violence \cle{physique} est l'usage de la \ldots{} contre le corps d'autrui pour le blesser ou le \ldots{}.
\item La violence \ldots{} est celle qui s'exerce par les structures sociales injustes (misère, exclusion) sans agresseur directement identifiable.
\item La violence \ldots{} atteint la dignité et l'équilibre mental d'une personne (humiliations, harcèlement).
\item La violence est une \ldots{} exercée sur autrui contre sa volonté.
\end{enumerate}
\end{exercice}

\begin{exercice}[Associer auteur et thèse]
Relie chaque auteur à la thèse qui lui correspond, puis rédige une phrase complète pour chaque association.
\begin{center}
\begin{tabularx}{\linewidth}{|l|X|}
\hline
\rowcolor{popblueL} \textbf{Auteur} & \textbf{Thèses proposées} \\
\hline
1. Hobbes & a) L'homme naît bon ; c'est la société qui le corrompt et engendre la violence. \\
\hline
2. Rousseau & b) La violence est l'échec du pouvoir : là où règne la violence, le pouvoir s'affaiblit. \\
\hline
3. Hannah Arendt & c) À l'état de nature, la guerre de tous contre tous rend la violence permanente. \\
\hline
\end{tabularx}
\end{center}
\end{exercice}

\courssub{Je m'entraîne}

\begin{exercice}[Identifier les types de violence]
Pour chacune des situations suivantes, identifie le ou les types de violence en jeu (physique, verbale, psychologique, structurelle, symbolique) et justifie ton choix en deux lignes.
\begin{enumerate}
\item À Douala, un élève est frappé chaque matin à l'entrée du lycée par un groupe de racketteurs.
\item Sur les réseaux sociaux, une élève de Première fait l'objet de moqueries répétées et de rumeurs.
\item Dans un village, des enfants ne peuvent pas aller à l'école parce que l'établissement le plus proche est à \SI{15}{km}.
\item Un chef de quartier humilie publiquement un habitant devant toute la communauté.
\end{enumerate}
\end{exercice}

\begin{exercice}[Hobbes ou Rousseau ?]
Lis le fait divers suivant : « À Yaoundé, une rixe entre bandes rivales de quartiers a fait plusieurs blessés. Les jeunes impliqués affirment se battre pour "défendre leur territoire et leur honneur". »
\begin{enumerate}
\item Explique comment Hobbes interpréterait cette situation (origine naturelle de la violence).
\item Explique comment Rousseau l'interpréterait (rôle de la société, de la pauvreté, de la comparaison sociale).
\item Quelle explication te semble la plus convaincante ? Justifie ta réponse en un paragraphe argumenté.
\end{enumerate}
\end{exercice}

\begin{exercice}[Argumenter : la peine de mort]
La peine de mort a été abolie dans de nombreux pays, mais certains la réclament encore pour les crimes les plus graves.
\begin{enumerate}
\item Rédige un argument en faveur de la peine de mort (répondre à la violence par la violence de l'État).
\item Rédige deux arguments contre, en mobilisant la distinction entre violence et justice.
\item Conclus en défendant ta position personnelle en cinq lignes minimum.
\end{enumerate}
\end{exercice}

\begin{exercice}[Rédiger une problématique]
À partir du sujet suivant : « La violence est-elle toujours condamnable ? »
\begin{enumerate}
\item Dégage le sens du problème posé et formule une problématique claire.
\item Propose un plan en deux ou trois parties, avec pour chaque partie une idée directrice et un auteur mobilisable.
\item Rédige l'accroche et la définition des termes de l'introduction.
\end{enumerate}
\end{exercice}

\courssub{Je me prépare au Bac}

\begin{exercice}[Dissertation : violence et politique]
\textbf{Sujet :} « La violence peut-elle être un moyen d'action politique légitime ? »

On pourra s'appuyer sur les repères suivants : la distinction entre \cle{pouvoir} et \cle{violence} chez Hannah Arendt ; la \cle{violence légitime} de l'État (monopole de la force) ; les luttes de libération et la désobéissance civile ; la non-violence comme force d'action.

\textbf{Travail demandé :}
\begin{enumerate}
\item Rédige une introduction complète (accroche, définition des termes, problématique, annonce du plan).
\item Développe la première partie : la violence semble parfois nécessaire en politique (deux arguments, un exemple tiré de l'histoire ou de l'actualité camerounaise ou africaine).
\item Développe la deuxième partie : la violence détruit le pouvoir et délégitime celui qui l'emploie (deux arguments, un auteur cité).
\item Rédige une conclusion personnelle et nuancée.
\end{enumerate}
\end{exercice}

\begin{exercice}[Explication de texte guidée : Hannah Arendt]
\textbf{Texte :} « Le pouvoir et la violence sont des contraires : là où l'un règne de façon absolue, l'autre est absent. La violence apparaît là où le pouvoir est en danger, mais, livrée à elle-même, elle aboutit à la disparition du pouvoir. [...] La violence peut détruire le pouvoir ; elle est totalement incapable de le créer. » (Hannah Arendt, \emph{Du mensonge à la violence}, adapté.)

\textbf{Questions :}
\begin{enumerate}
\item Dégage la thèse de l'auteur en une phrase.
\item Explique la distinction que fait Arendt entre pouvoir et violence.
\item Pourquoi la violence est-elle, selon le texte, « incapable de créer le pouvoir » ?
\item Illustre cette thèse par un exemple concret (un régime qui ne repose que sur la force).
\item Discussion : un État peut-il se passer totalement de la violence ? Justifie ta réponse.
\end{enumerate}
\end{exercice}

\begin{exercice}[Analyse de document et débat rédigé]
\textbf{Document :} « Violences scolaires à Yaoundé. Plusieurs établissements de la capitale signalent une hausse des bagarres, du racket et du harcèlement entre élèves. Les enseignants dénoncent le manque d'encadrement ; les élèves évoquent la pauvreté, l'influence des réseaux sociaux et les conflits de quartier. » (Article de presse, adapté.)

\textbf{Partie A — Analyse du document}
\begin{enumerate}
\item Identifie dans ce document trois types de violence et justifie chacun.
\item Explique les causes possibles de ces violences en mobilisant Hobbes (nature humaine) puis Rousseau (effets de la société).
\item Propose une solution argumentée, en distinguant ce qui relève de l'école, de la famille et de l'État.
\end{enumerate}

\textbf{Partie B — Débat rédigé : « L'État a-t-il le droit d'user de la violence ? »}
\begin{enumerate}
\item Rédige une introduction qui pose le problème.
\item Formule deux arguments \emph{pour} (ordre public, monopole de la violence légitime).
\item Formule deux arguments \emph{contre} (abus de force, violence d'État).
\item Rédige une conclusion personnelle justifiée.
\end{enumerate}
\end{exercice}



\newpage

\coursec{Corrigés des exercices}

\coursec{Atelier méthodologique : application de la leçon à des situations concrètes}

\begin{corrige}[Exercice 1 — QCM : les notions de base]
\begin{enumerate}
\item \textbf{b)} La violence est l'usage de la force (physique ou autre) pour contraindre autrui, le blesser ou détruire ce qu'il est ou possède. Un désaccord verbal n'est pas en soi une violence, et une émotion n'est pas un acte.
\item \textbf{b)} Dans le \emph{Léviathan}, Hobbes décrit l'état de nature comme une « guerre de tous contre tous » : \emph{l'homme est un loup pour l'homme} (\emph{homo homini lupus}).
\item \textbf{b)} Pour Rousseau, l'homme naît bon ; c'est la vie en société, la propriété et l'inégalité qui le corrompent et engendrent la violence.
\item \textbf{b)} Arendt distingue rigoureusement le \cle{pouvoir}, qui naît de l'action commune et du consentement, et la \cle{violence}, qui repose sur des instruments de force. Pour elle, la violence est l'opposé du pouvoir : « là où la violence règne absolument, le pouvoir est absent ».
\end{enumerate}
\textbf{Points de vigilance :} ne pas confondre violence et simple conflit ; bien retenir que pour Arendt, pouvoir et violence s'\emph{opposent} radicalement (la violence détruit le pouvoir) et ne sont pas des synonymes.
\end{corrige}

\begin{corrige}[Exercice 2 — Vrai ou faux : justifier]
\begin{enumerate}
\item \textbf{Faux.} La violence peut aussi être verbale, psychologique, symbolique ou structurelle : elle atteint autrui dans sa dignité ou ses conditions de vie, pas seulement dans son corps.
\item \textbf{Faux.} Les insultes, menaces et rumeurs blessent moralement, humilient et peuvent détruire l'équilibre psychique d'une personne : ce sont de véritables violences, parfois punies par la loi (harcèlement, menaces de mort).
\item \textbf{Vrai.} Le droit reconnaît la \cle{légitime défense} : répondre par une force proportionnée à une agression actuelle et injuste n'est pas condamnable.
\item \textbf{Faux.} L'État détient le \cle{monopole de la violence légitime} (police, justice) : l'usage encadré et légal de la force pour maintenir l'ordre se distingue de la violence arbitraire et abusive.
\end{enumerate}
\textbf{Points de vigilance :} la justification est exigée pour chaque item ; citer la condition de \emph{proportionnalité} pour la légitime défense et la notion de \emph{monopole de la violence légitime} (Weber) pour la question 4.
\end{corrige}

\begin{corrige}[Exercice 3 — Définitions à compléter]
\begin{enumerate}
\item La \cle{violence physique} est l'usage de la \textbf{force} contre le corps d'autrui pour le blesser ou le \textbf{contraindre}.
\item La violence \textbf{structurelle} est celle qui s'exerce par les structures sociales injustes (misère, exclusion) sans agresseur directement identifiable.
\item La violence \textbf{psychologique} atteint la dignité et l'équilibre mental d'une personne (humiliations, harcèlement).
\item La violence est une \textbf{contrainte} exercée sur autrui contre sa volonté. (Le mot « \emph{institutionnel} » était un intrus : la violence institutionnelle est commise par une institution identifiable, ce qui diffère de la violence structurelle.)
\end{enumerate}
\textbf{Points de vigilance :} bien distinguer violence \emph{structurelle} (injustice des structures sociales, sans agresseur identifiable) et violence \emph{institutionnelle} (commise par une institution identifiable) ; ici le mot attendu était « contrainte », forme nominale du verbe « contraindre ».
\end{corrige}

\begin{corrige}[Exercice 4 — Associer auteur et thèse]
\begin{itemize}
\item \textbf{1 -- c)} Hobbes soutient qu'à l'état de nature, l'égalité des forces et la défiance mutuelle plongent les hommes dans une guerre permanente de tous contre tous.
\item \textbf{2 -- a)} Rousseau affirme que l'homme naît bon et que c'est la société, avec la propriété et l'inégalité, qui le corrompt et engendre la violence.
\item \textbf{3 -- b)} Hannah Arendt montre que la violence est l'échec du pouvoir : un régime qui ne repose que sur la force voit son pouvoir s'affaiblir, car « la violence peut détruire le pouvoir ; elle est totalement incapable de le créer ».
\end{itemize}
\textbf{Points de vigilance :} rédiger une phrase complète pour chaque association (et non une simple lettre) ; ne pas inverser Hobbes (violence naturelle, remède : État fort) et Rousseau (violence culturelle, produit de la société inégalitaire).
\end{corrige}

\begin{corrige}[Exercice 5 — Identifier les types de violence]
\begin{enumerate}
\item \textbf{Violence physique} (coups) et \textbf{extorsion} (racket : violence pour obtenir des biens par la menace). Le corps de l'élève est directement atteint et ses biens lui sont pris contre sa volonté.
\item \textbf{Violence verbale et psychologique} (cyberharcèlement) : les moqueries et rumeurs répétées atteignent la dignité et l'équilibre mental de l'élève, sans contact physique.
\item \textbf{Violence structurelle} : l'éloignement de l'école résulte d'une organisation sociale injuste (manque d'infrastructures) qui prive des enfants de leur droit à l'éducation, sans agresseur identifiable.
\item \textbf{Violence symbolique et psychologique} : l'humiliation publique détruit la dignité et le statut social de la personne aux yeux de la communauté.
\end{enumerate}
\textbf{Points de vigilance :} une même situation peut combiner plusieurs types de violence ; le critère de la violence structurelle est l'\emph{absence d'agresseur directement identifiable} ; la justification en deux lignes est obligatoire.
\end{corrige}

\begin{corrige}[Exercice 6 — Hobbes ou Rousseau ?]
\begin{enumerate}
\item Pour Hobbes, ces jeunes illustrent la tendance naturelle de l'homme à la violence : en l'absence d'une autorité capable de les contenir, la rivalité, la peur et le désir de domination dégénèrent en guerre de tous contre tous. Seul un pouvoir fort (l'État, la police) peut rétablir la paix.
\item Pour Rousseau, ces jeunes ne sont pas mauvais par nature : c'est la société qui les a corrompus. La pauvreté, l'exclusion, la comparaison sociale et le besoin de reconnaissance (« l'honneur ») sont des produits de la vie en société inégalitaire.
\item \textbf{Réponse personnelle attendue (exemple) :} L'explication de Rousseau me semble plus convaincante car les violences de quartier touchent surtout les milieux défavorisés, ce qui montre le rôle déterminant des conditions sociales (pauvreté, exclusion, absence de perspectives). Toutefois, on ne peut nier avec Hobbes qu'une autorité légitime est nécessaire pour contenir les conflits immédiats et protéger les victimes.
\end{enumerate}
\textbf{Points de vigilance :} la question 3 exige un \emph{paragraphe argumenté} (thèse, argument, exemple), pas une simple opinion ; toute position est acceptée dès lors qu'elle est justifiée philosophiquement.
\end{corrige}

\begin{corrige}[Exercice 7 — Argumenter : la peine de mort]
\begin{enumerate}
\item \textbf{Pour (thèse de la rétribution et de la dissuasion) :} face à un criminel qui a détruit des vies, seule une peine équivalente semble réparer le mal (loi du talion) et dissuader les autres criminels ; supprimer la peine de mort serait faire preuve de faiblesse face à la violence extrême.
\item \textbf{Contre (thèse de la justice comme non-violence) :} (a) la \cle{justice} n'est pas la \cle{vengeance} : tuer au nom de la loi reproduit la violence qu'on prétend punir, alors que la justice vise à réparer et à protéger, non à détruire ; (b) l'\cle{erreur judiciaire} est irréversible : exécuter un innocent est une violence d'État irréparable, et aucune société n'est à l'abri d'un procès injuste.
\item \textbf{Conclusion personnelle (exemple) :} La peine de mort est une violence légalisée qui ne répare rien et ne rend pas la vie aux victimes. Elle place l'État sur le même plan que le criminel en usant du droit de vie et de mort. La justice doit se distinguer de la violence par la mesure et le respect de la vie humaine : la prison à perpétuité (avec possibilité de révision) protège la société sans tuer, tout en laissant une porte ouverte à la réparation morale.
\end{enumerate}
\textbf{Points de vigilance :} savoir défendre d'abord la thèse adverse avec loyauté ; mobiliser la distinction \cle{justice}/\cle{vengeance} ; respecter la longueur minimale de cinq lignes pour la conclusion ; ne pas introduire de nouvel argument en conclusion.
\end{corrige}

\begin{corrige}[Exercice 8 — Rédiger une problématique]
\begin{enumerate}
\item \textbf{Problématique :} Si la violence paraît toujours mauvaise parce qu'elle détruit autrui, certaines situations (légitime défense, résistance à l'oppression, force de l'État) semblent la rendre nécessaire, voire juste. La violence est-elle donc condamnable par essence, ou sa valeur morale dépend-elle de ses fins et de ses conditions d'exercice ?
\item \textbf{Plan possible :}
\begin{itemize}
\item \emph{Partie 1 :} La violence est condamnable parce qu'elle détruit la dignité humaine et le pouvoir (Arendt) ; elle engendre la violence.
\item \emph{Partie 2 :} Pourtant, certaines violences semblent légitimes : la force de l'État pour garantir l'ordre (Hobbes, monopole de la violence légitime), la légitime défense, la résistance à l'oppression.
\item \emph{Partie 3 :} Il faut distinguer violence destructrice et force légitime, encadrée par le droit, proportionnée à sa fin et soumise au contrôle démocratique.
\end{itemize}
\item \textbf{Accroche et définitions :} « Chaque jour, l'actualité nous parle de violence : bagarres, guerres, répressions. Mais la police qui arrête un malfaiteur use aussi de la force. La violence est l'usage de la force pour contraindre, blesser ou détruire ; condamnable signifie qui mérite d'être blâmé et puni. Dès lors, toute violence mérite-t-elle d'être condamnée ? »
\end{enumerate}
\textbf{Points de vigilance :} définir \emph{tous} les termes du sujet (« violence » et « condamnable ») ; la problématique doit être une vraie question, pas une affirmation ; chaque partie du plan doit annoncer une idée directrice et un auteur ou une référence théorique.
\end{corrige}

\begin{corrige}[Exercice 9 — Dissertation : violence et politique (éléments de correction)]
\begin{enumerate}
\item \textbf{Introduction :} Accroche sur les luttes politiques de l'histoire (indépendances africaines, lutte contre l'apartheid) ; Définitions : la \cle{violence} est l'usage de la force pour contraindre ou détruire ; \cle{légitime} signifie conforme au droit ou à la justice ; l'\cle{action politique} vise à gouverner la cité. Problématique : La violence, qui détruit, peut-elle servir la politique, qui vise à construire un ordre commun ? Annonce du plan en deux parties.
\item \textbf{Partie 1 : La violence comme moyen politique (libération, fondation de l'ordre).}
\begin{itemize}
\item (a) Argument : Face à l'oppression et à l'injustice, la violence peut être le seul moyen de se faire entendre et de se libérer (luttes anticoloniales, résistance à la tyrannie). Exemple : les luttes pour les indépendances en Afrique (Cameroun, Algérie) ou la résistance face au nazisme.
\item (b) Argument : L'État lui-même fonde l'ordre sur le monopole de la violence légitime (Hobbes : sans force, les lois restent lettre morte ; Weber : l'État est le détenteur du monopole de la violence physique légitime). Exemple : la force de police qui protège les citoyens contre les criminels.
\end{itemize}
\item \textbf{Partie 2 : Les limites et les dangers de la violence politique (Arendt, non-violence).}
\begin{itemize}
\item (a) Arendt : « La violence peut détruire le pouvoir ; elle est totalement incapable de le créer » : un régime qui ne repose que sur la force perd le consentement et s'affaiblit ; la violence détruit le dialogue politique.
\item (b) La violence engendre la violence : elle corrompt la cause qu'elle prétend servir et délégitime celui qui l'emploie. La \cle{non-violence} (désobéissance civile, Gandhi, Martin Luther King) peut être une force d'action plus efficace et plus durable car elle vise à convaincre, non à contraindre.
\end{itemize}
\item \textbf{Conclusion :} La violence n'est jamais un bien en soi ; seule une force strictement encadrée par le droit, proportionnée et au service de la justice peut être admise comme « force légitime ». La véritable action politique repose sur le \cle{pouvoir}, c'est-à-dire le consentement et la parole, non sur la violence.
\end{enumerate}
\textbf{Barème indicatif (sur 20) :} Introduction (accroche, définitions, problématique, plan) : 4 pts ; Partie 1 (deux arguments + exemple précis) : 5 pts ; Partie 2 (deux arguments + auteur cité nommément) : 5 pts ; Conclusion personnelle et nuancée : 3 pts ; Clarté, orthographe et organisation : 3 pts.

\textbf{Points de vigilance :} Citer au moins un auteur \emph{nommément} (Hobbes, Rousseau, Arendt, Weber, Gandhi) ; l'exemple doit être précis (histoire africaine ou camerounaise de préférence) ; éviter la paraphrase du sujet ; la conclusion doit trancher avec nuance, sans introduire de notion nouvelle.
\end{corrige}

\begin{corrige}[Exercice 10 — Explication de texte guidée : Hannah Arendt]
\begin{enumerate}
\item \textbf{Thèse :} Pouvoir et violence s'excluent : la violence détruit le pouvoir sans jamais pouvoir le fonder.
\item Le \cle{pouvoir} naît du consentement et de l'action commune des hommes (obéir parce qu'on reconnaît une autorité légitime) ; la \cle{violence} repose sur des instruments de force (armes, terreur) et contraint sans convaincre.
\item Parce que le pouvoir suppose l'adhésion libre des citoyens : la force peut obtenir une obéissance par la peur, mais jamais une reconnaissance durable ; dès que la terreur faiblit, l'ordre s'effondre (exemple de la chute des dictatures).
\item \textbf{Exemple concret :} Une dictature militaire qui ne tient que par la répression s'effondre dès que l'armée ou le peuple cesse d'obéir par peur ; à l'inverse, une autorité reconnue (chef traditionnel respecté, gouvernement élu) gouverne avec peu de force car elle repose sur le pouvoir (consentement).
\item \textbf{Discussion :} Non, un État ne peut se passer totalement de la force (police contre les criminels, armée pour la défense), mais il doit en limiter l'usage au strict nécessaire, sous le contrôle du droit et de la représentation nationale : c'est le \cle{monopole de la violence légitime}, distinct de la violence arbitraire. La force légitime protège le pouvoir, elle ne le remplace pas.
\end{enumerate}
\textbf{Barème indicatif (sur 20) :} Thèse dégagée : 3 pts ; Distinction pouvoir/violence : 4 pts ; Explication de l'incapacité de la violence à créer le pouvoir : 4 pts ; Exemple concret pertinent : 3 pts ; Discussion argumentée : 4 pts ; Expression : 2 pts.

\textbf{Points de vigilance :} S'appuyer sur le texte (citer de courts extraits entre guillemets) ; ne pas attribuer à Arendt l'idée que toute force est condamnable : elle vise la \emph{violence} (instrumentale, arbitraire), non la force légitime encadrée par le droit ; la question 5 exige une position justifiée, pas un simple oui ou non.
\end{corrige}

\begin{corrige}[Exercice 11 — Analyse de document et débat rédigé (éléments de correction)]
\textbf{Partie A : Analyse de la situation (violences scolaires)}
\begin{enumerate}
\item (a) \textbf{Violence physique} : les bagarres blessent les corps ; (b) \textbf{Extorsion / Violence avec contrainte} : le racket dépouille les élèves par la menace ; (c) \textbf{Violence psychologique} : le harcèlement (dont cyberharcèlement) détruit la dignité et l'équilibre mental des victimes.
\item Selon \textbf{Hobbes}, le manque d'encadrement laisse libre cours à la tendance naturelle des hommes à la rivalité et à la domination : sans autorité, la violence éclate (état de nature). Selon \textbf{Rousseau}, la pauvreté, les conflits de quartier, les modèles violents des réseaux sociaux et l'absence de repères éducatifs sont des produits de la société qui corrompent des élèves naturellement bons.
\item \textbf{École :} Renforcer la surveillance, le dialogue, la médiation par les pairs et l'éducation à la non-violence (citoyenneté). \textbf{Famille :} Assurer l'encadrement, l'éducation morale, la communication et le suivi de la scolarité. \textbf{État :} Lutter contre la pauvreté et l'exclusion (causes structurelles), garantir la sécurité aux abords des écoles, faire appliquer la loi avec mesure (protection des mineurs, justice restaurative).
\end{enumerate}
\textbf{Partie B : Débat — L'État a-t-il le droit d'user de la violence ?}
\begin{enumerate}
\item \textbf{Introduction :} L'État dispose de la police, de la gendarmerie, de l'armée ; il peut user de la force. Mais la force peut-elle être un droit ? Problématique : L'usage de la violence par l'État est-il légitime, et à quelles conditions strictes ?
\item \textbf{Pour (la légitimité de la force étatique) :} (a) Sans force publique, l'ordre s'effondre et les plus forts oppriment les plus faibles (Hobbes : la sécurité est la fin de l'État) ; (b) L'État détient le \cle{monopole de la violence légitime} (Weber) : encadrée par la loi, sa force protège les citoyens, fait respecter le droit et les décisions de justice.
\item \textbf{Contre (les risques de dérive) :} (a) L'État peut abuser de la force (répressions arbitraires, brutalités policières, torture) : la violence d'État devient alors la pire des violences car elle est organisée, systématique et impunie ; (b) Selon Arendt, un pouvoir qui ne repose que sur la violence se détruit lui-même en perdant le consentement des citoyens (crise de légitimité, révoltes).
\item \textbf{Conclusion :} L'État a le droit — et le devoir — d'user de la force, mais seulement de manière \emph{légale} (prévue par la loi), \emph{proportionnée} (adaptée à la menace), \emph{nécessaire} (dernier recours) et \emph{contrôlée} (par la justice, le parlement, l'opinion). La \cle{force légitime} se distingue de la \cle{violence arbitraire} par sa finalité (protéger, non opprimer) et par ses limites (l'État de droit).
\end{enumerate}
\textbf{Barème indicatif (sur 20) :} Partie A (Identification des violences justifiées : 3 pts ; Explication des causes Hobbes/Rousseau : 3 pts ; Solutions argumentées par acteur : 2 pts) : 8 pts ; Partie B (Introduction avec problématique : 2 pts ; Arguments Pour : 3 pts ; Arguments Contre : 3 pts ; Conclusion nuancée avec conditions : 2 pts) : 10 pts ; Expression et organisation : 2 pts.

\textbf{Points de vigilance :} En Partie A, chaque type de violence doit être \emph{justifié} à partir du document (citer les indices textuels) ; En Partie B, distinguer clairement \emph{force légitime} (légale, proportionnée, contrôlée, finalité protectrice) et \emph{violence arbitraire} (illégale, disproportionnée, finalité répressive) ; Citer explicitement les conditions d'application du monopole de la violence légitime (légalité, proportionnalité, nécessité, finalité de protection).
\end{corrige}

\newpage

\coursec{Fiche bilan}

\begin{popfigure}
\begin{center}\tcbox[colback=popdark,colframe=popdark,coltext=white,arc=9pt,boxsep=4pt,left=6pt,right=6pt]{\bfseries\small La violence}\end{center}
\vspace{-2pt}
\begin{tcbraster}[raster columns=3,raster equal height=rows,raster column skip=6pt,raster row skip=6pt,enhanced,blankest]
\begin{tcolorbox}[enhanced,colback=white,colframe=popblue,boxrule=0.9pt,arc=3pt,title={Définition},fonttitle=\bfseries\scriptsize,coltitle=white,colbacktitle=popblue,left=4pt,right=4pt,top=2pt,bottom=2pt,boxsep=1pt,toptitle=1.5pt,bottomtitle=1.5pt,halign=left]
\begin{itemize}[nosep,leftmargin=*,itemsep=1pt,font=\scriptsize]
\item Force illégitime
\item Atteinte à l'autre
\end{itemize}
\end{tcolorbox}
\begin{tcolorbox}[enhanced,colback=white,colframe=poporange,boxrule=0.9pt,arc=3pt,title={Origine},fonttitle=\bfseries\scriptsize,coltitle=white,colbacktitle=poporange,left=4pt,right=4pt,top=2pt,bottom=2pt,boxsep=1pt,toptitle=1.5pt,bottomtitle=1.5pt,halign=left]
\begin{itemize}[nosep,leftmargin=*,itemsep=1pt,font=\scriptsize]
\item Naturelle : Hobbes
\item Culturelle : Rousseau
\end{itemize}
\end{tcolorbox}
\begin{tcolorbox}[enhanced,colback=white,colframe=popgreen,boxrule=0.9pt,arc=3pt,title={Types},fonttitle=\bfseries\scriptsize,coltitle=white,colbacktitle=popgreen,left=4pt,right=4pt,top=2pt,bottom=2pt,boxsep=1pt,toptitle=1.5pt,bottomtitle=1.5pt,halign=left]
\begin{itemize}[nosep,leftmargin=*,itemsep=1pt,font=\scriptsize]
\item Physique
\item Psychologique
\item Symbolique
\end{itemize}
\end{tcolorbox}
\begin{tcolorbox}[enhanced,colback=white,colframe=poppurple,boxrule=0.9pt,arc=3pt,title={Condamnation},fonttitle=\bfseries\scriptsize,coltitle=white,colbacktitle=poppurple,left=4pt,right=4pt,top=2pt,bottom=2pt,boxsep=1pt,toptitle=1.5pt,bottomtitle=1.5pt,halign=left]
\begin{itemize}[nosep,leftmargin=*,itemsep=1pt,font=\scriptsize]
\item Kant : dignité
\item Droit et État
\end{itemize}
\end{tcolorbox}
\begin{tcolorbox}[enhanced,colback=white,colframe=poppink,boxrule=0.9pt,arc=3pt,title={Justification ?},fonttitle=\bfseries\scriptsize,coltitle=white,colbacktitle=poppink,left=4pt,right=4pt,top=2pt,bottom=2pt,boxsep=1pt,toptitle=1.5pt,bottomtitle=1.5pt,halign=left]
\begin{itemize}[nosep,leftmargin=*,itemsep=1pt,font=\scriptsize]
\item Légitime défense
\item Révolution
\end{itemize}
\end{tcolorbox}
\begin{tcolorbox}[enhanced,colback=white,colframe=popgold,boxrule=0.9pt,arc=3pt,title={Bilan},fonttitle=\bfseries\scriptsize,coltitle=white,colbacktitle=popgold,left=4pt,right=4pt,top=2pt,bottom=2pt,boxsep=1pt,toptitle=1.5pt,bottomtitle=1.5pt,halign=left]
\begin{itemize}[nosep,leftmargin=*,itemsep=1pt,font=\scriptsize]
\item Violence $\neq$ force
\item Primat du droit
\end{itemize}
\end{tcolorbox}
\begin{tcolorbox}[enhanced,colback=white,colframe=popmuted,boxrule=0.9pt,arc=3pt,title={Méthode},fonttitle=\bfseries\scriptsize,coltitle=white,colbacktitle=popmuted,left=4pt,right=4pt,top=2pt,bottom=2pt,boxsep=1pt,toptitle=1.5pt,bottomtitle=1.5pt,halign=left]
\begin{itemize}[nosep,leftmargin=*,itemsep=1pt,font=\scriptsize]
\item Distinguer / nuancer
\end{itemize}
\end{tcolorbox}
\end{tcbraster}

\legende{Carte mentale de la leçon : la violence (définition, origine, types, procès, méthode).}
\end{popfigure}

\begin{bilan}
\textbf{Définitions clés}
\begin{itemize}
\item \cle{Violence} : usage de la force qui porte atteinte à l'intégrité physique ou morale d'une personne, en violation du droit ou de la dignité.
\item \cle{Force} : capacité d'agir ; elle n'est violence que lorsqu'elle devient illégitime ou excessive.
\item \cle{Violence symbolique} (Bourdieu) : domination douce, imposée par le langage, l'éducation ou les normes sociales.
\item \cle{Légitime défense} : riposte proportionnée à une agression en cours, reconnue par le droit.
\end{itemize}

\textbf{Thèses des auteurs à connaître par cœur}
\begin{center}
\begin{tabularx}{\linewidth}{|l|X|}\hline
\rowcolor{popblueL} \textbf{Auteur} & \textbf{Thèse} \\ \hline
Hobbes & L'homme est naturellement violent : « l'homme est un loup pour l'homme » ; seul l'État (Léviathan) instaure la paix. \\ \hline
Rousseau & L'homme est naturellement bon ; c'est la société (propriété, inégalités) qui le rend violent. \\ \hline
Kant & L'homme doit être traité comme une fin, jamais comme un moyen : toute violence est contraire à la dignité. \\ \hline
Bourdieu & La violence la plus efficace est symbolique : elle s'exerce avec le consentement des dominés. \\ \hline
\end{tabularx}
\end{center}

\textbf{Types de violence}
\begin{itemize}
\item Physique (coups, guerre), psychologique (humiliation, harcèlement), verbale (insultes), économique et sociale (misère, exclusion), symbolique (domination culturelle).
\end{itemize}

\textbf{Méthodes express}
\begin{itemize}
\item Toujours \textbf{distinguer} violence et force, violence et conflit, violence subie et violence exercée.
\item Pour le procès : présenter d'abord la \textbf{condamnation} (morale et juridique), puis examiner les \textbf{cas limites} (légitime défense, résistance à l'oppression) avant de conclure avec nuance.
\item Illustrer par des exemples concrets camerounais : rixes scolaires, harcèlement, rôle de la police et des tribunaux, dialogue comme alternative.
\end{itemize}
\end{bilan}

\begin{aretenir}[Checklist avant le Bac]
\begin{itemize}
\item[$\square$] Je sais définir la violence et la distinguer de la simple force.
\item[$\square$] Je connais la thèse de Hobbes (violence naturelle) et celle de Rousseau (violence culturelle).
\item[$\square$] Je sais expliquer le rôle de l'État selon Hobbes.
\item[$\square$] Je peux citer au moins quatre types de violence avec un exemple chacun.
\item[$\square$] Je comprends la notion de violence symbolique (Bourdieu).
\item[$\square$] Je sais formuler la condamnation kantienne de la violence (la dignité).
\item[$\square$] Je sais ce qu'est la légitime défense et ses conditions.
\item[$\square$] Je sais discuter le cas de la révolte contre l'oppression sans tomber dans la justification de tout.
\item[$\square$] Je sais construire une problématique : « Toute violence est-elle condamnable ? »
\item[$\square$] Je sais rédiger une conclusion nuancée : primat du droit et du dialogue, sans nier les situations limites.
\item[$\square$] Je dispose d'exemples camerounais concrets pour illustrer chaque partie.
\item[$\square$] Je sais structurer ma dissertation : introduction problématisée, deux ou trois parties argumentées, conclusion.
\end{itemize}
\end{aretenir}

\findecours

\end{document}
